% 来自2026-10-21统一复核稿；题号、图号和解答保留。
\clearpage
\begingroup
\fancyhead[R]{\small 模块一：数学工具与场论语言}
\chapter{模块一：数学工具与场论语言}\label{m1:ux6a21ux5757ux4e00ux6570ux5b66ux5de5ux5177ux4e0eux573aux8bbaux8bedux8a00}

模块一介绍后文反复使用的矢量分析和场论工具。正文先给出定义、计算公式和积分定理，附录再讨论高阶恒等式、分量证明与分布理论。

{\def\LTcaptype{none} % do not increment counter
\begin{longtable}[]{@{}
  >{\raggedright\arraybackslash}p{(\linewidth - 4\tabcolsep) * \real{0.2500}}
  >{\raggedright\arraybackslash}p{(\linewidth - 4\tabcolsep) * \real{0.2500}}
  >{\raggedright\arraybackslash}p{(\linewidth - 4\tabcolsep) * \real{0.5000}}@{}}
\toprule\noalign{}
\begin{minipage}[b]{\linewidth}\raggedright
层级
\end{minipage} & \begin{minipage}[b]{\linewidth}\raggedright
内容
\end{minipage} & \begin{minipage}[b]{\linewidth}\raggedright
内容定位
\end{minipage} \\
\midrule\noalign{}
\endhead
\bottomrule\noalign{}
\endlastfoot
基础必讲 &
三重积、纳布拉算符的基本意义、梯度、散度、旋度、线积分、通量、高斯定理、斯托克斯定理、常用坐标微元
& 构成全书反复调用的基本语言与计算工具 \\
拔高选讲 &
纳布拉高阶恒等式、恒等式的分量证明、三重积与纳布拉算符的结构关联、\(\delta\)
函数的分布理论 & 解释场论公式的严格结构，并服务于综合问题 \\
\end{longtable}
}

基础部分侧重计算和使用条件；选讲部分补充证明，并处理奇异源等需要分布理论的问题。

\begin{center}
\includegraphics[width=\linewidth,height=.40\textheight,keepaspectratio]{\detokenize{assets/module-01/tikz-v11-module1-method-map.pdf}}
\par\smallskip{\small 图示：模块一从有向几何、微分运算和积分定理走向奇异源}
\end{center}

\section{专题一：矢量分析中的三重积（基础必讲）}\label{m1:ux4e13ux9898ux4e00ux77e2ux91cfux5206ux6790ux4e2dux7684ux4e09ux91cdux79efux57faux7840ux5fc5ux8bb2}

在经典力学与电动力学中，空间矢量的复杂交互常涉及三个矢量的乘积运算。三重积是处理角动量、力矩、矢量恒等式和场论推导的基本代数工具。本专题中的叉乘均指取定右手坐标系中的三维普通矢量叉乘；含有微分算子
\(\nabla\) 的恒等式在后续专题另行讨论。

\subsection{一、对象、定义与适用范围}\label{m1:ux4e00ux5bf9ux8c61ux5b9aux4e49ux4e0eux9002ux7528ux8303ux56f4}

对于三维欧几里得空间中的任意三个矢量 \(\mathbf{A}\)、\(\mathbf{B}\) 和
\(\mathbf{C}\)，可以构造两种具有明确几何与物理意义的三重积：

本节默认使用取定方向的右手正交基
\(\{\mathbf i,\mathbf j,\mathbf k\}\)。以下恒等式首先是普通自由矢量的代数关系；若
\(\mathbf A,\mathbf B,\mathbf C\)
是矢量场，则恒等式在同一空间点逐点成立。微分算子 \(\nabla\)
具有作用对象，不能当作普通矢量直接代入。

\subsubsection{1. 标量三重积}\label{m1:ux6807ux91cfux4e09ux91cdux79ef}

定义为一个点乘与一个叉乘的组合，结果为标量：

\[
\mathbf{A} \cdot (\mathbf{B} \times \mathbf{C})
\]

\textbf{轮换对称性定理}：标量三重积在矢量的循环轮换下保持不变，若交换任意两个矢量则变号：

\[
\mathbf{A} \cdot (\mathbf{B} \times \mathbf{C}) = \mathbf{B} \cdot (\mathbf{C} \times \mathbf{A}) = \mathbf{C} \cdot (\mathbf{A} \times \mathbf{B})
\]

\par\addvspace{.4\baselineskip}
\noindent\begin{minipage}{\linewidth}

\textbf{例 M1-01：右手基矢三重积的符号}

在三维欧氏空间中选定右手单位正交基
\(\mathbf i,\mathbf j,\mathbf k\)，并将三个基矢的起点移至同一原点。分别计算
\(\mathbf i\cdot(\mathbf j\times\mathbf k)\) 和
\(\mathbf i\cdot(\mathbf k\times\mathbf j)\)；用由基矢张成的平行六面体说明，交换后改变的是普通体积还是有向体积的符号。


\end{minipage}\par

\par\noindent\begin{minipage}{\linewidth}
\textbf{解答}

\begin{center}
\includegraphics[width=\linewidth,height=.40\textheight,keepaspectratio]{\detokenize{assets/module-01/examples/M1-01.pdf}}
\par\smallskip{\small 图 M1-01：右手基矢三重积的符号}
\end{center}
\end{minipage}\par

图
M1-01中的底面由两个基矢定向，另一个基矢读取法向投影。\(\mathbf j\times\mathbf k\)
沿 \(+\mathbf i\)，所以点乘为 \(+1\)；交换两个底面方向后，法向变成
\(-\mathbf i\)，点乘变为
\(-1\)。右侧两种取向的无向体积相同，变的是有向体积的符号。

\[
\mathbf i\cdot(\mathbf j\times\mathbf k)=1.
\]

若交换其中任意两个矢量，例如交换 \(\mathbf j,\mathbf k\)，则

\[
\mathbf i\cdot(\mathbf k\times\mathbf j)=-1.
\]

这正是有向体积随取向反转而变号的最小检验。

\par\addvspace{.4\baselineskip}
\noindent\begin{minipage}{\linewidth}

\textbf{例 M1-02：磁场力的瞬时功率}

质量为 \(m>0\)、电荷量为 \(q\) 的非相对论粒子，在某时刻的速度为
\(\mathbf v\)，所在处磁感应强度为 \(\mathbf B\)。本题只讨论磁场力
\(\mathbf F_B=q\mathbf v\times\mathbf B\) 的作用。求磁场力的瞬时功率
\(P_B=\mathbf F_B\cdot\mathbf v\)，用标量三重积说明它能否改变粒子的动能与速率，并讨论
\(q\) 变号是否改变该结论。


\end{minipage}\par

\par\noindent\begin{minipage}{\linewidth}
\textbf{解答}

\begin{center}
\includegraphics[width=\linewidth,height=.40\textheight,keepaspectratio]{\detokenize{assets/module-01/examples/M1-02.pdf}}
\par\smallskip{\small 图 M1-02：磁场力的瞬时功率}
\end{center}
\end{minipage}\par

图 M1-02中速度向右、磁场取纸外、正电荷受力向下，直角标记就是
\(\mathbf F_B\perp\mathbf v\)。求功率时读取沿速度的力分量，因此
\(\mathbf F_B\cdot\mathbf v=0\)。负电荷会反转力箭头，却仍不产生沿速度的分量；以下三重积计算把这个几何关系写成代数式。

带电粒子的磁场力为
\(\mathbf F_B=q\mathbf v\times\mathbf B\)。磁场力的瞬时功率为

\[
P_B=\mathbf F_B\cdot\mathbf v
=q(\mathbf v\times\mathbf B)\cdot\mathbf v
=q\mathbf v\cdot(\mathbf v\times\mathbf B)=0.
\]

最后一步是标量三重积的循环轮换；其中出现两个相同的矢量
\(\mathbf v\)，对应的行列式有两行相同，故结果为零。因此磁场力可以改变速度方向，却不能直接改变粒子的动能和速率。这里的结论只针对磁场力本身；若同时存在电场，电场力仍可做功。

\subsubsection{2. 矢量三重积}\label{m1:ux77e2ux91cfux4e09ux91cdux79ef}

定义为两个连续的叉乘运算，结果为矢量：

\[
\mathbf{A} \times (\mathbf{B} \times \mathbf{C})
\]

\textbf{矢量三重积恒等式}：矢量三重积可转化为原矢量的线性组合：

\[
\mathbf{A} \times (\mathbf{B} \times \mathbf{C}) = \mathbf{B}(\mathbf{A} \cdot \mathbf{C}) - \mathbf{C}(\mathbf{A} \cdot \mathbf{B})
\]

\par\addvspace{.4\baselineskip}
\noindent\begin{minipage}{\linewidth}

\textbf{例 M1-03：矢量三重积的基矢检验}

在右手单位正交基中取
\(\mathbf A=\mathbf i\)、\(\mathbf B=\mathbf j\)、\(\mathbf C=\mathbf i\)。先逐层计算
\(\mathbf A\times(\mathbf B\times\mathbf C)\)，再代入矢量三重积展开式核对。另取
\(\mathbf A=\mathbf i,\mathbf B=\mathbf i,\mathbf C=\mathbf j\)，比较两种括号次序的结果，判断叉乘能否省略括号，并说明第一组计算结果与
\(\mathbf B,\mathbf C\) 张成平面的关系。


\end{minipage}\par

\par\noindent\begin{minipage}{\linewidth}
\textbf{解答}

\begin{center}
\includegraphics[width=\linewidth,height=.40\textheight,keepaspectratio]{\detokenize{assets/module-01/examples/M1-03.pdf}}
\par\smallskip{\small 图 M1-03：矢量三重积的基矢检验}
\end{center}
\end{minipage}\par

沿图
M1-03中的基矢先读内层：\(\mathbf j\times\mathbf i=-\mathbf k\)。再用右手定则算外层
\(\mathbf i\times(-\mathbf k)=\mathbf j\)，结果回到
\(\mathbf i,\mathbf j\) 张成的平面。展开式中两个投影系数分别为 \(1\) 和
\(0\)，所以也只剩图中的 \(\mathbf j\) 方向。

取
\(\mathbf A=\mathbf i\)、\(\mathbf B=\mathbf j\)、\(\mathbf C=\mathbf i\)。左边为

\[
\mathbf i\times(\mathbf j\times\mathbf i)
=\mathbf i\times(-\mathbf k)=\mathbf j.
\]

右边为

\[
\mathbf j(\mathbf i\cdot\mathbf i)-\mathbf i(\mathbf i\cdot\mathbf j)
=\mathbf j.
\]

两边一致；同时结果 \(\mathbf j\) 落在 \(\mathbf B,\mathbf C\) 张成的
\(xy\) 平面内，符合后文的几何结论。

\textbf{顺序不能交换：} 括号表示先算内层叉乘，因而一般有

\[
\mathbf A\times(\mathbf B\times\mathbf C)
\ne
(\mathbf A\times\mathbf B)\times\mathbf C.
\]

左结合的三重叉积也可展开，但公式不同：

\[
(\mathbf A\times\mathbf B)\times\mathbf C
=\mathbf B(\mathbf A\cdot\mathbf C)
-\mathbf A(\mathbf B\cdot\mathbf C).
\]

它可由右结合公式和叉乘反交换律直接推出：

\[
\begin{aligned}
(\mathbf A\times\mathbf B)\times\mathbf C
&=-\mathbf C\times(\mathbf A\times\mathbf B)\\
&=-\big[\mathbf A(\mathbf C\cdot\mathbf B)
-\mathbf B(\mathbf C\cdot\mathbf A)\big].
\end{aligned}
\]

取
\(\mathbf A=\mathbf i\)、\(\mathbf B=\mathbf i\)、\(\mathbf C=\mathbf j\)，则

\[
\mathbf i\times(\mathbf i\times\mathbf j)=-\mathbf j,
\qquad
(\mathbf i\times\mathbf i)\times\mathbf j=\mathbf0.
\]

因此叉乘不满足结合律。两个公式都应从括号最内层开始识别，不能用同一个口诀机械替换。

\subsection{二、三重积公式的推导}\label{m1:ux4e8cux4e09ux91cdux79efux516cux5f0fux7684ux63a8ux5bfc}

\subsubsection{1.
标量三重积的行列式形式与符号规律}\label{m1:ux6807ux91cfux4e09ux91cdux79efux7684ux884cux5217ux5f0fux5f62ux5f0fux4e0eux7b26ux53f7ux89c4ux5f8b}

由叉乘的分量定义，

\[
\mathbf B\times\mathbf C
=
(B_yC_z-B_zC_y)\mathbf i
+(B_zC_x-B_xC_z)\mathbf j
+(B_xC_y-B_yC_x)\mathbf k.
\]

再与 \(\mathbf A\) 点乘，得到

\[
\begin{aligned}
\mathbf A\cdot(\mathbf B\times\mathbf C)
&=A_x(B_yC_z-B_zC_y)
+A_y(B_zC_x-B_xC_z)\\
&\quad+A_z(B_xC_y-B_yC_x)\\
&=
\begin{vmatrix}
A_x&A_y&A_z\\
B_x&B_y&B_z\\
C_x&C_y&C_z
\end{vmatrix}.
\end{aligned}
\]

循环轮换 \(\mathbf A,\mathbf B,\mathbf C\)
相当于行列式的偶置换，行列式不变；交换任意两个矢量相当于交换两行，行列式变号。这就严格得到标量三重积的轮换对称性与反对称性。特别地，三重积为零当且仅当三矢量线性相关，也就是它们不能张成非零体积。

由点乘与叉乘的不等式还可得到

\[
\big|\mathbf A\cdot(\mathbf B\times\mathbf C)\big|
\le |\mathbf A|\,|\mathbf B\times\mathbf C|
\le |\mathbf A|\,|\mathbf B|\,|\mathbf C|.
\]

对非零矢量，最后一个等号成立当且仅当 \(\mathbf B\perp\mathbf C\)，且
\(\mathbf A\parallel\mathbf B\times\mathbf C\)。这给出了三条棱长固定时平行六面体体积的最大值。

\begin{center}
\includegraphics[width=\linewidth,height=.40\textheight,keepaspectratio]{\detokenize{assets/module-01/tikz-m1-topic1-v3-01-signed-volume.png}}
\par\smallskip{\small 图示：标量三重积等于底面积乘有向高，矢量端点跨过底面时符号改变}
\end{center}

因此，``三重积为零''与``三重积为负''是两种完全不同的情形：前者表示三矢量线性相关，平行六面体退化；后者只表示取向与右手约定相反，体积大小仍然可以非零。

\subsubsection{2.
矢量三重积的分量证明}\label{m1:ux77e2ux91cfux4e09ux91cdux79efux7684ux5206ux91cfux8bc1ux660e}

下面只使用叉乘定义与代数分组。图中标出``补项''在推导中的作用；它只是把已有分量重组为两个点积。

\begin{center}
\includegraphics[width=\linewidth,height=.40\textheight,keepaspectratio]{\detokenize{assets/module-01/tikz-m1-topic1-v3-04-component-proof.png}}
\par\smallskip{\small 图示：矢量三重积的分量证明通过补项和分组识别两个点积}
\end{center}

\textbf{推导过程：}

设目标矢量
\(\mathbf{V} = \mathbf{A} \times (\mathbf{B} \times \mathbf{C})\)。

我们先单独表示出内层叉积 \(\mathbf{D} = \mathbf{B} \times \mathbf{C}\)
的三个直角坐标分量：

\[
D_x = B_y C_z - B_z C_y
\]

\[
D_y = B_z C_x - B_x C_z
\]

\[
D_z = B_x C_y - B_y C_x
\]

接下来计算外层叉积
\(\mathbf{V} = \mathbf{A} \times \mathbf{D}\)。三个分量的结构完全相同，先提取
\(x\) 轴分量 \(V_x\)：

\[
V_x = A_y D_z - A_z D_y
\]

将 \(D_y\) 和 \(D_z\) 的表达式代入上式并展开：

\[
V_x = A_y (B_x C_y - B_y C_x) - A_z (B_z C_x - B_x C_z)
\]

\[
V_x = A_y B_x C_y - A_y B_y C_x - A_z B_z C_x + A_z B_x C_z
\]

此时，运用代数中经典的``加零''技巧：我们在等式右侧同时加上并减去项
\(A_x B_x C_x\)：

\[
V_x = A_y B_x C_y - A_y B_y C_x - A_z B_z C_x + A_z B_x C_z + A_x B_x C_x - A_x B_x C_x
\]

对上述六项重新进行分组，将含有 \(B_x\) 的项提公因式 \(B_x\)，含有
\(C_x\) 的项提公因式 \(-C_x\)：

\[
V_x = B_x (A_x C_x + A_y C_y + A_z C_z) - C_x (A_x B_x + A_y B_y + A_z B_z)
\]

观察括号内的代数式，恰好是点乘的标准坐标展开形式：

\[
A_x C_x + A_y C_y + A_z C_z = \mathbf{A} \cdot \mathbf{C}
\]

\[
A_x B_x + A_y B_y + A_z B_z = \mathbf{A} \cdot \mathbf{B}
\]

代入后得到 \(x\) 分量的最终形式：

\[
V_x = B_x (\mathbf{A} \cdot \mathbf{C}) - C_x (\mathbf{A} \cdot \mathbf{B})
\]

对 \(y,z\) 分量重复完全相同的展开与分组，得到：

\[
V_y = B_y (\mathbf{A} \cdot \mathbf{C}) - C_y (\mathbf{A} \cdot \mathbf{B})
\]

\[
V_z = B_z (\mathbf{A} \cdot \mathbf{C}) - C_z (\mathbf{A} \cdot \mathbf{B})
\]

将三个分量组合回矢量形式，即证明了矢量三重积恒等式：

\[
\mathbf{V} = \mathbf{B}(\mathbf{A} \cdot \mathbf{C}) - \mathbf{C}(\mathbf{A} \cdot \mathbf{B})
\]

证毕。

\subsection{三、磁场中的速度分解}\label{m1:ux4e09ux78c1ux573aux4e2dux7684ux901fux5ea6ux5206ux89e3}

在非零磁场 \(\mathbf B\ne\mathbf0\) 中，任意速度可唯一分解为
\(\mathbf v=\mathbf v_{\parallel}+\mathbf v_{\perp}\)，其中

\[
\mathbf v_{\parallel}=\frac{\mathbf v\cdot\mathbf B}{B^2}\mathbf B,
\qquad
\mathbf v_{\perp}=\mathbf v-\mathbf v_{\parallel}.
\]

于是
\(\mathbf v_{\parallel}\parallel\mathbf B\)、\(\mathbf v_{\perp}\perp\mathbf B\)。由矢量三重积恒等式，

\[
\mathbf B\times(\mathbf v\times\mathbf B)
=\mathbf vB^2-\mathbf B(\mathbf B\cdot\mathbf v)
=B^2\mathbf v_{\perp}.
\]

这说明两次叉乘自动滤去平行于磁场的速度分量，只保留垂直分量。这一结构会在带电粒子回旋运动、磁镜和等离子体漂移中反复出现；若
\(\mathbf v\parallel\mathbf B\)，右端为零，也与叉乘的直接判断一致。若
\(\mathbf B=\mathbf0\)，原式两边都为零，但上述平行、垂直分解不再作为方向分解使用。

令 \(\mathbf b=\mathbf B/B\)
为磁场方向的单位矢量，则上式可写成更直接的投影公式

\[
\boxed{\mathbf b\times(\mathbf v\times\mathbf b)
=\mathbf v-(\mathbf v\cdot\mathbf b)\mathbf b
=\mathbf v_{\perp}}.
\]

因此，``单位矢量在两侧各叉乘一次''构成垂直于该方向的投影算子；平行分量被消去，垂直分量保持不变。若单位条件不满足，则必须保留因子
\(B^2\)，不能把 \(\mathbf B\) 直接当成 \(\mathbf b\)。

\begin{center}
\includegraphics[width=\linewidth,height=.40\textheight,keepaspectratio]{\detokenize{assets/module-01/tikz-m1-topic1-v3-06-double-cross-projection.jpg}}
\par\smallskip{\small 图示：两次叉乘消去平行于磁场的分量并保留垂直分量}
\end{center}

矢量三重积还立即给出\textbf{雅可比恒等式}：

\[
\boxed{
\mathbf A\times(\mathbf B\times\mathbf C)
+\mathbf B\times(\mathbf C\times\mathbf A)
+\mathbf C\times(\mathbf A\times\mathbf B)=\mathbf0 }.
\]

把三个三重积分别展开，得到

\[
\begin{aligned}
&\mathbf B(\mathbf A\cdot\mathbf C)-\mathbf C(\mathbf A\cdot\mathbf B)\\
&+\mathbf C(\mathbf B\cdot\mathbf A)-\mathbf A(\mathbf B\cdot\mathbf C)\\
&+\mathbf A(\mathbf C\cdot\mathbf B)-\mathbf B(\mathbf C\cdot\mathbf A)=\mathbf0,
\end{aligned}
\]

其中六项两两抵消。这个恒等式表达的是三个不同括号项之和为零，不能把它误读成叉乘具有结合律。

\subsection{四、物理图像与方法解释}\label{m1:ux56dbux7269ux7406ux56feux50cfux4e0eux65b9ux6cd5ux89e3ux91ca}

\begin{itemize}
\tightlist
\item
  \textbf{标量三重积的几何意义}：\\
  \(\mathbf{A} \cdot (\mathbf{B} \times \mathbf{C})\) 在几何上代表以
  \(\mathbf{A}\)、\(\mathbf{B}\)、\(\mathbf{C}\)
  为三条共点棱所构成的\textbf{平行六面体的有向体积}。令
  \(\mathbf n=\mathbf B\times\mathbf C\ne\mathbf0\)，并定义沿
  \(\mathbf n\) 的有向高为
  \(h_{\mathbf n}=\mathbf A\cdot\mathbf n/|\mathbf n|\)；则
  \(\mathbf A\cdot(\mathbf B\times\mathbf C)=|\mathbf B\times\mathbf C|h_{\mathbf n}\)。若结果为正，说明三者构成右手系。
\item
  \textbf{矢量三重积的共面约束}：\\
  当 \(\mathbf B,\mathbf C\) 不平行时，\(\mathbf{B} \times \mathbf{C}\)
  垂直于它们张成的平面；\(\mathbf{A} \times (\mathbf{B} \times \mathbf{C})\)
  又垂直于该法向，因此必落在这个平面内。这解释了为什么结果可写成
  \(\mathbf B,\mathbf C\) 的线性组合。若
  \(\mathbf B\parallel\mathbf C\)，内层叉乘为零，最终结果也为零；零矢量仍属于二者张成的子空间。
\end{itemize}

其中 \(h_{\mathbf n}\)
是可正可负的有向高，而不是通常取非负值的几何高度。对矢量三重积，令
\(\mathbf n=\mathbf B\times\mathbf C\)，则结果
\(\mathbf R=\mathbf A\times\mathbf n\) 同时满足

\[
\mathbf R\cdot\mathbf n=0,
\qquad
\mathbf R\in\operatorname{span}\{\mathbf B,\mathbf C\}.
\]

\begin{center}
\includegraphics[width=\linewidth,height=.40\textheight,keepaspectratio]{\detokenize{assets/module-01/tikz-m1-topic1-v3-05-vector-triple-plane.png}}
\par\smallskip{\small 图示：内层叉积给出平面法向，外层叉积的结果回到 B、C 张成的平面}
\end{center}

\begin{itemize}
\item
  \textbf{力学中的应用（向心加速度）}：\\
  对绕固定轴作匀角速度转动的质点，取原点在转轴上，位置矢量为
  \(\mathbf r\)，并记 \(\omega=|\boldsymbol\omega|\)，则
  \(\mathbf v=\boldsymbol\omega\times\mathbf r\)。其法向加速度为

  \[
  \mathbf{a}_n = \boldsymbol{\omega} \times \mathbf{v} = \boldsymbol{\omega} \times (\boldsymbol{\omega} \times \mathbf{r})
  \]
\end{itemize}

根据展开公式，\(\mathbf{a}_n = \boldsymbol{\omega}(\boldsymbol{\omega} \cdot \mathbf{r}) - \mathbf{r}\omega^2=-\omega^2\mathbf r_{\perp}\)，其中
\(\mathbf r_{\perp}\) 是垂直转轴的位矢分量。当 \(r_{\perp}>0\) 且
\(\omega\ne0\) 时，它严格指向转轴；若 \(r_{\perp}=0\) 或
\(\omega=0\)，则
\(\mathbf a_n=\mathbf0\)，零矢量没有方向。若转轴方向固定而角速度大小随时间变化，则总加速度还应加上切向项
\(\dot{\boldsymbol\omega}\times\mathbf r\)。

\begin{center}
\includegraphics[width=\linewidth,height=.40\textheight,keepaspectratio]{\detokenize{assets/module-01/tikz-m1-topic1-v3-07-rotation-acceleration.png}}
\par\smallskip{\small 图示：两次与角速度叉乘只保留垂直转轴的位矢分量并使方向反转}
\end{center}

\subsection{五、变式、使用条件与常见错误}\label{m1:ux4e94ux53d8ux5f0fux4f7fux7528ux6761ux4ef6ux4e0eux5e38ux89c1ux9519ux8bef}

\begin{enumerate}
\def\labelenumi{\arabic{enumi}.}
\tightlist
\item
  \textbf{先读最外层运算符。}
  \(\mathbf A\cdot(\mathbf B\times\mathbf C)\)
  的结果是标量，\(\mathbf A\times(\mathbf B\times\mathbf C)\)
  的结果是矢量；它们不能互相替代。
\item
  \textbf{叉乘没有结合律。} 不能去掉括号或改成
  \((\mathbf A\times\mathbf B)\times\mathbf C\)；只有在具体计算后偶然相等时才可得到同一结果。
\item
  \textbf{共面结论有前提。} 当 \(\mathbf B\times\mathbf C\ne\mathbf0\)
  时，结果落在 \(\mathbf B,\mathbf C\)
  张成的平面内；平行的退化情形应直接先判内层叉乘为零。
\item
  \textbf{不要把 \(\nabla\) 当作普通矢量。}
  对普通矢量的三重积恒等式不能直接把 \(\nabla\)
  填进去；一旦含有微分算子，必须说明它作用在哪个因子上，并使用后续专题的分量或乘积法则。
\item
  \textbf{量纲检查。} 等式右边两项必须具有同一量纲；例如
  \(\mathbf B(\mathbf A\cdot\mathbf C)\) 与
  \(\mathbf C(\mathbf A\cdot\mathbf B)\)
  都是三个矢量量纲的乘积。这能发现漏写点乘或误换叉乘顺序的错误。
\item
  \textbf{左右结合使用不同展开式。}
  \(\mathbf A\times(\mathbf B\times\mathbf C)\) 的结果属于
  \(\operatorname{span}\{\mathbf B,\mathbf C\}\)，而
  \((\mathbf A\times\mathbf B)\times\mathbf C\) 的结果属于
  \(\operatorname{span}\{\mathbf A,\mathbf B\}\)。
\item
  \textbf{零矢量没有方向。} 当三重积退化为 \(\mathbf0\)
  时，可以说它属于相应子空间，但不能再说它``沿着''某个确定方向。
\end{enumerate}

\subsection{六、练习与解答}\label{m1:ux516dux7ec3ux4e60ux4e0eux89e3ux7b54}

作业按``标量三重积---投影恒等式---角动量应用''递进。每题先判断结果类型和几何约束，再进行分量或恒等式计算。

\par\addvspace{.4\baselineskip}
\noindent\begin{minipage}{\linewidth}

\subsubsection{作业一：有向体积、符号与线性相关性}\label{m1:ux4f5cux4e1aux4e00ux6709ux5411ux4f53ux79efux7b26ux53f7ux4e0eux7ebfux6027ux76f8ux5173ux6027}

\textbf{题号：M1-04}

\textbf{层级：} 基础　　\textbf{方法：} 矢量分析

在固定的右手直角坐标系中，把三个矢量的起点平移到同一原点，分量均相对于该单位正交基给出。

给定

\[
\mathbf A=(1,2,0),\qquad
\mathbf B=(1,0,1),\qquad
\mathbf C=(0,1,1).
\]

\begin{enumerate}
\def\labelenumi{(\arabic{enumi})}
\item
  计算 \(\mathbf A\cdot(\mathbf B\times\mathbf C)\)。
\item
  求三矢量张成的平行六面体体积，并判断 \(\mathbf A,\mathbf B,\mathbf C\)
  的取向。
\item
  交换 \(\mathbf B\) 与 \(\mathbf C\)
  后结果如何变化？三矢量是否线性相关？
\end{enumerate}


\end{minipage}\par

\par\noindent\begin{minipage}{\linewidth}
\textbf{解答}

\begin{center}
\includegraphics[width=\linewidth,height=.40\textheight,keepaspectratio]{\detokenize{assets/module-01/examples/M1-04.pdf}}
\par\smallskip{\small 图 M1-04：有向体积与取向}
\end{center}
\end{minipage}\par

图 M1-04左图用 \(\mathbf B,\mathbf C\) 定义底面，其正法向为
\(\mathbf B\times\mathbf C=(-1,-1,1)\)。将图中的 \(\mathbf A=(1,2,0)\)
投影到这个法向，得到有向体积
\(-3\)；负号表示端点位于底面负侧，取绝对值得几何体积
\(3\)。交换底面的两条边只反转法向，以下行列式计算也应相应变号。

先计算

\[
\mathbf B\times\mathbf C
=
\begin{vmatrix}
\mathbf i&\mathbf j&\mathbf k\\
1&0&1\\
0&1&1
\end{vmatrix}
=(-1,-1,1).
\]

因此

\[
\mathbf A\cdot(\mathbf B\times\mathbf C)
=(1,2,0)\cdot(-1,-1,1)=-3.
\]

平行六面体的普通体积是有向体积的绝对值，故 \(V=3\)。结果为负，说明
\(\mathbf A,\mathbf B,\mathbf C\) 的取向与所选右手坐标系相反。交换
\(\mathbf B,\mathbf C\) 相当于交换行列式两行，结果变为
\(+3\)。由于三重积非零，三矢量线性无关。

\par\addvspace{.4\baselineskip}
\noindent\begin{minipage}{\linewidth}

\subsubsection{作业二：两次叉乘与正交投影}\label{m1:ux4f5cux4e1aux4e8cux4e24ux6b21ux53c9ux4e58ux4e0eux6b63ux4ea4ux6295ux5f71}

\textbf{题号：M1-05}

\textbf{层级：} 基础　　\textbf{方法：} 矢量分析

在三维欧氏空间中，\(\mathbf A\) 为任意自由矢量；过其起点取法向为
\(\mathbf n\) 的平面，且 \(|\mathbf n|=1\)。

设 \(\mathbf n\) 为单位矢量，定义

\[
P_{\perp}(\mathbf A)=\mathbf n\times(\mathbf A\times\mathbf n).
\]

\begin{enumerate}
\def\labelenumi{(\arabic{enumi})}
\item
  证明
  \(P_{\perp}(\mathbf A)=\mathbf A-(\mathbf A\cdot\mathbf n)\mathbf n\)。
\item
  证明 \(P_{\perp}(\mathbf A)\perp\mathbf n\)，并分别讨论
  \(\mathbf A\parallel\mathbf n\) 与 \(\mathbf A\perp\mathbf n\)。
\item
  证明投影的幂等性：\(P_{\perp}(P_{\perp}(\mathbf A))=P_{\perp}(\mathbf A)\)。
\end{enumerate}


\end{minipage}\par

\par\noindent\begin{minipage}{\linewidth}
\textbf{解答}

\begin{center}
\includegraphics[width=\linewidth,height=.40\textheight,keepaspectratio]{\detokenize{assets/module-01/examples/M1-05.pdf}}
\par\smallskip{\small 图 M1-05：两次叉乘与正交投影}
\end{center}
\end{minipage}\par

图 M1-05中先把 \(\mathbf A\) 分成沿单位法向 \(\mathbf n\)
的橙色分量与垂直于它的绿色分量。两次叉乘消去沿 \(\mathbf n\)
的分量，留下
\(\mathbf A-(\mathbf A\cdot\mathbf n)\mathbf n\)。绿色结果已经在
\(\mathbf n^{\perp}\)
内，故再次作同一投影不变；下面的代数证明正对应这条图示分解。

由矢量三重积公式和 \(\mathbf n\cdot\mathbf n=1\)，

\[
P_{\perp}(\mathbf A)
=\mathbf A(\mathbf n\cdot\mathbf n)
-\mathbf n(\mathbf n\cdot\mathbf A)
=\mathbf A-(\mathbf A\cdot\mathbf n)\mathbf n.
\]

与 \(\mathbf n\) 点乘得

\[
\mathbf n\cdot P_{\perp}(\mathbf A)
=\mathbf n\cdot\mathbf A
-(\mathbf A\cdot\mathbf n)(\mathbf n\cdot\mathbf n)=0,
\]

故结果垂直于 \(\mathbf n\)。若
\(\mathbf A\parallel\mathbf n\)，结果为零；若
\(\mathbf A\perp\mathbf n\)，结果等于 \(\mathbf A\)。最后，因为
\(P_{\perp}(\mathbf A)\) 已垂直于
\(\mathbf n\)，再次投影不会改变它，所以

\[
P_{\perp}(P_{\perp}(\mathbf A))=P_{\perp}(\mathbf A).
\]

\par\addvspace{.4\baselineskip}
\noindent\begin{minipage}{\linewidth}

\subsubsection{作业三：定轴转动中单个质点的角动量}\label{m1:ux4f5cux4e1aux4e09ux5b9aux8f74ux8f6cux52a8ux4e2dux5355ux4e2aux8d28ux70b9ux7684ux89d2ux52a8ux91cf}

\textbf{题号：M1-06}

\textbf{层级：} 基础　　\textbf{方法：} 矢量分析、角动量

以下只研究刚体上这一质量为 \(m>0\) 的质点；角动量一律相对于所给轴上原点
\(O\) 定义。

考虑一个质量为 \(m\) 的质点。取原点 \(O\)
位于固定转轴上，刚体以恒定角速度 \(\boldsymbol{\omega}\) 绕该轴转动，记
\(\omega=|\boldsymbol\omega|\)，质点相对于 \(O\) 的位置矢量为
\(\mathbf{r}\)。

\begin{enumerate}
\def\labelenumi{(\arabic{enumi})}
\item
  请写出该质点相对于原点 \(O\) 的角动量 \(\mathbf{L}\) 的矢量表达式。
\item
  利用矢量三重积展开公式，将 \(\mathbf{L}\) 表示为
  \(\boldsymbol{\omega}\) 和 \(\mathbf{r}\) 的函数；证明当
  \(\mathbf r\perp\boldsymbol\omega\) 且 \(r>0\)、\(\omega>0\)
  时，\(\mathbf L\) 与 \(\boldsymbol\omega\) 同向，并说明退化情形。
\end{enumerate}


\end{minipage}\par

\par\noindent\begin{minipage}{\linewidth}
\textbf{解答}

\begin{center}
\includegraphics[width=\linewidth,height=.40\textheight,keepaspectratio]{\detokenize{assets/module-01/examples/M1-06.pdf}}
\par\smallskip{\small 图 M1-06：定轴转动质点的角动量}
\end{center}
\end{minipage}\par

图 M1-06左图的 \(\mathbf r\) 含轴向与横向两部分，速度只来自横向位置
\(\mathbf r_{\perp}\)。中图沿转轴观察可直接读出
\(\mathbf v\perp\mathbf r_{\perp}\)。计算 \(\mathbf r\times m\mathbf v\)
时仍须使用完整的 \(\mathbf r\)，因此一般出现横向角动量；只有
\(\mathbf r\perp\boldsymbol\omega\) 时，右图公式中的第二项才消失。

\textbf{第 1 步：由基本定义写出角动量表达式}

质点相对于原点的角动量定义为位置矢量与动量的叉积：

\[
\mathbf{L} = \mathbf{r} \times \mathbf{p} = \mathbf{r} \times (m\mathbf{v})
\]

在定轴转动刚体中，质点的线速度与角速度的关系为：

\[
\mathbf{v} = \boldsymbol{\omega} \times \mathbf{r}
\]

代入角动量表达式，提取标量质量 \(m\)，得到包含矢量三重积的方程：

\[
\mathbf{L} = m[\mathbf{r} \times (\boldsymbol{\omega} \times \mathbf{r})]
\]

\textbf{第 2 步：利用矢量三重积恒等式展开三重积}

应用矢量三重积恒等式
\(\mathbf{A} \times (\mathbf{B} \times \mathbf{C}) = \mathbf{B}(\mathbf{A} \cdot \mathbf{C}) - \mathbf{C}(\mathbf{A} \cdot \mathbf{B})\)。

在此题中，对应关系为：\(\mathbf{A} \to \mathbf{r}\)，\(\mathbf{B} \to \boldsymbol{\omega}\)，\(\mathbf{C} \to \mathbf{r}\)。展开得：

\[
\mathbf{r} \times (\boldsymbol{\omega} \times \mathbf{r}) = \boldsymbol{\omega}(\mathbf{r} \cdot \mathbf{r}) - \mathbf{r}(\mathbf{r} \cdot \boldsymbol{\omega})
\]

由于 \(\mathbf{r} \cdot \mathbf{r} = r^2\)（其中 \(r\)
为质点到原点的距离），代回角动量表达式：

\[
\mathbf{L} = m [ r^2 \boldsymbol{\omega} - (\mathbf{r} \cdot \boldsymbol{\omega}) \mathbf{r} ] \tag{1}
\]

这便是定轴转动中单个质点相对于轴上固定原点 \(O\) 的角动量表达式。

\textbf{第 3 步：分析正交条件下的方向关系}

先讨论非退化情形 \(r>0\)、\(\omega>0\)。当位置矢量 \(\mathbf{r}\)
与角速度矢量 \(\boldsymbol{\omega}\) 垂直时，两矢量的点积为零：

\[
\mathbf{r} \cdot \boldsymbol{\omega} = 0
\]

将此条件代入上一步得到的式 (1) 中，等式右侧第二项消失，方程退化为：

\[
\mathbf{L} = m r^2 \boldsymbol{\omega}
\]

由于 \(m>0\)、\(r>0\)，比例系数 \(mr^2\) 为正，且已假定
\(\omega>0\)，所以 \(\mathbf L\) 非零并与 \(\boldsymbol\omega\)
同向；\(mr^2\) 正是该质点绕该轴的转动惯量 \(I\)。若 \(r=0\) 或
\(\omega=0\)，则 \(\mathbf L=\mathbf0\)，不能谈它的方向。当
\(\mathbf r\) 同时含有平行与垂直转轴的分量时，\(\mathbf L\) 一般不与
\(\boldsymbol\omega\) 共线；若
\(\mathbf r\parallel\boldsymbol\omega\)，则
\(\mathbf v=\mathbf0\)、\(\mathbf L=\mathbf0\)。这正是高等力学中引入惯量张量的物理起点。

\section{专题二：纳布拉算符、梯度、散度和旋度（基础必讲）}\label{m1:ux4e13ux9898ux4e8cux7eb3ux5e03ux62c9ux7b97ux7b26ux68afux5ea6ux6563ux5ea6ux548cux65cbux5ea6ux57faux7840ux5fc5ux8bb2}

\begin{quote}
阅读准备：偏导数只改变一个坐标 对 \(f(x,y,z)\) 求
\(\partial f/\partial x\) 时，先把 \(y,z\) 固定，只观察沿 \(x\)
方向的变化率。偏导数仍服从一元函数的乘积法则和链式法则；区别只在于场同时依赖多个坐标。
\end{quote}

\par\addvspace{.4\baselineskip}
\noindent\begin{minipage}{\linewidth}

\textbf{例 M1-07：两个坐标方向的偏导数}

在直角坐标系中给定标量场
\(f(x,y,z)=x^2y+z\)。分别在保持其余两个坐标不变的条件下，求
\(\partial f/\partial x\) 和
\(\partial f/\partial y\)，并说明它们各自衡量沿哪个坐标方向的局部变化率。判断同一位置的这两个变化率是否必须相同。


\end{minipage}\par

\par\noindent\begin{minipage}{\linewidth}
\textbf{解答}

\begin{center}
\includegraphics[width=\linewidth,height=.40\textheight,keepaspectratio]{\detokenize{assets/module-01/examples/M1-07.pdf}}
\par\smallskip{\small 图 M1-07：两个坐标方向的偏导数}
\end{center}
\end{minipage}\par

图 M1-07上图把 \(y,z\) 固定，读取关于 \(x\) 的抛物线切片，局部斜率为
\(2xy\)；下图把 \(x,z\) 固定，读取关于 \(y\) 的直线切片，斜率为
\(x^2\)。两图的横坐标分别是 \(x\) 与
\(y\)，不能把它们当作同一条曲线的两个斜率。图上的正参数切片只是示意，下面的偏导公式对整个场成立。

若 \(f=x^2y+z\)，则 \(\partial f/\partial x=2xy\)，而
\(\partial f/\partial y=x^2\)。同一个场沿不同坐标方向可以有不同变化率。

纳布拉算符 \(\nabla\)
由各坐标方向的一阶偏导数组成。它作用于标量场或矢量场时，分别构成梯度、散度和旋度；乘积法则及二阶恒等式还取决于算符的作用对象。专题三至五将分别展开三种一阶运算。

\subsection{一、作用对象、定义与输出类型}\label{m1:ux4e00ux4f5cux7528ux5bf9ux8c61ux5b9aux4e49ux4e0eux8f93ux51faux7c7bux578b}

设 \(U\subset\mathbb R^3\) 为开区域。在固定的右手直角坐标系中，

\[
\boxed{\nabla
=\mathbf i\frac{\partial}{\partial x}
+\mathbf j\frac{\partial}{\partial y}
+\mathbf k\frac{\partial}{\partial z}}.
\]

\(\nabla\)
是微分算符，不是脱离作用对象后可以自由移动的普通矢量。它的量纲为长度的倒数；若场量
\(F\) 的量纲为 \([F]\)，则一次空间微分的量纲为
\([F]/\mathrm L\)。上式还依赖``直角坐标基矢不随位置变化''这一事实；柱坐标和球坐标中必须使用含尺度因子的专门公式。

若 \(f\in C^1(U)\) 是标量场，\(\mathbf A=(A_x,A_y,A_z)\in C^1(U)\)
是矢量场，则三个基本一阶运算为

{\def\LTcaptype{none} % do not increment counter
\begin{longtable}[]{@{}
  >{\raggedright\arraybackslash}p{(\linewidth - 6\tabcolsep) * \real{0.2500}}
  >{\raggedright\arraybackslash}p{(\linewidth - 6\tabcolsep) * \real{0.2500}}
  >{\raggedright\arraybackslash}p{(\linewidth - 6\tabcolsep) * \real{0.2500}}
  >{\raggedright\arraybackslash}p{(\linewidth - 6\tabcolsep) * \real{0.2500}}@{}}
\toprule\noalign{}
\begin{minipage}[b]{\linewidth}\raggedright
运算
\end{minipage} & \begin{minipage}[b]{\linewidth}\raggedright
输入
\end{minipage} & \begin{minipage}[b]{\linewidth}\raggedright
定义
\end{minipage} & \begin{minipage}[b]{\linewidth}\raggedright
输出
\end{minipage} \\
\midrule\noalign{}
\endhead
\bottomrule\noalign{}
\endlastfoot
梯度 & 标量场 \(f\) & \(\nabla f=(\partial_xf,\partial_yf,\partial_zf)\)
& 矢量场 \\
散度 & 矢量场 \(\mathbf A\) &
\(\nabla\cdot\mathbf A=\partial_xA_x+\partial_yA_y+\partial_zA_z\) &
标量场 \\
旋度 & 矢量场 \(\mathbf A\) & \(\nabla\times\mathbf A\) & 矢量场 \\
\end{longtable}
}

旋度的分量定义为

\[
\boxed{
\nabla\times\mathbf A=
\left(
\frac{\partial A_z}{\partial y}-\frac{\partial A_y}{\partial z},
\frac{\partial A_x}{\partial z}-\frac{\partial A_z}{\partial x},
\frac{\partial A_y}{\partial x}-\frac{\partial A_x}{\partial y}
\right)}.
\]

常见的``行列式写法''只是记忆三个分量符号的工具，不是由普通数构成的行列式，不能据此随意交换行或把
\(\nabla\) 移过场量。

若 \(f,\mathbf A\in C^2(U)\)，直角坐标中的拉普拉斯算符为

\[
\nabla^2
=\frac{\partial^2}{\partial x^2}
+\frac{\partial^2}{\partial y^2}
+\frac{\partial^2}{\partial z^2},
\]

并规定

\[
\nabla^2f=\partial_x^2f+\partial_y^2f+\partial_z^2f,
\qquad
\nabla^2\mathbf A
=(\nabla^2A_x,\nabla^2A_y,\nabla^2A_z).
\]

因此，标量拉普拉斯的输出仍是标量场，矢量拉普拉斯的输出仍是矢量场。

\begin{center}
\includegraphics[width=\linewidth,height=.40\textheight,keepaspectratio]{\detokenize{assets/module-01/tikz-m1-topic2-v4-01-operator-map-safe-v7.png}}
\par\smallskip{\small 图示：先根据输入对象和连接符号判断纳布拉运算的输出类型}
\end{center}

\textbf{四个最小检验}

\begin{enumerate}
\def\labelenumi{\arabic{enumi}.}
\tightlist
\item
  对 \(f=x^2+2y-z\)，有 \(\nabla f=(2x,2,-1)\)，输出为矢量。
\item
  对 \(\mathbf A=(x,-y,0)\)，有
  \(\nabla\cdot\mathbf A=1-1=0\)。散度为零不等于场为零。
\item
  对 \(\mathbf C=(-y/2,x/2,0)\)，有
  \(\nabla\times\mathbf C=\mathbf k\)，输出方向垂直于 \(xy\) 平面。
\item
  对 \(f=x^2+y^2+z^2\)，有 \(\nabla^2f=6\)，输出为标量。
\end{enumerate}

\subsection{二、三种一阶运算的局部几何意义}\label{m1:ux4e8cux4e09ux79cdux4e00ux9636ux8fd0ux7b97ux7684ux5c40ux90e8ux51e0ux4f55ux610fux4e49}

\subsubsection{1.
梯度：方向导数的最大值}\label{m1:ux68afux5ea6ux65b9ux5411ux5bfcux6570ux7684ux6700ux5927ux503c}

设 \(\mathbf u\) 为单位矢量。标量场 \(f\) 在点 \(P\) 沿 \(\mathbf u\)
的方向导数为

\[
D_{\mathbf u}f(P)=\nabla f(P)\cdot\mathbf u.
\]

由柯西不等式，

\[
|D_{\mathbf u}f(P)|\le |\nabla f(P)|,
\]

当 \(\nabla f(P)\ne\mathbf0\) 时，方向导数的最大值在 \(\mathbf u\) 与
\(\nabla f(P)\)
同向时取得。因此梯度指向一阶变化最快的方向，其模长就是最大的方向导数。若
\(\mathrm d\mathbf r_t\) 沿等值面切向，则

\[
\mathrm df=\nabla f\cdot\mathrm d\mathbf r_t=0,
\qquad
\nabla f\perp\mathrm d\mathbf r_t.
\]

当 \(\nabla f(P)=\mathbf0\)
时，所有一阶方向导数都为零，不能再由梯度确定``最陡方向''；该点可能是极值点，也可能是鞍点或更高阶退化点。

\subsubsection{2.
散度：单位体积的净流出极限}\label{m1:ux6563ux5ea6ux5355ux4f4dux4f53ux79efux7684ux51c0ux6d41ux51faux6781ux9650}

对在 \(P\) 附近连续可微的矢量场，取一族规则缩向 \(P\) 的小体积
\(V\)，闭合曲面 \(\partial V\) 统一取外法向，定义

\[
\boxed{
\nabla\cdot\mathbf A(P)
=\lim_{V\to P}
\frac{1}{\Delta V}
\oint_{\partial V}\mathbf A\cdot\mathrm d\mathbf S }.
\]

正散度表示局部净流出，负散度表示局部净流入，零散度只表示一阶净通量为零。它既不要求
\(\mathbf A=\mathbf0\)，也不排除场线弯曲。

\par\addvspace{.4\baselineskip}
\noindent\begin{minipage}{\linewidth}

\textbf{例 M1-08：线性径向场的小盒通量}

空间中给定线性径向矢量场
\(\mathbf A=a(x\mathbf e_x+y\mathbf e_y+z\mathbf e_z)\)，其中 \(a\)
为常量。在任意一点附近取边长为
\(\Delta x,\Delta y,\Delta z\)、各边平行于坐标轴的小长方体，所有面元均取体积外法向。求
\(\nabla\cdot\mathbf A\)；分别计算三对相对面的净通量，并用总通量除以体积核对散度的局部定义。


\end{minipage}\par

\par\noindent\begin{minipage}{\linewidth}
\textbf{解答}

\begin{center}
\includegraphics[width=\linewidth,height=.40\textheight,keepaspectratio]{\detokenize{assets/module-01/examples/M1-08.pdf}}
\par\smallskip{\small 图 M1-08：线性径向场的小盒通量}
\end{center}
\end{minipage}\par

图 M1-08中的径向箭头取 \(a>0\)
作示意。按小盒的三个坐标方向配对，每一对面的外向通量差都是
\(a\Delta V\)；六面相加给 \(3a\Delta V\)。所以散度为
\(3a\)，而不是单个截面读出的 \(a\)；\(a<0\)
时全部箭头反向，结果随之变号。

对 \(\mathbf A=a(x,y,z)\)，有
\(\nabla\cdot\mathbf A=3a\)。任取以原点为中心的小长方体，六个面的总外向通量除以体积均为
\(3a\)。

\subsubsection{3.
旋度：单位面积的定向环量极限}\label{m1:ux65cbux5ea6ux5355ux4f4dux9762ux79efux7684ux5b9aux5411ux73afux91cfux6781ux9650}

取通过 \(P\)、法向为单位矢量 \(\hat{\mathbf n}\) 的小平面片 \(S\)。边界
\(\partial S\) 的正向由右手定则与 \(\hat{\mathbf n}\) 绑定，则

\[
\boxed{
\hat{\mathbf n}\cdot(\nabla\times\mathbf A)(P)
=\lim_{S\to P}
\frac{1}{\Delta S}
\oint_{\partial S}\mathbf A\cdot\mathrm d\mathbf l }.
\]

这个极限给出旋度沿 \(\hat{\mathbf n}\)
的分量；对不同法向重复判断即可确定旋度矢量。若改变法向，边界正向也必须同时反向。对光滑速度场，局部自转角速度矢量为
\(\frac12\nabla\times\mathbf v\)。

\par\addvspace{.4\baselineskip}
\noindent\begin{minipage}{\linewidth}

\textbf{例 M1-09：刚性转动场的旋度}

流体的二维速度场为
\(\mathbf v(x,y,z)=-\Omega y\mathbf e_x+\Omega x\mathbf e_y\)，其中
\(\Omega\) 为不随位置变化的常量，速度与 \(z\)
无关。求其旋度及方向；将此速度场视为绕 \(z\)
轴的刚体转动，说明旋度与角速度矢量之间的关系，以及 \(\Omega\)
变号时方向如何改变。


\end{minipage}\par

\par\noindent\begin{minipage}{\linewidth}
\textbf{解答}

\begin{center}
\includegraphics[width=\linewidth,height=.40\textheight,keepaspectratio]{\detokenize{assets/module-01/examples/M1-09.pdf}}
\par\smallskip{\small 图 M1-09：刚性转动场的旋度}
\end{center}
\end{minipage}\par

图 M1-09中的切向箭头绕 \(z\) 轴转动，逆时针对应纸外正法向。\(x\)、\(y\)
两方向的速度梯度各贡献 \(\Omega\)，相加后旋度是
\(2\Omega\mathbf e_z\)。这里读的是单位面积环量；刚体的角速度只有它的一半。

刚性转动速度场 \(\mathbf v=(-\Omega y,\Omega x,0)\) 满足

\[
\nabla\times\mathbf v=2\Omega\mathbf k,
\]

所以局部角速度为 \(\Omega\mathbf k\)，与直接的刚体转动图像一致。

在真空静电学的光滑区域中，三种运算构成一条清晰的关系链：

\[
\mathbf E=-\nabla\varphi,\qquad
\nabla\cdot\mathbf E=\frac{\rho}{\varepsilon_0},\qquad
\nabla\times\mathbf E=\mathbf0.
\]

对第一式取散度并代入第二式，得到泊松方程

\[
\boxed{\nabla^2\varphi=-\frac{\rho}{\varepsilon_0}},
\]

在不含电荷的光滑区域内进一步化为拉普拉斯方程
\(\nabla^2\varphi=0\)。这些关系是静电学关系；当磁场随时间变化时，法拉第定律给出

\[
\nabla\times\mathbf E=-\frac{\partial\mathbf B}{\partial t},
\]

此时电场一般不能在整个区域内写成单值标量势的负梯度。

\subsection{三、乘积法则与二阶恒等式}\label{m1:ux4e09ux4e58ux79efux6cd5ux5219ux4e0eux4e8cux9636ux6052ux7b49ux5f0f}

\subsubsection{1.
线性与乘积法则}\label{m1:ux7ebfux6027ux4e0eux4e58ux79efux6cd5ux5219}

对常数 \(a,b\) 和同类型的场，\(\nabla\)
的梯度、散度、旋度均满足线性关系。例如

\[
\nabla(af+bg)=a\nabla f+b\nabla g,
\qquad
\nabla\cdot(a\mathbf A+b\mathbf B)
=a\nabla\cdot\mathbf A+b\nabla\cdot\mathbf B.
\]

若 \(f,g\in C^1(U)\)、\(\mathbf A\in C^1(U)\)，则

\[
\boxed{\nabla(fg)=f\nabla g+g\nabla f},
\]

\[
\boxed{\nabla\cdot(f\mathbf A)
=f(\nabla\cdot\mathbf A)+\mathbf A\cdot\nabla f},
\]

\[
\boxed{\nabla\times(f\mathbf A)
=f(\nabla\times\mathbf A)+(\nabla f)\times\mathbf A}.
\]

这些公式的本质是一元乘积法则在三个坐标方向上逐项使用。以散度为例，

\[
\begin{aligned}
\nabla\cdot(f\mathbf A)
&=\partial_x(fA_x)+\partial_y(fA_y)+\partial_z(fA_z)\\
&=f(\partial_xA_x+\partial_yA_y+\partial_zA_z)
+A_x\partial_xf+A_y\partial_yf+A_z\partial_zf.
\end{aligned}
\]

旋度的 \(x\) 分量同样给出

\[
\begin{aligned}
[\nabla\times(f\mathbf A)]_x
&=\partial_y(fA_z)-\partial_z(fA_y)\\
&=f(\partial_yA_z-\partial_zA_y)
+(\partial_yf)A_z-(\partial_zf)A_y,
\end{aligned}
\]

其中后两项恰为
\([(\nabla f)\times\mathbf A]_x\)；另两个分量循环轮换即可。

\par\addvspace{.4\baselineskip}
\noindent\begin{minipage}{\linewidth}

\textbf{例 M1-10：标量因子变化产生的散度}

在右手直角坐标系中，标量场为 \(f(x,y,z)=x\)，矢量场为常矢量
\(\mathbf A=\mathbf i\)。分别将乘积 \(f\mathbf A\)
展开成直角坐标分量直接求散度，并使用散度乘积法则再算一次。比较两种结果，说明非零散度来自标量因子的空间变化还是常矢量场本身。


\end{minipage}\par

\par\noindent\begin{minipage}{\linewidth}
\textbf{解答}

\begin{center}
\includegraphics[width=\linewidth,height=.40\textheight,keepaspectratio]{\detokenize{assets/module-01/examples/M1-10.pdf}}
\par\smallskip{\small 图 M1-10：标量因子变化产生的散度}
\end{center}
\end{minipage}\par

图 M1-10中 \(\mathbf A=\mathbf e_x\) 本身均匀，改变的是标量权重
\(f=x\)；右半轴上乘积场的箭头随位置增长。对乘积取散度会读出
\(\partial_x x=1\)。用分配法计算时，均匀场的散度项为零，全部贡献来自沿场方向的标量变化
\(\mathbf A\cdot\nabla f=1\)。

取 \(f=x\)、\(\mathbf A=\mathbf i\)。直接计算
\(\nabla\cdot(x\mathbf i)=1\)；乘积法则也给出

\[
x(\nabla\cdot\mathbf i)+\mathbf i\cdot\nabla x=0+1=1.
\]

\subsubsection{2.
双重旋度恒等式（拔高）}\label{m1:ux53ccux91cdux65cbux5ea6ux6052ux7b49ux5f0fux62d4ux9ad8}

设
\(\mathbf A\in C^2(U)\)。二阶连续可微保证混合偏导可以交换。目标恒等式为

\[
\boxed{
\nabla\times(\nabla\times\mathbf A)
=\nabla(\nabla\cdot\mathbf A)-\nabla^2\mathbf A }.
\]

先证明 \(x\) 分量。由旋度定义，

\[
\begin{aligned}
[\nabla\times(\nabla\times\mathbf A)]_x
&=\partial_y(\partial_xA_y-\partial_yA_x)
-\partial_z(\partial_zA_x-\partial_xA_z)\\
&=\partial_x(\partial_yA_y+\partial_zA_z)
-\partial_y^2A_x-\partial_z^2A_x.
\end{aligned}
\]

在右侧加上并减去 \(\partial_x^2A_x\)，得到

\[
\begin{aligned}
[\nabla\times(\nabla\times\mathbf A)]_x
&=\partial_x(\partial_xA_x+\partial_yA_y+\partial_zA_z)\\
&\quad-(\partial_x^2+\partial_y^2+\partial_z^2)A_x\\
&=\partial_x(\nabla\cdot\mathbf A)-\nabla^2A_x.
\end{aligned}
\]

按 \(x\to y\to z\to x\)
循环置换，可得另外两个分量，合并即得到目标恒等式。

\begin{center}
\includegraphics[width=\linewidth,height=.40\textheight,keepaspectratio]{\detokenize{assets/module-01/tikz-m1-topic2-v4-05-curl-curl-route.png}}
\par\smallskip{\small 图示：双重旋度证明依次经过取一个分量、展开、补项、重组和循环推广}
\end{center}

\par\addvspace{.4\baselineskip}
\noindent\begin{minipage}{\linewidth}

\textbf{例 M1-11：双重旋度恒等式的分量核对}

给定光滑矢量场 \(\mathbf A(x,y,z)=x^2\mathbf e_y\)。分别计算
\(\nabla\times\mathbf A\)、\(\nabla\times(\nabla\times\mathbf A)\)、\(\nabla(\nabla\cdot\mathbf A)\)，以及对各分量逐一作用的矢量拉普拉斯
\(\nabla^2\mathbf A\)。据此核对双重旋度恒等式两边的分量和符号，注意保持微分运算的先后次序。


\end{minipage}\par

\par\noindent\begin{minipage}{\linewidth}
\textbf{解答}

\begin{center}
\includegraphics[width=\linewidth,height=.40\textheight,keepaspectratio]{\detokenize{assets/module-01/examples/M1-11.pdf}}
\par\smallskip{\small 图 M1-11：双重旋度恒等式的分量核对}
\end{center}
\end{minipage}\par

图 M1-11中的场箭头始终沿 \(+y\)，但其大小为 \(x^2\)；因此对 \(y\)
求导并不产生散度。沿 \(x\) 的变化使第一次旋度成为
\(2x\mathbf e_z\)，再求一次旋度时，\(y\) 分量带上负的 \(x\) 导数，得到
\(-2\mathbf e_y\)。这与右图的 \(-\nabla^2\mathbf A\)
相符，符号不能只凭抛物线开口判断。

取 \(\mathbf A=(0,x^2,0)\)，则

\[
\nabla\cdot\mathbf A=0,\qquad
\nabla^2\mathbf A=(0,2,0),
\]

所以恒等式右侧为 \((0,-2,0)\)。另一方面，

\[
\nabla\times\mathbf A=(0,0,2x),\qquad
\nabla\times(\nabla\times\mathbf A)=(0,-2,0).
\]

两边一致。双重旋度不等于矢量拉普拉斯；两者还相差梯度项
\(\nabla(\nabla\cdot\mathbf A)\)。

\subsection{四、运算顺序与使用条件}\label{m1:ux56dbux8fd0ux7b97ux987aux5e8fux4e0eux4f7fux7528ux6761ux4ef6}

\begin{enumerate}
\def\labelenumi{\arabic{enumi}.}
\tightlist
\item
  \textbf{先判输入和输出。} \(\nabla f\)、\(\nabla\times\mathbf A\)
  是矢量；\(\nabla\cdot\mathbf A\)、\(\nabla^2f\)
  是标量；\(\nabla^2\mathbf A\) 是矢量。
\item
  \textbf{再判坐标系。}
  本专题所有分量展开均为直角坐标公式。柱坐标、球坐标中的基矢随位置变化，不能只把
  \(x,y,z\) 换成其他坐标。
\item
  \textbf{检查正则性。} 一阶运算通常要求
  \(C^1\)，交换混合偏导和使用二阶恒等式要求相应场至少属于 \(C^2\)。
\item
  \textbf{检查奇点与界面。}
  点电荷、细导线、面电荷和介质界面处，场可能不连续或发散。普通微分式应在各光滑子区域内使用，并与积分形式、边界条件或分布理论结合。
\item
  \textbf{区分局部与整体。} \(\nabla\times\mathbf E=\mathbf0\)
  是局部结论；要在整个区域内推出单值势函数，还需适当的拓扑条件，例如区域单连通。
\item
  \textbf{固定定向。} 闭合曲面的 \(\mathrm d\mathbf S\)
  取外法向；曲面法向与边界 \(\mathrm d\mathbf l\)
  的正向由右手定则共同确定。
\item
  \textbf{最后做量纲与极限检查。}
  每次空间微分多引入一个长度倒数；对称点、远场和退化情形应与物理直觉一致。
\end{enumerate}

\begin{center}
\includegraphics[width=\linewidth,height=.40\textheight,keepaspectratio]{\detokenize{assets/module-01/tikz-m1-topic2-v4-06-checklist.png}}
\par\smallskip{\small 图示：纳布拉运算依次检查对象、符号、坐标、定义域、规则和结果}
\end{center}

\subsection{例题}\label{m1:ux4f8bux9898}

\par\addvspace{.4\baselineskip}
\noindent\begin{minipage}{\linewidth}

\subsubsection{先判断输出对象，再开始求导（M1-12）}\label{m1:ux5148ux5224ux65adux8f93ux51faux5bf9ux8c61ux518dux5f00ux59cbux6c42ux5bfcm1-12}

在右手直角坐标系中给定标量场 \(f=x^2y\) 和矢量场
\(\mathbf A=-y\mathbf e_x+x\mathbf e_y\)，两者均与 \(z\) 无关。求
\(\nabla f\)、\(\nabla\cdot\mathbf A\) 和
\(\nabla\times\mathbf A\)；在计算前分别指出三种运算的输入对象和输出类型，并说明不能把这些不同类型的结果互相替代。


\end{minipage}\par

\par\noindent\begin{minipage}{\linewidth}
\textbf{解答}

\begin{center}
\includegraphics[width=\linewidth,height=.40\textheight,keepaspectratio]{\detokenize{assets/module-01/examples/M1-12.pdf}}
\par\smallskip{\small 图 M1-12：纳布拉运算的输出类型}
\end{center}
\end{minipage}\par

图 M1-12 上半部取 \(z=0\) 切片，蓝色开曲线满足 \(x^2y=C\)；在
\(P=(1,1,0)\)，绿色梯度 \((2,1,0)\) 垂直于经过 \(P\)
的等值线。下半部直接画本题的
\(\mathbf A=(-y,x,0)\)：圆周切向场逆时针环绕 \(z\)
轴，闭合面无净流出，而两个交叉导数相减给出纸外旋度 \(2\mathbf e_z\)。

取标量场 \(f=x^2y\) 和矢量场 \(\mathbf A=(-y,x,0)\)。由定义

\[
\nabla f=(2xy,x^2,0),\qquad
\nabla\cdot\mathbf A=0,\qquad
\nabla\times\mathbf A=2\mathbf e_z.
\]

同一个 \(\nabla\)
分别给出矢量、标量和矢量；连接方式而不是分量个数决定运算类型。

\par\addvspace{.4\baselineskip}
\noindent\begin{minipage}{\linewidth}

\subsubsection{用两条路径核对散度乘积法则（M1-13）}\label{m1:ux7528ux4e24ux6761ux8defux5f84ux6838ux5bf9ux6563ux5ea6ux4e58ux79efux6cd5ux5219m1-13}

设标量场 \(f(x,y,z)=x\)，矢量场
\(\mathbf A(x,y,z)=y\mathbf e_x+z\mathbf e_y+x\mathbf e_z\)。在同一右手直角坐标系中，先写出
\(f\mathbf A\) 的三个分量并直接求散度，再分别计算 \(\nabla f\) 与
\(\nabla\cdot\mathbf A\)，利用乘积法则核对结果。所有求导均在一般位置
\((x,y,z)\) 进行。


\end{minipage}\par

\par\noindent\begin{minipage}{\linewidth}
\textbf{解答}

\begin{center}
\includegraphics[width=\linewidth,height=.40\textheight,keepaspectratio]{\detokenize{assets/module-01/examples/M1-13.pdf}}
\par\smallskip{\small 图 M1-13：散度乘积法则的两条计算路线}
\end{center}
\end{minipage}\par

图 M1-13上图的 \(f=x\) 说明权重随 \(x\)
改变；下排连接的是``先相乘''的运算关系，不是场流线。对完整矢量场
\(\mathbf A=(y,z,x)\)，三个同指标偏导都为零，故
\(\nabla\cdot\mathbf A=0\)。分配求导后只剩
\(\mathbf A\cdot\mathbf e_x=y\)；这与先得到 \((xy,xz,x^2)\)
再取散度一致。

令 \(f=x\)、\(\mathbf A=(y,z,x)\)。直接先乘得
\(f\mathbf A=(xy,xz,x^2)\)，故

\[
\nabla\cdot(f\mathbf A)=y.
\]

另一方面，\(\nabla f=\mathbf e_x\)、\(\nabla\cdot\mathbf A=0\)，于是

\[
f\nabla\cdot\mathbf A+(\nabla f)\cdot\mathbf A=0+y=y.
\]

两条路径一致。乘积法则的附加项来自标量因子的空间变化；若 \(f\)
为常量，该项才消失。

\subsection{五、练习与解答}\label{m1:ux4e94ux7ec3ux4e60ux4e0eux89e3ux7b54}

\par\addvspace{.4\baselineskip}
\noindent\begin{minipage}{\linewidth}

\subsubsection{作业一：四种运算与输出类型}\label{m1:ux4f5cux4e1aux4e00ux56dbux79cdux8fd0ux7b97ux4e0eux8f93ux51faux7c7bux578b}

\textbf{题号：M1-14}

\textbf{层级：} 基础　　\textbf{方法：} 矢量分析

使用右手直角坐标系，在一般观察点计算；矢量拉普拉斯或旋度均按相应分量定义，不把不同输出类型混同。

给定

\[
f=x^2y+z^2,\qquad
\mathbf A=(x^2,yz,xz).
\]

计算 \(\nabla f\)、\(\nabla^2f\)、\(\nabla\cdot\mathbf A\) 和
\(\nabla\times\mathbf A\)，并说明每个结果是标量还是矢量。


\end{minipage}\par

\par\noindent\begin{minipage}{\linewidth}
\textbf{解答}

\begin{center}
\includegraphics[width=\linewidth,height=.40\textheight,keepaspectratio]{\detokenize{assets/module-01/examples/M1-14.pdf}}
\par\smallskip{\small 图 M1-14：四种运算与输出类型}
\end{center}
\end{minipage}\par

图 M1-14 上半部明确取 \(z=0\)，此时 \(f=x^2y+z^2\) 的等值线为
\(x^2y=C\)，在 \(P=(1,1,0)\) 的梯度是 \((2,1,0)\)。下半部同一切片上
\(\mathbf A=x^2\mathbf e_x\)，但完整三维旋度在 \(P\) 为
\(-\mathbf e_x\)，图中红箭头因此朝左；求旋度必须先对完整场求导，再代入
\(z=0\)。散度和拉普拉斯同样须保留未画出的 \(z\) 方向贡献，得到右侧的
\(3x+z\) 与 \(2y+2\)。

由定义，

\[
\nabla f
=(2xy,x^2,2z),
\qquad
\nabla^2f
=\partial_x^2f+\partial_y^2f+\partial_z^2f
=2y+2.
\]

对矢量场，

\[
\nabla\cdot\mathbf A
=\partial_x(x^2)+\partial_y(yz)+\partial_z(xz)
=2x+z+x=3x+z,
\]

\[
\begin{aligned}
\nabla\times\mathbf A
&=\bigl(\partial_y(xz)-\partial_z(yz),\,
\partial_z(x^2)-\partial_x(xz),\\
&\qquad\partial_x(yz)-\partial_y(x^2)\bigr)\\
&=(-y,-z,0).
\end{aligned}
\]

因此 \(\nabla f\) 与 \(\nabla\times\mathbf A\) 是矢量场，\(\nabla^2f\)
与 \(\nabla\cdot\mathbf A\) 是标量场。

\par\addvspace{.4\baselineskip}
\noindent\begin{minipage}{\linewidth}

\subsubsection{作业二：同时验证散度与旋度的乘积法则}\label{m1:ux4f5cux4e1aux4e8cux540cux65f6ux9a8cux8bc1ux6563ux5ea6ux4e0eux65cbux5ea6ux7684ux4e58ux79efux6cd5ux5219}

\textbf{题号：M1-15}

\textbf{层级：} 基础　　\textbf{方法：} 高斯定理、旋度与环量

在同一右手直角坐标系中求导，各分量均为坐标的光滑函数；两种方法应使用相同的观察点与运算次序。

令

\[
f=x,\qquad \mathbf A=(y,z,x).
\]

分别用直接分量计算和乘积法则计算
\(\nabla\cdot(f\mathbf A)\)、\(\nabla\times(f\mathbf A)\)。


\end{minipage}\par

\par\noindent\begin{minipage}{\linewidth}
\textbf{解答}

\begin{center}
\includegraphics[width=\linewidth,height=.40\textheight,keepaspectratio]{\detokenize{assets/module-01/examples/M1-15.pdf}}
\par\smallskip{\small 图 M1-15：散度和旋度的乘积法则}
\end{center}
\end{minipage}\par

图 M1-15中的线性切片提示
\(\nabla f=\mathbf e_x\)，所以两条计算路线的差别只在于何时把导数分配给
\(f\)。散度的附加项读取 \(\mathbf A\) 的 \(x\) 分量，等于
\(y\)；旋度的附加项读取叉乘
\(\mathbf e_x\times\mathbf A=(0,-x,z)\)。把这两个图示贡献分别与 \(f\)
乘以原场散度、旋度的部分相加，再与直接展开核对。

因为 \(f\mathbf A=(xy,xz,x^2)\)，直接计算得到

\[
\nabla\cdot(f\mathbf A)
=\partial_x(xy)+\partial_y(xz)+\partial_z(x^2)=y,
\]

\[
\begin{aligned}
\nabla\times(f\mathbf A)
&=\bigl(\partial_y(x^2)-\partial_z(xz),\,
\partial_z(xy)-\partial_x(x^2),\\
&\qquad \partial_x(xz)-\partial_y(xy)\bigr)\\
&=(-x,-2x,z-x).
\end{aligned}
\]

另一方面，

\[
\nabla f=\mathbf i,\qquad
\nabla\cdot\mathbf A=0,\qquad
\nabla\times\mathbf A=(-1,-1,-1).
\]

散度乘积法则给出

\[
f(\nabla\cdot\mathbf A)+\mathbf A\cdot\nabla f
=0+y=y.
\]

又有

\[
(\nabla f)\times\mathbf A
=\mathbf i\times(y\mathbf i+z\mathbf j+x\mathbf k)
=(0,-x,z),
\]

故旋度乘积法则给出

\[
f(\nabla\times\mathbf A)+(\nabla f)\times\mathbf A
=(-x,-x,-x)+(0,-x,z)
=(-x,-2x,z-x),
\]

与直接计算一致。

\par\addvspace{.4\baselineskip}
\noindent\begin{minipage}{\linewidth}

\subsubsection{作业三：由理想电偶极子的电势求电场}\label{m1:ux4f5cux4e1aux4e09ux7531ux7406ux60f3ux7535ux5076ux6781ux5b50ux7684ux7535ux52bfux6c42ux7535ux573a}

\textbf{题号：M1-16}

\textbf{层级：} 基础　　\textbf{方法：} 多极展开

介质为真空，\(\varepsilon_0\) 是真空电容率；\(\mathbf p\)
为不随位置变化的有向偶极矩。偶极轴线是 \(\mathbf p\)
所在直线，赤道面是过原点且垂直于 \(\mathbf p\) 的平面。

原点处有电偶极矩为常矢量 \(\mathbf p\) 的理想电偶极子。在
\(\mathbf r\ne\mathbf0\) 处，电势为

\[
\varphi(\mathbf r)
=\frac{1}{4\pi\varepsilon_0}
\frac{\mathbf p\cdot\mathbf r}{r^3},
\qquad r=|\mathbf r|.
\]

利用 \(\mathbf E=-\nabla\varphi\)
求电场，并检验偶极轴线与赤道面上的方向和大小。


\end{minipage}\par

\par\noindent\begin{minipage}{\linewidth}
\textbf{解答}

\begin{center}
\includegraphics[width=\linewidth,height=.40\textheight,keepaspectratio]{\detokenize{assets/module-01/examples/M1-16.pdf}}
\par\smallskip{\small 图 M1-16：由偶极电势求电场}
\end{center}
\end{minipage}\par

图 M1-16左图只用 \(\mathbf p,\mathbf r,\theta\)
定义一般场点，没有把未知电场预先假定为径向。求导时把
\(\mathbf p\cdot\mathbf r\) 与 \(r^{-3}\) 当作两个标量因子，分别产生沿
\(\mathbf p\) 和沿 \(\mathbf r\)
的项。右图的轴线与赤道是结果检验：轴线上两项合成
\(+2\mathbf p\)，赤道上只剩 \(-\mathbf p\)，相同距离的场强比为 \(2:1\)。

先把被求梯度的量写成两个标量因子的乘积：

\[
\frac{\mathbf p\cdot\mathbf r}{r^3}
=(\mathbf p\cdot\mathbf r)\,r^{-3}.
\]

由于 \(\mathbf p\) 与位置无关，

\[
\nabla(\mathbf p\cdot\mathbf r)=\mathbf p.
\]

又因

\[
\nabla r
=\left(\frac{x}{r},\frac{y}{r},\frac{z}{r}\right)
=\frac{\mathbf r}{r}
=\hat{\mathbf r},
\]

由链式法则可得

\[
\nabla(r^{-3})
=-3r^{-4}\nabla r
=-\frac{3\mathbf r}{r^5}.
\]

于是

\[
\begin{aligned}
\nabla\left[(\mathbf p\cdot\mathbf r)r^{-3}\right]
&=(\mathbf p\cdot\mathbf r)\nabla(r^{-3})
+r^{-3}\nabla(\mathbf p\cdot\mathbf r)\\
&=-\frac{3(\mathbf p\cdot\mathbf r)\mathbf r}{r^5}
+\frac{\mathbf p}{r^3}.
\end{aligned}
\]

因此

\[
\boxed{
\mathbf E(\mathbf r)
=\frac{1}{4\pi\varepsilon_0r^3}
\left[
3(\mathbf p\cdot\hat{\mathbf r})\hat{\mathbf r}
-\mathbf p
\right]},\qquad r>0.
\]

若以 \(\mathbf p\) 的方向为极轴，并记
\(\theta=\angle(\mathbf p,\mathbf r)\)，则

\[
\mathbf E
=\frac{p}{4\pi\varepsilon_0r^3}
\left(2\cos\theta\,\hat{\mathbf r}
+\sin\theta\,\hat{\boldsymbol\theta}\right).
\]

回到图 M1-16 的偶极轴线，\(\hat{\mathbf r}=\pm\hat{\mathbf p}\)，所以

\[
\mathbf E_{\text{轴}}
=\frac{1}{4\pi\varepsilon_0}\frac{2p}{r^3}\hat{\mathbf p}.
\]

再看图 M1-16 的赤道面，\(\mathbf p\cdot\hat{\mathbf r}=0\)，所以

\[
\mathbf E_{\text{赤道}}
=-\frac{1}{4\pi\varepsilon_0}\frac{\mathbf p}{r^3}.
\]

两处场强均按 \(1/r^3\) 衰减；在相同 \(r\)
处，轴线场强大小是赤道面场强大小的 \(2\) 倍，方向分别沿 \(\mathbf p\)
和逆 \(\mathbf p\)。一般位置的 \(\mathbf E\)
既有径向分量也有极角方向分量，不能把它预先画成沿
\(\mathbf r\)。上述公式只表示 \(r>0\)
的普通函数部分；理想点偶极子在原点的完整分布表达还含接触项，将在
\(\delta\) 函数相关专题中讨论。

\section{专题二附录一：纳布拉算符的运算性质（拔高选讲）}\label{m1:ux4e13ux9898ux4e8cux9644ux5f55ux4e00ux7eb3ux5e03ux62c9ux7b97ux7b26ux7684ux8fd0ux7b97ux6027ux8d28ux62d4ux9ad8ux9009ux8bb2}

本附录把常见高阶恒等式集中成表。应先掌握梯度、散度、旋度的物理意义与基本计算，再调用这些恒等式；完整分量证明见后续附录。

\begin{center}
\includegraphics[width=\linewidth,height=.40\textheight,keepaspectratio]{\detokenize{assets/module-01/tikz-m1m2-29-nabla-identities-map.pdf}}
\par\smallskip{\small 图示：纳布拉恒等式按线性、乘积、二阶与方向导数分类}
\end{center}

纳布拉算符 \(\nabla\)
具有双重身份：它既是一个\textbf{空间微分算子}，又在形式上表现为一个\textbf{矢量}。因此，涉及
\(\nabla\)
的运算必须同时遵循\textbf{微积分中的求导法则（如乘积法则）以及矢量代数规则}。

计算含 \(\nabla\)
的式子时要同时检查两点：点乘或叉乘决定输出类型，每个偏导数的作用对象决定乘积法则产生哪些项。

\begin{center}
\includegraphics[width=\linewidth,height=.40\textheight,keepaspectratio]{\detokenize{assets/module-01/tikz-m1-nabla-id-01-dual-role.pdf}}
\par\smallskip{\small 图示：纳布拉恒等式要同时追踪矢量结构和微分作用}
\end{center}

下列恒等式均在三维直角坐标系中使用，并假设标量场 \(u,v\) 与矢量场
\(\mathbf A,\mathbf B\) 具有所需阶次的连续偏导数。

\subsubsection{一、线性性质}\label{m1:ux4e00ux7ebfux6027ux6027ux8d28}

作为微分算子，\(\nabla\)
对标量场和矢量场的加法与数乘满足基础的线性分布律（\(c_1, c_2\)
为常数）：

\begin{enumerate}
\def\labelenumi{\arabic{enumi}.}
\item
  \textbf{梯度的线性律}：

  \[
  \nabla (c_1 u + c_2 v) = c_1 \nabla u + c_2 \nabla v
  \]
\end{enumerate}

\begin{itemize}
\item
  \textbf{散度的线性律}：

  \[
  \nabla \cdot (c_1 \mathbf{A} + c_2 \mathbf{B}) = c_1 (\nabla \cdot \mathbf{A}) + c_2 (\nabla \cdot \mathbf{B})
  \]
\item
  \textbf{旋度的线性律}：

  \[
  \nabla \times (c_1 \mathbf{A} + c_2 \mathbf{B}) = c_1 (\nabla \times \mathbf{A}) + c_2 (\nabla \times \mathbf{B})
  \]
\end{itemize}

\subsubsection{二、乘积法则}\label{m1:ux4e8cux4e58ux79efux6cd5ux5219}

当 \(\nabla\)
作用于两个场的乘积时，必须采用微积分中的\textbf{莱布尼茨法则（求导的乘积法则）}
，并将 \(\nabla\)
依次作用于每一个因子。运算时需严格保持矢量点乘与叉乘的代数形式。

使用乘积法则时，先标出 \(\nabla\)
作用的因子。下图并列梯度、散度和旋度的乘积法则，可以直接看出第二项的来源，以及点乘、叉乘次序的限制。

\begin{center}
\includegraphics[width=\linewidth,height=.40\textheight,keepaspectratio]{\detokenize{assets/module-01/tikz-m1-nabla-id-02-product-rule-splitting.pdf}}
\par\smallskip{\small 图示：乘积法则要求纳布拉依次作用到每一个因子}
\end{center}

\begin{enumerate}
\def\labelenumi{\arabic{enumi}.}
\item
  \textbf{标量与标量的乘积（梯度）} ：

  \[
  \nabla (uv) = u\nabla v + v\nabla u
  \]
\end{enumerate}

\begin{itemize}
\item
  \textbf{标量与矢量的乘积（散度与旋度）} ：

  \begin{itemize}
  \item
    散度：

    \[
    \nabla \cdot (u\mathbf{A}) = u(\nabla \cdot \mathbf{A}) + \mathbf{A} \cdot (\nabla u)
    \]
  \item
    旋度：

    \[
    \nabla \times (u\mathbf{A}) = u(\nabla \times \mathbf{A}) + (\nabla u) \times \mathbf{A}
    \]
  \end{itemize}

  \emph{(注：推导时应分别展开 \(\nabla\) 对 \(u\) 与 \(\mathbf A\)
  的偏导数，再按点乘或叉乘结构重新分组；不能把 \(\nabla\)
  当普通矢量而按``先后顺序''搬动。)}
\item
  \textbf{矢量与矢量的点乘（梯度）} ：

  \[
  \nabla (\mathbf{A} \cdot \mathbf{B}) = (\mathbf{A} \cdot \nabla)\mathbf{B} + (\mathbf{B} \cdot \nabla)\mathbf{A} + \mathbf{A} \times (\nabla \times \mathbf{B}) + \mathbf{B} \times (\nabla \times \mathbf{A})
  \]
\end{itemize}

\emph{(注：此公式在流体力学推导伯努利方程，以及电动力学推导洛伦兹力做功时极为重要。其中}
\emph{\((\mathbf{A} \cdot \nabla)\)} \emph{是一个标量微分算符。)}

\begin{itemize}
\item
  \textbf{矢量与矢量的叉乘（散度与旋度）} ：

  \begin{itemize}
  \item
    散度：

    \[
    \nabla \cdot (\mathbf{A} \times \mathbf{B}) = \mathbf{B} \cdot (\nabla \times \mathbf{A}) - \mathbf{A} \cdot (\nabla \times \mathbf{B})
    \]
  \end{itemize}

  \emph{(注：负号应由分量展开确认：展开左式后，含 \(\partial_iB_j\)
  的三项正好组合为 \(-\mathbf A\cdot(\nabla\times\mathbf B)\)。这里
  \(\nabla\) 必须作用于右侧场，不能把它像普通矢量一样移到 \(\mathbf B\)
  的另一边。)}

  \begin{itemize}
  \item
    旋度：

    \[
    \nabla \times (\mathbf{A} \times \mathbf{B}) = \mathbf{A}(\nabla \cdot \mathbf{B}) - \mathbf{B}(\nabla \cdot \mathbf{A}) + (\mathbf{B} \cdot \nabla)\mathbf{A} - (\mathbf{A} \cdot \nabla)\mathbf{B}
    \]
  \end{itemize}
\end{itemize}

\subsubsection{三、二阶微分恒等式}\label{m1:ux4e09ux4e8cux9636ux5faeux5206ux6052ux7b49ux5f0f}

二阶微分恒等式涉及 \(\nabla\)
算符连续两次作用。根据偏导数可交换的克莱罗定理（Clairaut's
theorem），其中包含两个重要的``零恒等式''。

这两个``零恒等式''是局部光滑场中的结论。竞赛题里若出现点电荷、细导线、挖孔区域、多值势函数等奇异结构，要先确认讨论点是否真的在场的光滑定义域内，再决定能否直接套用微分恒等式。

\begin{center}
\includegraphics[width=\linewidth,height=.40\textheight,keepaspectratio]{\detokenize{assets/module-01/figure-028-revised.pdf}}
\par\smallskip{\small 图示：二阶纳布拉恒等式的局部抵消与适用边界}
\end{center}

\begin{enumerate}
\def\labelenumi{\arabic{enumi}.}
\item
  \textbf{标量场的拉普拉斯（Laplacian）} ：\\
  梯度的散度构成拉普拉斯算符。

  \[
  \nabla \cdot (\nabla u) = \nabla^2 u = \Delta u
  \]
\end{enumerate}

\begin{itemize}
\item
  \textbf{梯度的旋度恒为零（无旋场定理）} ：

  \[
  \nabla \times (\nabla u) = \mathbf{0}
  \]
\end{itemize}

\emph{(物理意义：任何保守力场（如静电场、重力场）均可表示为某个标量势的梯度，因此保守场必定是无旋场。反过来，若要由``无旋''推出``可写成标量势梯度''，还需要区域单连通且场足够光滑。)}

\begin{itemize}
\item
  \textbf{旋度场的散度恒为零（局部恒等式）} ：

  \[
  \nabla \cdot (\nabla \times \mathbf{A}) = 0
  \]
\end{itemize}

\emph{(物理意义：在通常不含磁单极子的麦克斯韦理论中，物理公设是
\(\nabla\cdot\mathbf B=0\)。在无磁荷的局部单连通区域内，才可进一步引入矢量势使
\(\mathbf B=\nabla\times\mathbf A\)；后式是前一无散条件的局部表示，不是用来证明``没有磁单极子''的原因。)}

\begin{itemize}
\item
  \textbf{旋度的旋度（双重旋度）} ：

  \[
  \nabla \times (\nabla \times \mathbf{A}) = \nabla (\nabla \cdot \mathbf{A}) - \nabla^2 \mathbf{A}
  \]
\end{itemize}

\emph{(注：推导电磁波波动方程时会用到这个恒等式。)}

\subsubsection{\texorpdfstring{四、 特殊算符拓展：方向导数与
\((\mathbf{A} \cdot \nabla)\)}{四、 特殊算符拓展：方向导数与 (\textbackslash mathbf\{A\} \textbackslash cdot \textbackslash nabla)}}\label{m1:ux56db-ux7279ux6b8aux7b97ux7b26ux62d3ux5c55ux65b9ux5411ux5bfcux6570ux4e0e-mathbfa-cdot-nabla}

当算符 \(\nabla\) 与一个给定矢量 \(\mathbf{A}\)
作标量积时，会产生一个新的\textbf{标量微分算符}
\((\mathbf{A} \cdot \nabla)\)：

\[
(\mathbf{A} \cdot \nabla) = A_x \frac{\partial}{\partial x} + A_y \frac{\partial}{\partial y} + A_z \frac{\partial}{\partial z}
\]

这个点乘仍是微分算符，会继续作用于后面的对象。\(\mathbf{A}\)
给出取变化率的方向，\(\nabla\) 负责空间求导。

\begin{center}
\includegraphics[width=\linewidth,height=.40\textheight,keepaspectratio]{\detokenize{assets/module-01/tikz-m1-nabla-id-04-directional-operator.pdf}}
\par\smallskip{\small 图示：方向导数算符表示沿给定矢量方向观察场的变化}
\end{center}

\begin{enumerate}
\def\labelenumi{\arabic{enumi}.}
\item
  \textbf{物理意义}：当 \(\mathbf{A}\) 为单位矢量 \(\mathbf{n}\)
  时，\((\mathbf{n} \cdot \nabla) u = \frac{\partial u}{\partial n}\)，即为标量场
  \(u\) 沿方向 \(\mathbf{n}\) 的\textbf{方向导数}。
\item
  \textbf{作用于矢量场}：算符 \((\mathbf{A} \cdot \nabla)\)
  也可以直接作用于另一个矢量场
  \(\mathbf{B}\)，结果是一个新的矢量场。这在流体动力学中的``对流导数（Convective
  derivative）''项 \((\mathbf{v} \cdot \nabla)\mathbf{v}\)
  中有着关键应用。

  \[
  \begin{aligned}
  (\mathbf{A} \cdot \nabla)\mathbf{B}
  &=
  \left(A_x\frac{\partial B_x}{\partial x}
  +A_y\frac{\partial B_x}{\partial y}
  +A_z\frac{\partial B_x}{\partial z}\right)\mathbf i\\
  &\quad+
  \left(A_x\frac{\partial B_y}{\partial x}
  +A_y\frac{\partial B_y}{\partial y}
  +A_z\frac{\partial B_y}{\partial z}\right)\mathbf j\\
  &\quad+
  \left(A_x\frac{\partial B_z}{\partial x}
  +A_y\frac{\partial B_z}{\partial y}
  +A_z\frac{\partial B_z}{\partial z}\right)\mathbf k.
  \end{aligned}
  \]
\end{enumerate}

\section{专题二附录二：纳布拉运算性质的直角坐标分量证明（拔高选讲）}\label{m1:ux4e13ux9898ux4e8cux9644ux5f55ux4e8cux7eb3ux5e03ux62c9ux8fd0ux7b97ux6027ux8d28ux7684ux76f4ux89d2ux5750ux6807ux5206ux91cfux8bc1ux660eux62d4ux9ad8ux9009ux8bb2}

本附录给出数学物理层面的分量证明，用以检查矢量恒等式的符号、求导次序和适用条件。使用基本结论时可按需查阅相关证明。

\begin{center}
\includegraphics[width=\linewidth,height=.40\textheight,keepaspectratio]{\detokenize{assets/module-01/tikz-m1m2-30-component-proof-route.pdf}}
\par\smallskip{\small 图示：分量法证明矢量恒等式时，先证明一个分量，再利用轮换对称推广}
\end{center}

图中的``局部''二字不可省略：分量展开默认函数在考察区域内足够光滑，混合偏导可以交换。带点电荷、细导线或多值势函数的题目，恒等式可以在去掉奇点后的区域成立，却未必能直接推出跨越整个区域的全局结论。

为保证推导的严谨性，且不引入爱因斯坦求和约定，以下证明均在三维直角坐标系
\((x, y, z)\) 中进行，采用\textbf{分量展开法}。算符定义为
\(\nabla = \mathbf{i}\frac{\partial}{\partial x} + \mathbf{j}\frac{\partial}{\partial y} + \mathbf{k}\frac{\partial}{\partial z}\)。对于复杂的矢量等式，我们将重点证明其
\(x\) 分量，其余分量可由空间坐标的\textbf{轮换对称性}直接得出。

\begin{center}
\includegraphics[width=\linewidth,height=.40\textheight,keepaspectratio]{\detokenize{assets/module-01/figure-031-revised.pdf}}
\par\smallskip{\small 图示：分量法证明纳布拉恒等式时，需要同时追踪对象、分量和结构三本账}
\end{center}

\subsubsection{一、线性性质的证明}\label{m1:ux4e00ux7ebfux6027ux6027ux8d28ux7684ux8bc1ux660e}

基于微积分中偏导数算符自身的线性性质
\(\frac{\partial}{\partial x}(c_1 f + c_2 g) = c_1 \frac{\partial f}{\partial x} + c_2 \frac{\partial g}{\partial x}\)，\(\nabla\)
算符直接继承了该性质。

\textbf{1. 梯度的线性律：}
\(\nabla (c_1 u + c_2 v) = c_1 \nabla u + c_2 \nabla v\)

\textbf{证明：}

\[
\nabla (c_1 u + c_2 v) = \mathbf{i}\frac{\partial}{\partial x}(c_1 u + c_2 v) + \mathbf{j}\frac{\partial}{\partial y}(c_1 u + c_2 v) + \mathbf{k}\frac{\partial}{\partial z}(c_1 u + c_2 v)
\]

\[
= \mathbf{i}\left(c_1\frac{\partial u}{\partial x} + c_2\frac{\partial v}{\partial x}\right) + \mathbf{j}\left(c_1\frac{\partial u}{\partial y} + c_2\frac{\partial v}{\partial y}\right) + \mathbf{k}\left(c_1\frac{\partial u}{\partial z} + c_2\frac{\partial v}{\partial z}\right) = c_1 \nabla u + c_2 \nabla v
\]

散度与旋度的线性律同理可由单项偏导数的线性展开直接得证，此处不再赘述。

\subsubsection{二、乘积法则的证明}\label{m1:ux4e8cux4e58ux79efux6cd5ux5219ux7684ux8bc1ux660e}

此类性质依赖于微积分中的乘积求导法则
\(\frac{\partial}{\partial x}(fg) = \frac{\partial f}{\partial x}g + f\frac{\partial g}{\partial x}\)。

\begin{center}
\includegraphics[width=\linewidth,height=.40\textheight,keepaspectratio]{\detokenize{assets/module-01/figure-032-revised.pdf}}
\par\smallskip{\small 图示：纳布拉算符的乘积法则来自每个因子都被微分一次}
\end{center}

\textbf{1. 标量与标量的乘积（梯度）：}
\(\nabla (uv) = u\nabla v + v\nabla u\)

\textbf{证明：}

\[
\nabla (uv) = \mathbf{i}\frac{\partial (uv)}{\partial x} + \mathbf{j}\frac{\partial (uv)}{\partial y} + \mathbf{k}\frac{\partial (uv)}{\partial z}
\]

\[
= \mathbf{i}\left(u\frac{\partial v}{\partial x} + v\frac{\partial u}{\partial x}\right) + \mathbf{j}\left(u\frac{\partial v}{\partial y} + v\frac{\partial u}{\partial y}\right) + \mathbf{k}\left(u\frac{\partial v}{\partial z} + v\frac{\partial u}{\partial z}\right)
\]

将含有 \(u\) 和 \(v\) 的项分别合并提取：

\[
= u\left(\mathbf{i}\frac{\partial v}{\partial x} + \mathbf{j}\frac{\partial v}{\partial y} + \mathbf{k}\frac{\partial v}{\partial z}\right) + v\left(\mathbf{i}\frac{\partial u}{\partial x} + \mathbf{j}\frac{\partial u}{\partial y} + \mathbf{k}\frac{\partial u}{\partial z}\right) = u\nabla v + v\nabla u
\]

\textbf{2. 标量与矢量的乘积（散度）：}
\(\nabla \cdot (u\mathbf{A}) = u(\nabla \cdot \mathbf{A}) + \mathbf{A} \cdot (\nabla u)\)

\textbf{证明：}

\[
\nabla \cdot (u\mathbf{A}) = \frac{\partial (uA_x)}{\partial x} + \frac{\partial (uA_y)}{\partial y} + \frac{\partial (uA_z)}{\partial z}
\]

对每一项运用乘积法则：

\[
= \left(u\frac{\partial A_x}{\partial x} + A_x\frac{\partial u}{\partial x}\right) + \left(u\frac{\partial A_y}{\partial y} + A_y\frac{\partial u}{\partial y}\right) + \left(u\frac{\partial A_z}{\partial z} + A_z\frac{\partial u}{\partial z}\right)
\]

重新分组组合：

\[
= u\left(\frac{\partial A_x}{\partial x} + \frac{\partial A_y}{\partial y} + \frac{\partial A_z}{\partial z}\right) + \left(A_x\frac{\partial u}{\partial x} + A_y\frac{\partial u}{\partial y} + A_z\frac{\partial u}{\partial z}\right) = u(\nabla \cdot \mathbf{A}) + \mathbf{A} \cdot (\nabla u)
\]

\textbf{3. 标量与矢量的乘积（旋度）：}
\(\nabla \times (u\mathbf{A}) = u(\nabla \times \mathbf{A}) + (\nabla u) \times \mathbf{A}\)

\textbf{证明：}

考察左式 \(\nabla \times (u\mathbf{A})\) 的 \(x\) 分量：

\[
[\nabla \times (u\mathbf{A})]_x = \frac{\partial (uA_z)}{\partial y} - \frac{\partial (uA_y)}{\partial z}
\]

展开偏导数：

\[
= \left(u\frac{\partial A_z}{\partial y} + A_z\frac{\partial u}{\partial y}\right) - \left(u\frac{\partial A_y}{\partial z} + A_y\frac{\partial u}{\partial z}\right)
\]

重新分组为两部分：

\[
= u\left(\frac{\partial A_z}{\partial y} - \frac{\partial A_y}{\partial z}\right) + \left(\frac{\partial u}{\partial y}A_z - \frac{\partial u}{\partial z}A_y\right)
\]

第一部分正是 \(u(\nabla \times \mathbf{A})_x\)；第二部分恰好是矢量叉乘
\((\nabla u \times \mathbf{A})\) 的 \(x\) 分量。同理可证 \(y, z\)
分量。原式得证。

\begin{center}
\includegraphics[width=\linewidth,height=.40\textheight,keepaspectratio]{\detokenize{assets/module-01/figure-033-revised.pdf}}
\par\smallskip{\small 图示：复杂乘积恒等式要靠展开后的多余项抵消或按分量重新收集}
\end{center}

\textbf{4. 矢量与矢量的点乘（梯度）：}
\(\nabla (\mathbf{A} \cdot \mathbf{B}) = (\mathbf{A} \cdot \nabla)\mathbf{B} + (\mathbf{B} \cdot \nabla)\mathbf{A} + \mathbf{A} \times (\nabla \times \mathbf{B}) + \mathbf{B} \times (\nabla \times \mathbf{A})\)

\textbf{证明：}

考察左式 \(\nabla(\mathbf{A} \cdot \mathbf{B})\) 的 \(x\) 分量：

\[
[\nabla(\mathbf{A} \cdot \mathbf{B})]_x = \frac{\partial}{\partial x}(A_x B_x + A_y B_y + A_z B_z)
\]

\[
= A_x \frac{\partial B_x}{\partial x} + B_x \frac{\partial A_x}{\partial x} + A_y \frac{\partial B_y}{\partial x} + B_y \frac{\partial A_y}{\partial x} + A_z \frac{\partial B_z}{\partial x} + B_z \frac{\partial A_z}{\partial x}
\]

现在分别考察右式中各项的 \(x\) 分量：

\begin{itemize}
\tightlist
\item
  \[
  [(\mathbf{A} \cdot \nabla)\mathbf{B}]_x = A_x\frac{\partial B_x}{\partial x} + A_y\frac{\partial B_x}{\partial y} + A_z\frac{\partial B_x}{\partial z}
  \]
\item
  \[
  [(\mathbf{B} \cdot \nabla)\mathbf{A}]_x = B_x\frac{\partial A_x}{\partial x} + B_y\frac{\partial A_x}{\partial y} + B_z\frac{\partial A_x}{\partial z}
  \]
\item
  \[
  [\mathbf{A} \times (\nabla \times \mathbf{B})]_x = A_y(\nabla \times \mathbf{B})_z - A_z(\nabla \times \mathbf{B})_y = A_y\left(\frac{\partial B_y}{\partial x} - \frac{\partial B_x}{\partial y}\right) - A_z\left(\frac{\partial B_x}{\partial z} - \frac{\partial B_z}{\partial x}\right)
  \]
\item
  \[
  [\mathbf{B} \times (\nabla \times \mathbf{A})]_x = B_y\left(\frac{\partial A_y}{\partial x} - \frac{\partial A_x}{\partial y}\right) - B_z\left(\frac{\partial A_x}{\partial z} - \frac{\partial A_z}{\partial x}\right)
  \]
\end{itemize}

将右侧这四个 \(x\) 分量相加。我们发现交叉项会成对抵消：

\(A_y\frac{\partial B_x}{\partial y}\) 与
\(-A_y\frac{\partial B_x}{\partial y}\)
抵消；\(A_z\frac{\partial B_x}{\partial z}\) 与
\(-A_z\frac{\partial B_x}{\partial z}\) 抵消。

同理，含有 \(B\) 和 \(A\) 的偏导数交叉项也会完全抵消。

剩余项恰好为：

\[
A_x \frac{\partial B_x}{\partial x} + B_x \frac{\partial A_x}{\partial x} + A_y \frac{\partial B_y}{\partial x} + A_z \frac{\partial B_z}{\partial x} + B_y \frac{\partial A_y}{\partial x} + B_z \frac{\partial A_z}{\partial x}
\]

此结果与左式完全一致。同理可证 \(y, z\) 分量。

\textbf{5. 矢量与矢量的叉乘（散度）：}
\(\nabla \cdot (\mathbf{A} \times \mathbf{B}) = \mathbf{B} \cdot (\nabla \times \mathbf{A}) - \mathbf{A} \cdot (\nabla \times \mathbf{B})\)

\textbf{证明：}

左式展开为：

\[
\nabla \cdot (\mathbf{A} \times \mathbf{B}) = \frac{\partial}{\partial x}(A_y B_z - A_z B_y) + \frac{\partial}{\partial y}(A_z B_x - A_x B_z) + \frac{\partial}{\partial z}(A_x B_y - A_y B_x)
\]

应用乘积法则全部展开，共产生 12 项：

\[
= \frac{\partial A_y}{\partial x}B_z + A_y\frac{\partial B_z}{\partial x} - \frac{\partial A_z}{\partial x}B_y - A_z\frac{\partial B_y}{\partial x} + \frac{\partial A_z}{\partial y}B_x + A_z\frac{\partial B_x}{\partial y} - \frac{\partial A_x}{\partial y}B_z - A_x\frac{\partial B_z}{\partial y} + \frac{\partial A_x}{\partial z}B_y + A_x\frac{\partial B_y}{\partial z} - \frac{\partial A_y}{\partial z}B_x - A_y\frac{\partial B_x}{\partial z}
\]

按照 \(\mathbf{B}\) 和 \(\mathbf{A}\) 的分量进行重新同类项合并：

\[
= B_x\left(\frac{\partial A_z}{\partial y} - \frac{\partial A_y}{\partial z}\right) + B_y\left(\frac{\partial A_x}{\partial z} - \frac{\partial A_z}{\partial x}\right) + B_z\left(\frac{\partial A_y}{\partial x} - \frac{\partial A_x}{\partial y}\right) - A_x\left(\frac{\partial B_z}{\partial y} - \frac{\partial B_y}{\partial z}\right) - A_y\left(\frac{\partial B_x}{\partial z} - \frac{\partial B_z}{\partial x}\right) - A_z\left(\frac{\partial B_y}{\partial x} - \frac{\partial B_x}{\partial y}\right)
\]

观察括号内部分，前三项恰为
\(\mathbf{B} \cdot (\nabla \times \mathbf{A})\)，后三项恰为
\(-\mathbf{A} \cdot (\nabla \times \mathbf{B})\)。原式得证。

\textbf{6. 矢量与矢量的叉乘（旋度）：}
\(\nabla \times (\mathbf{A} \times \mathbf{B}) = \mathbf{A}(\nabla \cdot \mathbf{B}) - \mathbf{B}(\nabla \cdot \mathbf{A}) + (\mathbf{B} \cdot \nabla)\mathbf{A} - (\mathbf{A} \cdot \nabla)\mathbf{B}\)

\textbf{证明：}

考察左式 \(\nabla \times (\mathbf{A} \times \mathbf{B})\) 的 \(x\)
分量：

\[
[\nabla \times (\mathbf{A} \times \mathbf{B})]_x = \frac{\partial}{\partial y}(\mathbf{A} \times \mathbf{B})_z - \frac{\partial}{\partial z}(\mathbf{A} \times \mathbf{B})_y = \frac{\partial}{\partial y}(A_x B_y - A_y B_x) - \frac{\partial}{\partial z}(A_z B_x - A_x B_z)
\]

展开偏导数（共 8 项）：

\[
= \frac{\partial A_x}{\partial y}B_y + A_x\frac{\partial B_y}{\partial y} - \frac{\partial A_y}{\partial y}B_x - A_y\frac{\partial B_x}{\partial y} - \frac{\partial A_z}{\partial z}B_x - A_z\frac{\partial B_x}{\partial z} + \frac{\partial A_x}{\partial z}B_z + A_x\frac{\partial B_z}{\partial z}
\]

为了凑出散度 \(\nabla \cdot \mathbf{B}\) 和
\(\nabla \cdot \mathbf{A}\)，我们在式子中\textbf{加上并减去}项
\(A_x\frac{\partial B_x}{\partial x}\) 与
\(B_x\frac{\partial A_x}{\partial x}\)：

重新提取 \(A_x\) 和 \(-B_x\) 为公因式：

\[
= A_x\left(\frac{\partial B_x}{\partial x} + \frac{\partial B_y}{\partial y} + \frac{\partial B_z}{\partial z}\right) - B_x\left(\frac{\partial A_x}{\partial x} + \frac{\partial A_y}{\partial y} + \frac{\partial A_z}{\partial z}\right) + \left(B_x\frac{\partial A_x}{\partial x} + B_y\frac{\partial A_x}{\partial y} + B_z\frac{\partial A_x}{\partial z}\right) - \left(A_x\frac{\partial B_x}{\partial x} + A_y\frac{\partial B_x}{\partial y} + A_z\frac{\partial B_x}{\partial z}\right)
\]

上式四部分完美对应于右式的 \(x\) 分量：

\[
= A_x(\nabla \cdot \mathbf{B}) - B_x(\nabla \cdot \mathbf{A}) + [(\mathbf{B} \cdot \nabla)\mathbf{A}]_x - [(\mathbf{A} \cdot \nabla)\mathbf{B}]_x
\]

同理可证 \(y, z\) 分量。

\subsubsection{三、 二阶微分恒等式 (Second-Order Derivatives)
的证明}\label{m1:ux4e09-ux4e8cux9636ux5faeux5206ux6052ux7b49ux5f0f-second-order-derivatives-ux7684ux8bc1ux660e}

二阶性质的证明依赖于函数足够光滑时偏导数可交换的\textbf{克莱罗定理}： \[
\frac{\partial^2 u}{\partial x \partial y} = \frac{\partial^2 u}{\partial y \partial x}.
\]

\begin{center}
\includegraphics[width=\linewidth,height=.40\textheight,keepaspectratio]{\detokenize{assets/module-01/figure-034-revised.pdf}}
\par\smallskip{\small 图示：二阶纳布拉恒等式的核心是混合偏导抵消和补项重组}
\end{center}

\textbf{1. 标量场的拉普拉斯：} \(\nabla \cdot (\nabla u) = \nabla^2 u\)

\textbf{证明：}

\[
\nabla \cdot (\nabla u) = \nabla \cdot \left(\mathbf{i}\frac{\partial u}{\partial x} + \mathbf{j}\frac{\partial u}{\partial y} + \mathbf{k}\frac{\partial u}{\partial z}\right)
\]

\[
= \frac{\partial}{\partial x}\left(\frac{\partial u}{\partial x}\right) + \frac{\partial}{\partial y}\left(\frac{\partial u}{\partial y}\right) + \frac{\partial}{\partial z}\left(\frac{\partial u}{\partial z}\right) = \frac{\partial^2 u}{\partial x^2} + \frac{\partial^2 u}{\partial y^2} + \frac{\partial^2 u}{\partial z^2} = \nabla^2 u
\]

\textbf{2. 梯度的旋度恒为零：} \(\nabla \times (\nabla u) = \mathbf{0}\)

\textbf{证明：}

考察左式的 \(x\) 分量：

\[
[\nabla \times (\nabla u)]_x = \frac{\partial}{\partial y}(\nabla u)_z - \frac{\partial}{\partial z}(\nabla u)_y = \frac{\partial}{\partial y}\left(\frac{\partial u}{\partial z}\right) - \frac{\partial}{\partial z}\left(\frac{\partial u}{\partial y}\right) = \frac{\partial^2 u}{\partial y \partial z} - \frac{\partial^2 u}{\partial z \partial y}
\]

根据克莱罗定理，混合偏导数相等，故该项为 \(0\)。同理 \(y, z\) 分量均为
\(0\)。

\textbf{3. 旋度的散度恒为零：}
\(\nabla \cdot (\nabla \times \mathbf{A}) = 0\)

\textbf{证明：}

左式展开为：

\[
\nabla \cdot (\nabla \times \mathbf{A}) = \frac{\partial}{\partial x}(\nabla \times \mathbf{A})_x + \frac{\partial}{\partial y}(\nabla \times \mathbf{A})_y + \frac{\partial}{\partial z}(\nabla \times \mathbf{A})_z
\]

\[
= \frac{\partial}{\partial x}\left(\frac{\partial A_z}{\partial y} - \frac{\partial A_y}{\partial z}\right) + \frac{\partial}{\partial y}\left(\frac{\partial A_x}{\partial z} - \frac{\partial A_z}{\partial x}\right) + \frac{\partial}{\partial z}\left(\frac{\partial A_y}{\partial x} - \frac{\partial A_x}{\partial y}\right)
\]

展开后得到六个二阶偏导数项：

\[
= \frac{\partial^2 A_z}{\partial x \partial y} - \frac{\partial^2 A_y}{\partial x \partial z} + \frac{\partial^2 A_x}{\partial y \partial z} - \frac{\partial^2 A_z}{\partial y \partial x} + \frac{\partial^2 A_y}{\partial z \partial x} - \frac{\partial^2 A_x}{\partial z \partial y}
\]

重新组合可抵消的项：

\[
= \left(\frac{\partial^2 A_z}{\partial x \partial y} - \frac{\partial^2 A_z}{\partial y \partial x}\right) + \left(\frac{\partial^2 A_x}{\partial y \partial z} - \frac{\partial^2 A_x}{\partial z \partial y}\right) + \left(\frac{\partial^2 A_y}{\partial z \partial x} - \frac{\partial^2 A_y}{\partial x \partial z}\right) = 0 + 0 + 0 = 0
\]

\textbf{4. 旋度的旋度（双重旋度）}

\[
\nabla \times (\nabla \times \mathbf{A}) = \nabla (\nabla \cdot \mathbf{A}) - \nabla^2 \mathbf{A}
\]

\textbf{证明：} \emph{(注：正文第三节已经给出分量证明，这里不再重复。)}

\section{专题二附录三：三重积恒等式在纳布拉运算中的应用（拔高选讲）}\label{m1:ux4e13ux9898ux4e8cux9644ux5f55ux4e09ux4e09ux91cdux79efux6052ux7b49ux5f0fux5728ux7eb3ux5e03ux62c9ux8fd0ux7b97ux4e2dux7684ux5e94ux7528ux62d4ux9ad8ux9009ux8bb2}

本附录说明三重积的代数结构怎样帮助记忆含 \(\nabla\)
的乘积恒等式，适合作为已经掌握基本矢量微积分后的统一整理。必须先强调：\(\nabla\)
是微分算符，不是普通矢量；普通三重积只能提供结构提示，严格证明必须落实到分量与乘积法则。

\begin{center}
\includegraphics[width=\linewidth,height=.40\textheight,keepaspectratio]{\detokenize{assets/module-01/tikz-m1-triple-nabla-rigorous-v2-01-proof-route.png}}
\par\smallskip{\small 图示：含纳布拉的乘积恒等式应从分量乘积法则和指标缩并严格推出}
\end{center}

本附录的所有推导都按同一个顺序读：先逐分量使用乘积法则，明确偏导数作用于哪个场；再用
Levi-Civita
符号整理指标；最后把分量结果还原为散度、旋度或方向导数。图中的三重积结构只用于记忆结果，不替代证明。

\subsection{A1.
数学基础回顾}\label{m1:a1.-ux6570ux5b66ux57faux7840ux56deux987e}

在探讨 \(\nabla\) 算符之前，需先明确普通矢量的三重积代数性质。设
\(\vec{A}\)、\(\vec{B}\)、\(\vec{C}\) 为任意三个常矢量：

\textbf{1. 标量三重积}

标量三重积的几何意义是三个矢量构成的平行六面体体积。其最重要的代数性质是\textbf{轮换对称性}：

\[
\vec{A} \cdot (\vec{B} \times \vec{C}) = \vec{B} \cdot (\vec{C} \times \vec{A}) = \vec{C} \cdot (\vec{A} \times \vec{B})
\]

互换叉乘与点乘的位置不改变结果：

\[
\vec{A} \cdot (\vec{B} \times \vec{C}) = (\vec{A} \times \vec{B}) \cdot \vec{C}
\]

\textbf{2. 矢量三重积}

矢量三重积的结果是一个矢量。其展开式为：

\[
\vec{A} \times (\vec{B} \times \vec{C}) = \vec{B}(\vec{A} \cdot \vec{C}) - \vec{C}(\vec{A} \cdot \vec{B})
\]

\subsection{\texorpdfstring{A2. \(\nabla\)
算符与``作用对象''记账法}{A2. \textbackslash nabla 算符与``作用对象''记账法}}\label{m1:a2.-nabla-ux7b97ux7b26ux4e0eux4f5cux7528ux5bf9ux8c61ux8bb0ux8d26ux6cd5}

在直角坐标系中

\[
\nabla=\mathbf e_i\partial_i,
\]

其中重复指标按 \(i=1,2,3\) 求和。\(\nabla\) 的形式像矢量，但
\(\partial_i\)
必须作用在其右侧的函数上。两个场相乘时，合法的起点只有逐分量乘积法则，例如

\[
\partial_i(A_jB_k)=(\partial_iA_j)B_k+A_j(\partial_iB_k).
\]

有时为了标记两项分别来自``对 \(\mathbf A\) 求导''与``对 \(\mathbf B\)
求导''，可以在演算草稿中写
\(\nabla^{(A)}\)、\(\nabla^{(B)}\)。这些上标只是\textbf{记账标签}，不能脱离被作用的乘积写成算符恒等式
\(\nabla=\nabla^{(A)}+\nabla^{(B)}\)，也不能把带标签的 \(\nabla\)
直接代入普通三重积公式。

\subsection{\texorpdfstring{A3. 标量三重积与 \(\nabla\)
算符（散度恒等式）}{A3. 标量三重积与 \textbackslash nabla 算符（散度恒等式）}}\label{m1:a3.-ux6807ux91cfux4e09ux91cdux79efux4e0e-nabla-ux7b97ux7b26ux6563ux5ea6ux6052ux7b49ux5f0f}

标量三重积的轮换对称性提示了\textbf{叉乘场的散度}应具有怎样的轮换结构；严格结论由下面的分量计算给出。

\textbf{目标公式：} 展开 \(\nabla \cdot (\vec{A} \times \vec{B})\)

下图把乘积法则产生的两类项分开表示：一类含 \(\partial_iA_j\)，另一类含
\(\partial_iB_k\)。图示用于辨认结构，下面的指标计算给出严格证明。

\begin{center}
\includegraphics[width=\linewidth,height=.40\textheight,keepaspectratio]{\detokenize{assets/module-01/tikz-m1-triple-nabla-rigorous-v2-02-divergence-cross.png}}
\par\smallskip{\small 图示：叉乘场的散度由乘积法则产生两项，负号来自Levi-Civita符号的反对称性}
\end{center}

\textbf{严格推导：} 设
\((\mathbf A\times\mathbf B)_i=\varepsilon_{ijk}A_jB_k\)。逐分量使用乘积法则：

\[
\begin{aligned}
\nabla\cdot(\mathbf A\times\mathbf B)
&=\partial_i(\varepsilon_{ijk}A_jB_k)\\
&=\varepsilon_{ijk}(\partial_iA_j)B_k
  +\varepsilon_{ijk}A_j(\partial_iB_k).
\end{aligned}
\]

第一项中作指标的循环改名，第二项中交换两个指标并利用
\(\varepsilon_{ijk}\) 的反对称性，得到

\[
\varepsilon_{ijk}(\partial_iA_j)B_k
=B_k\varepsilon_{kij}\partial_iA_j
=\mathbf B\cdot(\nabla\times\mathbf A),
\]

\[
\varepsilon_{ijk}A_j(\partial_iB_k)
=-A_j\varepsilon_{jik}\partial_iB_k
=-\mathbf A\cdot(\nabla\times\mathbf B).
\]

\textbf{结论：}

\[
\nabla \cdot (\mathbf A \times \mathbf B) = \mathbf B \cdot (\nabla \times \mathbf A) - \mathbf A \cdot (\nabla \times \mathbf B)
\]

负号最终来自 Levi-Civita
符号交换指标时的反对称性；标量三重积的轮换性质可以帮助记住这一结果，但不是把
\(\nabla\) 当普通矢量的依据。

\subsection{\texorpdfstring{A4. 矢量三重积与 \(\nabla\)
算符（旋度与梯度恒等式）}{A4. 矢量三重积与 \textbackslash nabla 算符（旋度与梯度恒等式）}}\label{m1:a4.-ux77e2ux91cfux4e09ux91cdux79efux4e0e-nabla-ux7b97ux7b26ux65cbux5ea6ux4e0eux68afux5ea6ux6052ux7b49ux5f0f}

矢量三重积的展开形式可以帮助记忆\textbf{叉乘场的旋度}以及\textbf{点乘场的梯度}，严格推导仍从分量出发。

\subsubsection{1.
叉乘场的旋度}\label{m1:ux53c9ux4e58ux573aux7684ux65cbux5ea6}

\textbf{目标公式：} 展开 \(\nabla \times (\vec{A} \times \vec{B})\)

这一步最容易漏掉方向导数项。下图把``对 \(\mathbf A\) 求导''与``对
\(\mathbf B\) 求导''产生的项分别归类，合起来正好是四项。

\begin{center}
\includegraphics[width=\linewidth,height=.40\textheight,keepaspectratio]{\detokenize{assets/module-01/tikz-m1-triple-nabla-rigorous-v2-03-curl-cross.png}}
\par\smallskip{\small 图示：叉乘场的旋度通过Levi-Civita缩并和乘积法则展开为四项}
\end{center}

\textbf{严格推导：} 利用

\[
\varepsilon_{ijk}\varepsilon_{klm}
=\delta_{il}\delta_{jm}-\delta_{im}\delta_{jl},
\]

有

\[
\begin{aligned}
[\nabla\times(\mathbf A\times\mathbf B)]_i
&=\varepsilon_{ijk}\partial_j(\varepsilon_{klm}A_lB_m)\\
&=(\delta_{il}\delta_{jm}-\delta_{im}\delta_{jl})
  \partial_j(A_lB_m)\\
&=\partial_j(A_iB_j)-\partial_j(A_jB_i)\\
&=B_j\partial_jA_i+A_i\partial_jB_j
  -B_i\partial_jA_j-A_j\partial_jB_i.
\end{aligned}
\]

四项依次对应
\((\mathbf B\cdot\nabla)\mathbf A\)、\(\mathbf A(\nabla\cdot\mathbf B)\)、\(-\mathbf B(\nabla\cdot\mathbf A)\)
与 \(-(\mathbf A\cdot\nabla)\mathbf B\)。

\textbf{结论：}

\[
\nabla \times (\mathbf A \times \mathbf B) = (\mathbf B \cdot \nabla)\mathbf A - \mathbf B(\nabla \cdot \mathbf A) + \mathbf A(\nabla \cdot \mathbf B) - (\mathbf A \cdot \nabla)\mathbf B
\]

\subsubsection{2.
点乘场的梯度}\label{m1:ux70b9ux4e58ux573aux7684ux68afux5ea6}

\textbf{目标公式：} 展开 \(\nabla(\vec{A} \cdot \vec{B})\)

点乘场的梯度也可以按``先看 \(\vec A\) 的变化、再看 \(\vec B\)
的变化''拆成两半。每一半各给出一个旋度项和一个方向导数项。

\begin{center}
\includegraphics[width=\linewidth,height=.40\textheight,keepaspectratio]{\detokenize{assets/module-01/tikz-m1-triple-nabla-rigorous-v2-04-gradient-dot.png}}
\par\smallskip{\small 图示：点乘场梯度的第i个分量由乘积法则分成两个严格可核对的部分}
\end{center}

\begin{center}
\includegraphics[width=\linewidth,height=.40\textheight,keepaspectratio]{\detokenize{assets/module-01/figure-039-revised.pdf}}
\par\smallskip{\small 图示：算例通过分别计算两侧，核对散度叉乘恒等式中的符号与作用对象}
\end{center}

这类恒等式的算例不宜直接背成``叉乘换序就变号''：先把两个给定场的旋度、叉乘和散度分开算，再让恒等式两侧在同一个结果上会合，才能检查每个算符究竟作用于哪个对象。

\textbf{严格推导：} 左端第 \(i\) 个分量为

\[
[\nabla(\mathbf A\cdot\mathbf B)]_i
=\partial_i(A_jB_j)
=B_j\partial_iA_j+A_j\partial_iB_j.
\]

右端涉及 \(\mathbf A\) 的两项满足

\[
\begin{aligned}
[\mathbf B\times(\nabla\times\mathbf A)]_i
+[(\mathbf B\cdot\nabla)\mathbf A]_i
&=\varepsilon_{ijk}B_j\varepsilon_{klm}\partial_lA_m
  +B_j\partial_jA_i\\
&=B_j(\partial_iA_j-\partial_jA_i)+B_j\partial_jA_i\\
&=B_j\partial_iA_j.
\end{aligned}
\]

交换 \(\mathbf A\)、\(\mathbf B\) 可得另外两项之和为
\(A_j\partial_iB_j\)。两部分相加，于是

\[
\nabla(\mathbf A \cdot \mathbf B) = \mathbf A \times (\nabla \times \mathbf B) + \mathbf B \times (\nabla \times \mathbf A) + (\mathbf A \cdot \nabla)\mathbf B + (\mathbf B \cdot \nabla)\mathbf A
\]

\subsection{A5. 总结}\label{m1:a5.-ux603bux7ed3}

含 \(\nabla\)
的乘积恒等式同时受两套规则约束：矢量代数决定指标怎样收缩，微分乘积法则决定每个偏导数作用于哪个因子。普通三重积适合帮助记忆结果的项数与符号；遇到疑问时，应回到
\(\partial_i\)、\(\varepsilon_{ijk}\) 和 \(\delta_{ij}\)
的分量计算。尤其不能把 \(\nabla\) 当成可任意换位的常矢量。

\section{专题三：梯度（基础必讲）}\label{m1:ux4e13ux9898ux4e09ux68afux5ea6ux57faux7840ux5fc5ux8bb2}

\begin{quote}
阅读准备：全微分是局部的一阶近似 可微标量场在一点附近满足
\(\Delta f\approx f_x\Delta x+f_y\Delta y+f_z\Delta z=\nabla f\cdot\Delta\mathbf r\)。梯度把三个坐标方向的变化率合成为一个几何矢量；舍去的是位移的二阶及更高阶小量。
\end{quote}

\par\addvspace{.4\baselineskip}
\noindent\begin{minipage}{\linewidth}

\textbf{例 M1-17：小位移中的局部线性化}

在 \(xy\) 平面上给定标量场 \(f(x,y)=x^2+y^2\)。观察点最初位于
\(P=(1,2)\)，随后发生小位移
\(d\mathbf r=0.01\mathbf e_x-0.01\mathbf e_y\)。求 \(P\)
处的梯度，并用一阶局部线性化估计函数增量
\(\Delta f\)；说明所用估计忽略了位移的哪些阶次，以及增量的正负对应怎样的变化。


\end{minipage}\par

\par\noindent\begin{minipage}{\linewidth}
\textbf{解答}

\begin{center}
\includegraphics[width=\linewidth,height=.40\textheight,keepaspectratio]{\detokenize{assets/module-01/examples/M1-17.pdf}}
\par\smallskip{\small 图 M1-17：小位移中的局部线性化}
\end{center}
\end{minipage}\par

图 M1-17中圆形等值线的法向就是 \(\nabla f(P)=(2,4)\)。橙色小位移沿
\((1,-1)\)，含有与外法向相反的分量，所以点乘为负。将放大示意的箭头换回实际位移
\((.01,-.01)\)，一阶增量为 \(-.02\)；箭头的放大不改变点乘的符号。

\(f=x^2+y^2\) 在 \((1,2)\) 处有 \(\nabla f=(2,4)\)。沿 \((0.01,-0.01)\)
移动时，\(\Delta f\approx2(0.01)+4(-0.01)=-0.02\)。

梯度把``标量在空间中怎样变化''压缩成一个矢量。它同时回答三个问题：沿给定方向每单位长度变化多少、哪个方向增长最快、等值面在该点的法向指向哪里。电势与电场、势能与保守力之间的联系，正是梯度最重要的物理应用。

\subsection{一、可微性、局部线性化与梯度}\label{m1:ux4e00ux53efux5faeux6027ux5c40ux90e8ux7ebfux6027ux5316ux4e0eux68afux5ea6}

设标量场

\[
\phi:\Omega\subset\mathbb R^3\longrightarrow\mathbb R,
\]

点 \(P\in\Omega\) 是内点。若存在一个关于位移 \(\Delta\mathbf r\)
的线性函数 \(L\)，使

\[
\phi(P+\Delta\mathbf r)-\phi(P)
=L(\Delta\mathbf r)+o(|\Delta\mathbf r|),
\qquad \Delta\mathbf r\to\mathbf0,
\]

就称 \(\phi\) 在 \(P\) 可微。欧氏空间的点积给出唯一矢量
\(\nabla\phi(P)\)，使

\[
L(\Delta\mathbf r)=\nabla\phi(P)\cdot\Delta\mathbf r.
\]

因此梯度最本质的定义是

\[
\boxed{
\phi(P+\Delta\mathbf r)
=\phi(P)+\nabla\phi(P)\cdot\Delta\mathbf r
+o(|\Delta\mathbf r|)
}.
\]

在固定的三维直角坐标系中，如果 \(\phi\) 在 \(P\)
的邻域内具有连续一阶偏导数，即 \(\phi\in C^1\)，则

\[
\boxed{
\nabla\phi
=\frac{\partial\phi}{\partial x}\mathbf e_x
+\frac{\partial\phi}{\partial y}\mathbf e_y
+\frac{\partial\phi}{\partial z}\mathbf e_z
}.
\]

\(C^1\)
是保证可微的常用充分条件，但不是可微的必要条件。梯度是几何矢量，分量公式则依赖所选坐标系与基矢；不能把直角坐标分量公式原样搬到柱坐标或球坐标中。

全微分是局部线性项的记号：

\[
d\phi=\nabla\phi\cdot d\mathbf r.
\]

梯度的量纲比原标量场多一个长度倒数：

\[
[\nabla\phi]=\frac{[\phi]}{\mathrm L}.
\]

例如电势的单位为伏特，电势梯度的单位为伏特每米，与电场单位一致。

\begin{center}
\includegraphics[width=\linewidth,height=.40\textheight,keepaspectratio]{\detokenize{assets/module-01/tikz-m1-topic3-v3-01-local-linearization.jpg}}
\par\smallskip{\small 图示：梯度是标量场在一点处局部线性变化的矢量表示}
\end{center}

\par\addvspace{.4\baselineskip}
\noindent\begin{minipage}{\linewidth}

\textbf{例 M1-18：线性标量场的常梯度}

在直角坐标系中给定线性标量场 \(\phi(x,y,z)=ax+by+cz+d\)，其中
\(a,b,c,d\) 均为常量。求任意位置的梯度，指出它是否随观察点变化；解释
\(d\) 改变时标量场的数值和局部变化规律分别如何改变。


\end{minipage}\par

\par\noindent\begin{minipage}{\linewidth}
\textbf{解答}

\begin{center}
\includegraphics[width=\linewidth,height=.40\textheight,keepaspectratio]{\detokenize{assets/module-01/examples/M1-18.pdf}}
\par\smallskip{\small 图 M1-18：线性标量场的常梯度}
\end{center}
\end{minipage}\par

图
M1-18中等值线平行、等距，三个位置的梯度箭头也平行且等长，说明局部线性变化规律不随位置改变。沿各坐标方向读取斜率分别是
\(a,b,c\)，故梯度为 \((a,b,c)\)。常数 \(d\)
只改等值线的数值标号，不改箭头方向或长度。

设 \(\phi=ax+by+cz+d\)，则

\[
\nabla\phi=a\mathbf e_x+b\mathbf e_y+c\mathbf e_z.
\]

这是常矢量，说明线性标量场在各处的局部变化规律相同；常数 \(d\)
只整体平移标量值，不改变梯度。

\subsection{二、方向导数与最陡方向}\label{m1:ux4e8cux65b9ux5411ux5bfcux6570ux4e0eux6700ux9661ux65b9ux5411}

\subsubsection{1.
沿曲线的变化率}\label{m1:ux6cbfux66f2ux7ebfux7684ux53d8ux5316ux7387}

设空间曲线写成 \(\mathbf r=\mathbf r(t)\)。复合函数链式法则给出

\[
\frac{d}{dt}\phi(\mathbf r(t))
=\nabla\phi\cdot\frac{d\mathbf r}{dt}.
\]

如果曲线以弧长 \(s\) 为参数，则

\[
\mathbf t=\frac{d\mathbf r}{ds},
\qquad |\mathbf t|=1,
\]

故标量场沿曲线每单位路程的变化率为

\[
\boxed{
\frac{d\phi}{ds}=\nabla\phi\cdot\mathbf t
}.
\]

\subsubsection{2.
给定单位方向的方向导数}\label{m1:ux7ed9ux5b9aux5355ux4f4dux65b9ux5411ux7684ux65b9ux5411ux5bfcux6570}

取单位矢量 \(\mathbf l\)，沿 \(\mathbf l\) 的方向导数定义为

\[
D_{\mathbf l}\phi(P)
=\lim_{h\to0}
\frac{\phi(P+h\mathbf l)-\phi(P)}{h}.
\]

由可微性，令 \(\Delta\mathbf r=h\mathbf l\)，有

\[
\phi(P+h\mathbf l)-\phi(P)
=h\nabla\phi(P)\cdot\mathbf l+o(|h|).
\]

两边除以 \(h\) 并取极限，得到

\[
\boxed{
D_{\mathbf l}\phi(P)=\nabla\phi(P)\cdot\mathbf l
=|\nabla\phi(P)|\cos\theta
}.
\]

其中 \(\theta\) 是 \(\nabla\phi(P)\) 与 \(\mathbf l\) 的夹角。

若给的是非单位矢量 \(\mathbf v\)，必须先分清题意：

\[
\left.\frac{d}{dt}\phi(P+t\mathbf v)\right|_{t=0}
=\nabla\phi(P)\cdot\mathbf v
\]

是关于参数 \(t\) 的变化率；沿 \(\mathbf v\)
方向每单位长度的方向导数则应使用 \(\mathbf l=\mathbf v/|\mathbf v|\)。

\begin{center}
\includegraphics[width=\linewidth,height=.40\textheight,keepaspectratio]{\detokenize{assets/module-01/tikz-m1-topic3-v3-02-directional-projection.jpg}}
\par\smallskip{\small 图示：方向导数是梯度在指定单位方向上的有向投影}
\end{center}

\subsubsection{3. 最大增长率}\label{m1:ux6700ux5927ux589eux957fux7387}

由柯西不等式，

\[
-|\nabla\phi|
\le \nabla\phi\cdot\mathbf l
\le |\nabla\phi|,
\qquad |\mathbf l|=1.
\]

当 \(\nabla\phi(P)\ne\mathbf0\) 时：

\[
\begin{array}{c|c|c}
\text{方向} & \mathbf l & D_{\mathbf l}\phi(P)\\ \hline
\text{增长最快} & \displaystyle \frac{\nabla\phi}{|\nabla\phi|} & |\nabla\phi|\\[6pt]
\text{下降最快} & \displaystyle -\frac{\nabla\phi}{|\nabla\phi|} & -|\nabla\phi|\\[6pt]
\text{一阶不变} & \mathbf l\perp\nabla\phi & 0
\end{array}
\]

这里``最大''指带符号的方向导数最大。若问 \(|D_{\mathbf l}\phi|\)
的最大值，则 \(\pm\nabla\phi\) 两个方向都达到 \(|\nabla\phi|\)。

当 \(\nabla\phi(P)=\mathbf0\)
时，所有方向的一阶方向导数都为零；此时不能说梯度给出了唯一的最陡方向，必须考察二阶或更高阶项。

\begin{center}
\includegraphics[width=\linewidth,height=.40\textheight,keepaspectratio]{\detokenize{assets/module-01/tikz-m1-topic3-v3-03-steepest-directions.jpg}}
\par\smallskip{\small 图示：梯度非零时增长最快、下降最快和一阶不变的方向彼此对应}
\end{center}

\par\addvspace{.4\baselineskip}
\noindent\begin{minipage}{\linewidth}

\textbf{例 M1-19：指定方向的变化率}

给定标量场 \(\phi=x^2y+3z\)，观察点为 \(P(1,-1,0)\)。指定有向单位矢量
\(\mathbf l=(2\mathbf e_x-\mathbf e_y+2\mathbf e_z)/3\)，方向导数按沿
\(\mathbf l\) 每移动单位长度的变化率定义。求 \(P\)
处沿该方向的方向导数，并与同一点所有单位方向中的最大增长率比较。


\end{minipage}\par

\par\noindent\begin{minipage}{\linewidth}
\textbf{解答}

\begin{center}
\includegraphics[width=\linewidth,height=.40\textheight,keepaspectratio]{\detokenize{assets/module-01/examples/M1-19.pdf}}
\par\smallskip{\small 图 M1-19：指定方向的变化率}
\end{center}
\end{minipage}\par

图 M1-19中 \(\mathbf l\)
几乎垂直于梯度，但仍有很小的正投影；这提示变化率应为正且远小于梯度模长。用给定分量计算该投影得
\(1/3\)，而取沿梯度的单位方向才达到 \(\sqrt{14}\)。这里 \(\mathbf l\)
已归一化，不应再把结果除以一次长度。

设 \(\phi=x^2y+3z\)，在 \(P(1,-1,0)\) 处沿
\(\displaystyle \mathbf l=\frac13(2,-1,2)\) 的方向求变化率。由

\[
\nabla\phi=(2xy,x^2,3),
\qquad
\nabla\phi(P)=(-2,1,3),
\]

得

\[
D_{\mathbf l}\phi(P)
=(-2,1,3)\cdot\frac13(2,-1,2)
=\frac13.
\]

该点的最大增长率为 \(\sqrt{14}\)，显然大于给定方向上的 \(1/3\)。

\subsection{三、等值面、法向与等值面疏密}\label{m1:ux4e09ux7b49ux503cux9762ux6cd5ux5411ux4e0eux7b49ux503cux9762ux758fux5bc6}

\subsubsection{1.
正则等值面的法向}\label{m1:ux6b63ux5219ux7b49ux503cux9762ux7684ux6cd5ux5411}

等值面由

\[
\phi(\mathbf r)=C
\]

定义。设 \(\mathbf r(s)\) 是等值面上经过 \(P\) 的任意光滑曲线，则

\[
\phi(\mathbf r(s))=C.
\]

对 \(s\) 求导：

\[
0=\frac{d\phi}{ds}
=\nabla\phi(P)\cdot\frac{d\mathbf r}{ds}.
\]

\(d\mathbf r/ds\) 可以取为等值面在 \(P\) 处的任意切向量。因此，若
\(\nabla\phi(P)\ne\mathbf0\)，则

\[
\boxed{
\nabla\phi(P)\perp\text{等值面在 }P\text{ 处的切平面}
}.
\]

梯度指向 \(\phi\) 增大的一侧。等值面在 \(P\) 处的单位法向可取

\[
\hat{\mathbf n}=\pm\frac{\nabla\phi}{|\nabla\phi|};
\]

正负号由题目规定的定向决定，而不是由``法向''二字自动决定。

\begin{center}
\includegraphics[width=\linewidth,height=.40\textheight,keepaspectratio]{\detokenize{assets/module-01/tikz-m1-topic3-v3-04-level-surface-normal.jpg}}
\par\smallskip{\small 图示：正则等值面处梯度垂直切平面并指向标量增大的一侧}
\end{center}

\subsubsection{2.
等值面疏密与梯度大小}\label{m1:ux7b49ux503cux9762ux758fux5bc6ux4e0eux68afux5ea6ux5927ux5c0f}

沿单位法向 \(\hat{\mathbf n}=\nabla\phi/|\nabla\phi|\) 移动小距离
\(\Delta n\)，一阶近似给出

\[
\Delta\phi
=\nabla\phi\cdot(\hat{\mathbf n}\Delta n)
+o(\Delta n)
=|\nabla\phi|\Delta n+o(\Delta n).
\]

若相邻等值面的标量差固定为 \(\Delta\phi\)，则

\[
\boxed{
\Delta n\approx\frac{|\Delta\phi|}{|\nabla\phi|}
}.
\]

所以等值面越密，\(|\nabla\phi|\) 越大；等值面越疏，\(|\nabla\phi|\)
越小。这个结论只在局部、相邻等值面足够接近且梯度非零时使用。

\subsubsection{3.
临界点不能套用``梯度给法向''}\label{m1:ux4e34ux754cux70b9ux4e0dux80fdux5957ux7528ux68afux5ea6ux7ed9ux6cd5ux5411}

满足 \(\nabla\phi(P)=\mathbf0\)
的点称为临界点。此时等值集可能仍然光滑，也可能退化、相交或缩成一点，不能仅凭一阶信息判断。

\begin{itemize}
\tightlist
\item
  对 \(\phi=x^2+y^2\)，等值集 \(\phi=0\) 在原点缩成一个点。
\item
  对 \(\phi=x^2-y^2\)，等值集 \(\phi=0\) 是两条相交直线
  \(y=\pm x\)，原点没有唯一切线和唯一法向。
\item
  对 \(\phi=x^2\)，三维零等值集 \(x=0\) 本身仍是光滑平面，但其上处处有
  \(\nabla\phi=\mathbf0\)；这时必须由平面的几何方程另求法向，不能从该梯度取得法向。
\end{itemize}

\begin{center}
\includegraphics[width=\linewidth,height=.40\textheight,keepaspectratio]{\detokenize{assets/module-01/tikz-m1-topic3-v3-05-spacing-critical-points.jpg}}
\par\smallskip{\small 图示：相同等值差下曲线越密梯度越大，临界点处可能没有唯一法向}
\end{center}

\par\addvspace{.4\baselineskip}
\noindent\begin{minipage}{\linewidth}

\textbf{例 M1-20：球面梯度与切平面}

三维空间中给定标量场 \(\phi=x^2+y^2+z^2\)。取常数 \(C>0\)，并在等值面
\(\phi=C\) 上任选一点 \(P\)，其位置矢量为
\(\mathbf r_P\)。说明该等值面的几何形状，求 \(P\)
处梯度，并写出该点切平面的法向关系，解释梯度为什么垂直于等值面。


\end{minipage}\par

\par\noindent\begin{minipage}{\linewidth}
\textbf{解答}

\begin{center}
\includegraphics[width=\linewidth,height=.40\textheight,keepaspectratio]{\detokenize{assets/module-01/examples/M1-20.pdf}}
\par\smallskip{\small 图 M1-20：球面梯度与切平面}
\end{center}
\end{minipage}\par

图 M1-20中的球面是 \(\phi=C>0\) 的等值面，半径指向的梯度为
\(2\mathbf r\)。法截面里的切线代表空间切平面，故以 \(2\mathbf r_0\)
为法向写出切平面方程。原点处梯度虽然为零，却不在这张正半径球面上，不能把它当作球面法向退化的点。

对

\[
\phi=x^2+y^2+z^2,
\]

等值面 \(\phi=C>0\) 是半径 \(\sqrt C\) 的球面，且

\[
\nabla\phi=2x\mathbf e_x+2y\mathbf e_y+2z\mathbf e_z=2\mathbf r.
\]

在球面上一点 \(P(x_0,y_0,z_0)\) 处，切平面为

\[
x_0(x-x_0)+y_0(y-y_0)+z_0(z-z_0)=0,
\]

其法向与 \(\nabla\phi(P)=2(x_0,y_0,z_0)\) 平行。原点是临界点，不属于任何
\(C>0\) 的球面。

\subsection{四、正交曲线坐标中的梯度}\label{m1:ux56dbux6b63ux4ea4ux66f2ux7ebfux5750ux6807ux4e2dux7684ux68afux5ea6}

设 \((q_1,q_2,q_3)\) 是正交曲线坐标，单位基矢为
\(\mathbf e_1,\mathbf e_2,\mathbf e_3\)，尺度因子定义为

\[
h_i=\left|\frac{\partial\mathbf r}{\partial q_i}\right|.
\]

坐标微小变化对应的真实位移为

\[
d\mathbf r
=h_1dq_1\mathbf e_1
+h_2dq_2\mathbf e_2
+h_3dq_3\mathbf e_3.
\]

另一方面，

\[
df
=\frac{\partial f}{\partial q_1}dq_1
+\frac{\partial f}{\partial q_2}dq_2
+\frac{\partial f}{\partial q_3}dq_3.
\]

令 \(\nabla f=G_1\mathbf e_1+G_2\mathbf e_2+G_3\mathbf e_3\)，代入
\(df=\nabla f\cdot d\mathbf r\) 并比较三个独立微分的系数，得到

\[
\boxed{
\nabla f
=\sum_{i=1}^{3}\frac1{h_i}
\frac{\partial f}{\partial q_i}\mathbf e_i
}.
\]

\begin{center}
\includegraphics[width=\linewidth,height=.40\textheight,keepaspectratio]{\detokenize{assets/module-01/figure-045-revised.pdf}}
\par\smallskip{\small 图示：正交坐标中的尺度因子把坐标增量换成真实长度后才能得到梯度分量}
\end{center}

三种常用坐标中的公式为：

\textbf{直角坐标} \((x,y,z)\)：

\[
\nabla f
=\frac{\partial f}{\partial x}\mathbf e_x
+\frac{\partial f}{\partial y}\mathbf e_y
+\frac{\partial f}{\partial z}\mathbf e_z.
\]

\textbf{柱坐标} \((\rho,\phi,z)\)：

\[
\boxed{
\nabla f
=\frac{\partial f}{\partial\rho}\mathbf e_{\rho}
+\frac1\rho\frac{\partial f}{\partial\phi}\mathbf e_{\phi}
+\frac{\partial f}{\partial z}\mathbf e_z
},
\qquad \rho>0.
\]

\textbf{球坐标} \((r,\theta,\phi)\)：

\[
\boxed{
\nabla f
=\frac{\partial f}{\partial r}\mathbf e_r
+\frac1r\frac{\partial f}{\partial\theta}\mathbf e_{\theta}
+\frac1{r\sin\theta}\frac{\partial f}{\partial\phi}\mathbf e_{\phi}
},
\]

其中该坐标表达式在 \(r=0\) 以及极轴 \(\sin\theta=0\)
处有坐标奇异性。坐标公式失效不一定表示物理场奇异，应改用直角坐标或取适当极限检查。

\par\addvspace{.4\baselineskip}
\noindent\begin{minipage}{\linewidth}

\textbf{例 M1-21：对数柱坐标场的梯度}

采用圆柱坐标 \((\rho,\phi,z)\)，其中 \(\rho\) 是到 \(z\)
轴的距离，\(\mathbf e_{\rho},\mathbf e_{\phi},\mathbf e_z\)
为局部单位基矢。在 \(\rho>0\) 区域给定
\(f(\rho,z)=A\ln(\rho/\rho_0)+Bz\)，\(A,B\) 为常量，\(\rho_0>0\)
为参考长度。求梯度的三个柱坐标分量，并说明结果中 \(1/\rho\)
因子的具体来源。


\end{minipage}\par

\par\noindent\begin{minipage}{\linewidth}
\textbf{解答}

\begin{center}
\includegraphics[width=\linewidth,height=.40\textheight,keepaspectratio]{\detokenize{assets/module-01/examples/M1-21.pdf}}
\par\smallskip{\small 图 M1-21：对数柱坐标场的梯度}
\end{center}
\end{minipage}\par

图 M1-21的子午截面把位移分成径向 \(\mathbf e_{\rho}\) 与轴向
\(\mathbf e_z\)；方位角不改变函数值，因此没有 \(\mathbf e_{\phi}\)
分量。分别读取对数项与线性项的导数，得到 \(A/\rho\) 和
\(B\)。图示曲线仅选正参数表达等值关系，通式中的两个系数仍按题目实际符号取值。

若

\[
f(\rho,z)=A\ln\frac{\rho}{\rho_0}+Bz,
\qquad \rho>0,
\]

则

\[
\nabla f=\frac A\rho\mathbf e_{\rho}+B\mathbf e_z.
\]

这里 \(1/\rho\) 来自对数求导，而不是柱坐标梯度公式额外产生；由于 \(f\)
与方位角无关，\(\mathbf e_{\phi}\) 分量为零。

\subsection{五、复合函数、径向场与源点梯度}\label{m1:ux4e94ux590dux5408ux51fdux6570ux5f84ux5411ux573aux4e0eux6e90ux70b9ux68afux5ea6}

\subsubsection{1.
标量复合函数的链式法则}\label{m1:ux6807ux91cfux590dux5408ux51fdux6570ux7684ux94feux5f0fux6cd5ux5219}

若 \(u=g(\mathbf r)\)，\(\phi=F(u)\)，则

\[
d\phi=F'(u)\,du
=F'(g)\nabla g\cdot d\mathbf r,
\]

所以

\[
\boxed{
\nabla F(g)=F'(g)\nabla g
}.
\]

\subsubsection{2.
球对称标量场}\label{m1:ux7403ux5bf9ux79f0ux6807ux91cfux573a}

令

\[
r=|\mathbf r|=\sqrt{x^2+y^2+z^2}.
\]

在 \(r>0\) 处，

\[
\nabla r
=\frac{x}{r}\mathbf e_x
+\frac{y}{r}\mathbf e_y
+\frac{z}{r}\mathbf e_z
=\hat{\mathbf r},
\]

故

\[
\boxed{
\nabla f(r)=f'(r)\hat{\mathbf r}
},
\qquad r>0.
\]

因此方向由 \(f'(r)\) 的符号决定：\(f'(r)>0\) 时梯度沿外径向，\(f'(r)<0\)
时梯度沿内径向。不能看到``球对称''就默认梯度向外。

\subsubsection{3.
相对于任意源点的距离}\label{m1:ux76f8ux5bf9ux4e8eux4efbux610fux6e90ux70b9ux7684ux8dddux79bb}

若源点为 \(\mathbf r'\)，场点为 \(\mathbf r\)，令

\[
R=|\mathbf r-\mathbf r'|,
\qquad
\hat{\mathbf R}=\frac{\mathbf r-\mathbf r'}{R}.
\]

对场点坐标求梯度时，\(\mathbf r'\) 固定，于是

\[
\nabla_{\mathbf r}R=\hat{\mathbf R},
\qquad
\boxed{
\nabla_{\mathbf r}\frac1R
=-\frac{\mathbf r-\mathbf r'}{R^3}
},
\qquad R>0.
\]

若改为对源点坐标求梯度，符号相反：

\[
\nabla_{\mathbf r'}\frac1R
=+\frac{\mathbf r-\mathbf r'}{R^3}.
\]

这个符号差异在连续电荷分布、多极展开和格林函数推导中十分重要。

\begin{center}
\includegraphics[width=\linewidth,height=.40\textheight,keepaspectratio]{\detokenize{assets/module-01/tikz-m1-topic3-v3-07-radial-sign-source.jpg}}
\par\smallskip{\small 图示：径向函数的梯度方向由径向导数符号决定，场点梯度与源点梯度符号相反}
\end{center}

\par\addvspace{.4\baselineskip}
\noindent\begin{minipage}{\linewidth}

\textbf{例 M1-22：点电荷电势的负梯度}

真空中电荷量为 \(q\) 的点电荷固定在原点，取无穷远电势为零。对
\(r=|\mathbf r|>0\)，其电势为
\(\varphi=q/(4\pi\varepsilon_0r)\)，\(\mathbf e_r\)
规定为背离原点的径向单位矢量。由 \(\mathbf E=-\nabla\varphi\)
求电场，分别说明 \(q>0\) 与 \(q<0\)
时的实际场方向，并指出原点为何不属于普通求导区域。


\end{minipage}\par

\par\noindent\begin{minipage}{\linewidth}
\textbf{解答}

\begin{center}
\includegraphics[width=\linewidth,height=.40\textheight,keepaspectratio]{\detokenize{assets/module-01/examples/M1-22.pdf}}
\par\smallskip{\small 图 M1-22：点电荷电势的负梯度}
\end{center}
\end{minipage}\par

图 M1-22所示的正电荷电势向外降低，所以 \(\nabla\varphi\)
指向内径向；电场是它的反向，指向外径向。若把电荷改成负值，两条箭头的系数同时反号，电场自然改为向内。下面只对
\(r>0\) 求导，原点集中源另由分布理论处理。

点电荷 \(q\) 在原点产生电势

\[
\varphi(r)=\frac{q}{4\pi\varepsilon_0r}.
\]

在 \(r>0\) 处，

\[
\nabla\varphi
=-\frac{q}{4\pi\varepsilon_0r^2}\hat{\mathbf r},
\qquad
\mathbf E=-\nabla\varphi
=\frac{q}{4\pi\varepsilon_0r^2}\hat{\mathbf r}.
\]

若
\(q<0\)，系数自动使电场沿内径向；不应在公式外再人为补一个``向内''的负号。

\subsection{六、势能、电势、保守力与功}\label{m1:ux516dux52bfux80fdux7535ux52bfux4fddux5b88ux529bux4e0eux529f}

\subsubsection{1.
势能与保守力}\label{m1:ux52bfux80fdux4e0eux4fddux5b88ux529b}

若给定粒子的保守力可由势能 \(U(\mathbf r)\) 描述，则保守力做元功

\[
dW_{\mathrm c}=\mathbf F\cdot d\mathbf r=-dU.
\]

因为 \(dU=\nabla U\cdot d\mathbf r\) 对任意微小位移成立，所以

\[
\boxed{
\mathbf F=-\nabla U
}.
\]

从 \(A\) 到 \(B\)，

\[
W_{A\to B}
=\int_A^B\mathbf F\cdot d\mathbf l
=U(A)-U(B).
\]

给 \(U\) 加任意常数不改变 \(\mathbf F\)。

\subsubsection{2. 电势与电场}\label{m1:ux7535ux52bfux4e0eux7535ux573a}

静电场中，电势 \(\varphi\) 定义为单位正试探电荷的势能：

\[
U=q\varphi.
\]

于是

\[
\boxed{
\mathbf E=-\nabla\varphi,
\qquad
\mathbf F=q\mathbf E=-\nabla(q\varphi)
}.
\]

电场总是指向电势下降最快的方向；正电荷受力沿电场，负电荷受力与电场反向。因而``电场指向低电势''不能直接替换为``任意带电粒子都向低势能运动''。

电势差满足

\[
\varphi(B)-\varphi(A)
=-\int_A^B\mathbf E\cdot d\mathbf l.
\]

\begin{center}
\includegraphics[width=\linewidth,height=.40\textheight,keepaspectratio]{\detokenize{assets/module-01/figure-047-revised.pdf}}
\par\smallskip{\small 图示：势能与保守力、电势与电场之间都由负梯度联系但物理对象不同}
\end{center}

\par\addvspace{.4\baselineskip}
\noindent\begin{minipage}{\linewidth}

\textbf{例 M1-23：匀强电场的等势面}

真空中某区域的电场为匀强场 \(\mathbf E=E_0\mathbf e_x\)，其中 \(E_0>0\)
为常量。允许自行选定电势参考值。求与该场相容的电势，说明等势面的形状与法向；比较沿
\(+x\)
前进时的电势变化，并判断相邻等势面若取相同电势差，其几何间距是否相同。


\end{minipage}\par

\par\noindent\begin{minipage}{\linewidth}
\textbf{解答}

\begin{center}
\includegraphics[width=\linewidth,height=.40\textheight,keepaspectratio]{\detokenize{assets/module-01/examples/M1-23.pdf}}
\par\smallskip{\small 图 M1-23：匀强电场的等势面}
\end{center}
\end{minipage}\par

图 M1-23中等势面为垂直于 \(x\) 轴的一组平行截线，电场箭头沿
\(+x\)。因此电势沿 \(x\) 的斜率必须为 \(-E_0\)，积分得到
\(-E_0x+C\)。相同的电势差配上恒定斜率，就给出图中的相同间距；平移零势常数不改变这些几何关系。

匀强电场 \(\mathbf E=E_0\mathbf e_x\) 可取电势

\[
\varphi=-E_0x+C.
\]

等势面为 \(x=\text{常数}\)
的平面，彼此等距；\(-\nabla\varphi=E_0\mathbf e_x\)。常数 \(C\)
不影响电场。

\subsection{七、梯度公式的使用条件}\label{m1:ux4e03ux68afux5ea6ux516cux5f0fux7684ux4f7fux7528ux6761ux4ef6}

使用梯度时依次检查：

\begin{enumerate}
\def\labelenumi{\arabic{enumi}.}
\tightlist
\item
  \textbf{对象}：被求导的是势能、电势、温度还是别的标量场，量纲不能混淆。
\item
  \textbf{可微性}：普通分量公式只用于场在该点可微的情形；尖点、界面、点源处须分区或取单侧极限。
\item
  \textbf{方向单位化}：每单位长度的方向导数必须使用单位方向。
\item
  \textbf{坐标尺度}：柱、球坐标中的角向分量含尺度因子，基矢也随位置变化。
\item
  \textbf{定义域}：\(1/r\)、\(\ln\rho\)
  等公式必须排除奇点；坐标奇异与物理奇异要分开。
\item
  \textbf{正负号}：先由标量随距离的增减判断 \(\nabla\phi\)，再根据
  \(\mathbf F=-\nabla U\) 或 \(\mathbf E=-\nabla\varphi\) 取反。
\item
  \textbf{对称性与极限}：球对称标量场的梯度只能沿径向；最后检查近区、远区与量纲。
\end{enumerate}

\subsection{八、练习与解答}\label{m1:ux516bux7ec3ux4e60ux4e0eux89e3ux7b54}

\par\addvspace{.4\baselineskip}
\noindent\begin{minipage}{\linewidth}

\subsubsection{作业一：方向导数与切平面}\label{m1:ux4f5cux4e1aux4e00ux65b9ux5411ux5bfcux6570ux4e0eux5207ux5e73ux9762}

\textbf{题号：M1-24}

\textbf{层级：} 基础　　\textbf{方法：} 梯度与局部线性化

所有分量相对于右手单位正交基给出；``沿 \(\mathbf v\)
的方向''只规定方向，不表示以 \(|\mathbf v|\) 作为步长。

设

\[
f(x,y,z)=x^2+2y^2+3z^2.
\]

在点 \(P(1,-1,1)\) 处：

\begin{enumerate}
\def\labelenumi{\arabic{enumi}.}
\tightlist
\item
  求增长最快的单位方向与最大增长率；
\item
  求沿 \(\mathbf v=(2,1,-2)\) 方向每单位长度的方向导数；
\item
  求等值面 \(f=f(P)\) 在 \(P\) 处的切平面。
\end{enumerate}


\end{minipage}\par

\par\noindent\begin{minipage}{\linewidth}
\textbf{解答}

\begin{center}
\includegraphics[width=\linewidth,height=.40\textheight,keepaspectratio]{\detokenize{assets/module-01/examples/M1-24.pdf}}
\par\smallskip{\small 图 M1-24：方向导数与切平面}
\end{center}
\end{minipage}\par

图
M1-24左图的水平线是切平面的法截面，蓝色梯度沿正法向，绿色给定方向则带有向下的负法向分量。先由梯度求最大增长率与单位法向，再把
\(\mathbf v\) 除以 \(3\) 后作投影，得到
\(-4\)。切平面方程使用同一个梯度作为法向，所以图中的垂直关系和下面的平面方程必须同时满足。

先求梯度：

\[
\nabla f=(2x,4y,6z),
\qquad
\nabla f(P)=(2,-4,6).
\]

其模为

\[
|\nabla f(P)|=\sqrt{2^2+(-4)^2+6^2}=2\sqrt{14}.
\]

因此增长最快的单位方向为

\[
\mathbf l_{\max}
=\frac{(2,-4,6)}{2\sqrt{14}}
=\frac{(1,-2,3)}{\sqrt{14}},
\]

最大增长率为 \(2\sqrt{14}\)。

\(\mathbf v\) 的模为 \(3\)，对应单位方向为

\[
\mathbf l=\frac13(2,1,-2).
\]

所以

\[
D_{\mathbf l}f(P)
=(2,-4,6)\cdot\frac13(2,1,-2)
=-4.
\]

负号说明沿给定方向标量场在该点下降。

由于 \(\nabla f(P)\ne\mathbf0\)，它是等值面的法向。图 M1-24
的水平切平面使用蓝色梯度作法向，因此切平面为

\[
\nabla f(P)\cdot[(x,y,z)-(1,-1,1)]=0,
\]

即

\[
2(x-1)-4(y+1)+6(z-1)=0,
\]

化简得

\[
\boxed{x-2y+3z=6}.
\]

\par\addvspace{.4\baselineskip}
\noindent\begin{minipage}{\linewidth}

\subsubsection{作业二：柱坐标中的电势梯度}\label{m1:ux4f5cux4e1aux4e8cux67f1ux5750ux6807ux4e2dux7684ux7535ux52bfux68afux5ea6}

\textbf{题号：M1-25}

\textbf{层级：} 基础　　\textbf{方法：} 梯度与局部线性化

圆柱坐标的 \(\rho\) 表示到轴线的距离，\(\rho_0>0\)
为参考长度，\(\mathbf e_{\rho}\) 指向离轴方向。只讨论 \(\rho>0\)
的普通区域。

在 \(\rho>0\) 的区域，电势为

\[
\varphi(\rho,z)
=A\ln\frac{\rho}{\rho_0}-E_0z,
\]

其中 \(A,\rho_0,E_0\) 为常量。求电场，并验证电场垂直等势面。


\end{minipage}\par

\par\noindent\begin{minipage}{\linewidth}
\textbf{解答}

\begin{center}
\includegraphics[width=\linewidth,height=.40\textheight,keepaspectratio]{\detokenize{assets/module-01/examples/M1-25.pdf}}
\par\smallskip{\small 图 M1-25：柱坐标电势的等势面}
\end{center}
\end{minipage}\par

图 M1-25中的等势曲线位于 \(\rho z\) 子午截面，轴线 \(\rho=0\)
被明确排除。沿径向求导给 \(A/\rho\)，沿轴向求导给
\(-E_0\)，整体取负号后在图示 \(A>0,E_0>0\)
的情形，得到电场的左向与上向分量；一般情形按系数符号确定方向。沿图中的等势切线有
\(d\varphi=0\)，因此电场与切线正交；这里正交来自点乘为零，不是从画面目测得来。

电势与方位角 \(\phi\) 无关。柱坐标梯度为

\[
\nabla\varphi
=\frac{\partial\varphi}{\partial\rho}\mathbf e_{\rho}
+\frac1\rho\frac{\partial\varphi}{\partial\phi}\mathbf e_{\phi}
+\frac{\partial\varphi}{\partial z}\mathbf e_z.
\]

三个偏导数为

\[
\frac{\partial\varphi}{\partial\rho}=\frac A\rho,
\qquad
\frac{\partial\varphi}{\partial\phi}=0,
\qquad
\frac{\partial\varphi}{\partial z}=-E_0.
\]

故

\[
\nabla\varphi
=\frac A\rho\mathbf e_{\rho}-E_0\mathbf e_z,
\]

\[
\boxed{
\mathbf E=-\nabla\varphi
=-\frac A\rho\mathbf e_{\rho}+E_0\mathbf e_z
}.
\]

等势面满足

\[
A\ln\frac{\rho}{\rho_0}-E_0z=C.
\]

沿图 M1-25 的等势线作任意切向位移 \(d\mathbf r_t\)，有
\(d\varphi=0\)，故

\[
\nabla\varphi\cdot d\mathbf r_t=0.
\]

所以 \(\nabla\varphi\) 是等势面的法向，而 \(\mathbf E=-\nabla\varphi\)
沿相反法向，同样垂直等势面。定义域 \(\rho>0\) 不可省略，因为 \(\ln\rho\)
与 \(1/\rho\) 在轴线上奇异。

\par\addvspace{.4\baselineskip}
\noindent\begin{minipage}{\linewidth}

\subsubsection{作业三：汤川型势能中的保守力}\label{m1:ux4f5cux4e1aux4e09ux6c64ux5dddux578bux52bfux80fdux4e2dux7684ux4fddux5b88ux529b}

\textbf{题号：M1-26}

\textbf{层级：} 基础　　\textbf{方法：} 梯度与局部线性化

取相对位置矢量 \(\mathbf r\)
从作为力源的一粒子指向受力粒子，\(\mathbf e_r=\mathbf r/r\)
规定为正径向；近区、远区分别按 \(\alpha r\ll1\) 与 \(\alpha r\gg1\)
比较。

设两粒子的相对距离为 \(r>0\)，相互作用势能为

\[
U(r)=U_0\frac{e^{-\alpha r}}{r},
\]

其中 \(U_0>0\)、\(\alpha>0\)。求保守力
\(\mathbf F(\mathbf r)=-\nabla U\)，并讨论其方向、近区与远区行为。


\end{minipage}\par

\par\noindent\begin{minipage}{\linewidth}
\textbf{解答}

\begin{center}
\includegraphics[width=\linewidth,height=.40\textheight,keepaspectratio]{\detokenize{assets/module-01/examples/M1-26.pdf}}
\par\smallskip{\small 图 M1-26：汤川势能与保守力}
\end{center}
\end{minipage}\par

图 M1-26左图把相互作用中心与场点用 \(\mathbf r\)
连起来，右图的力曲线则由同一径向系数给出。势能随距离下降，故
\(dU/dr<0\)，负梯度沿外径向。近区与远区要分别比较图中 \(1/r^2\)
行为和指数衰减，不能把远区也继续当作单纯平方反比。

因为 \(U\) 只依赖 \(r\)，在 \(r>0\) 处

\[
\nabla U=\frac{dU}{dr}\hat{\mathbf r}.
\]

先计算一元导数：

\[
\begin{aligned}
\frac{dU}{dr}
&=U_0\frac{d}{dr}\left(e^{-\alpha r}r^{-1}\right)\\
&=U_0\left[-\alpha e^{-\alpha r}r^{-1}
-e^{-\alpha r}r^{-2}\right]\\
&=-U_0\frac{(1+\alpha r)e^{-\alpha r}}{r^2}.
\end{aligned}
\]

所以

\[
\nabla U
=-U_0\frac{(1+\alpha r)e^{-\alpha r}}{r^2}
\hat{\mathbf r},
\]

\[
\boxed{
\mathbf F(\mathbf r)
=U_0\frac{(1+\alpha r)e^{-\alpha r}}{r^2}
\hat{\mathbf r}
}.
\]

由于 \(U_0,\alpha,r>0\)，力沿 \(+\hat{\mathbf r}\)，对应排斥型势能。若
\(U_0<0\)，整体系数反号，力沿内径向。

先读图 M1-26 的近区。当 \(r\to0\) 时，

\[
e^{-\alpha r}=1-\alpha r+O(r^2),
\]

故

\[
U(r)\sim\frac{U_0}{r},
\qquad
|\mathbf F|\sim\frac{U_0}{r^2}.
\]

再读图 M1-26 的远区。当 \(r\to\infty\) 时，

\[
U(r)\sim U_0\frac{e^{-\alpha r}}r,
\qquad
|\mathbf F|\sim U_0\alpha\frac{e^{-\alpha r}}r.
\]

指数因子使作用在距离远大于 \(1/\alpha\) 时迅速衰减；\(1/\alpha\)
是该模型的特征作用距离。这里 \(U\) 是势能，若题目给的是电势
\(\varphi\)，则粒子受力应写为 \(\mathbf F=q(-\nabla\varphi)\)。

\section{专题四：散度（基础必讲）}\label{m1:ux4e13ux9898ux56dbux6563ux5ea6ux57faux7840ux5fc5ux8bb2}

散度把矢量场在一点附近的``净流出''压缩成一个标量。它回答的是：围住该点的微小闭合面上，向外穿出的通量减去向内穿入的通量后，每单位体积还剩多少。电荷密度、质量守恒、稳恒电流和不可压缩条件都通过这一局部量组织起来。

\subsection{一、对象、局部定义与基本性质}\label{m1:ux4e00ux5bf9ux8c61ux5c40ux90e8ux5b9aux4e49ux4e0eux57faux672cux6027ux8d28}

设矢量场

\[
\mathbf A:\Omega\subset\mathbb R^3\longrightarrow\mathbb R^3
\]

在点 \(P\in\Omega\) 的邻域内属于 \(C^1\)。在固定直角坐标系中写成

\[
\mathbf A=A_x\mathbf e_x+A_y\mathbf e_y+A_z\mathbf e_z.
\]

散度定义为

\[
\boxed{
\nabla\cdot\mathbf A
=\frac{\partial A_x}{\partial x}
+\frac{\partial A_y}{\partial y}
+\frac{\partial A_z}{\partial z}
}.
\]

散度的输出是标量场，而不是矢量场。它的量纲为

\[
[\nabla\cdot\mathbf A]=\frac{[\mathbf A]}{\mathrm L}.
\]

\subsubsection{1.
坐标无关的局部通量定义}\label{m1:ux5750ux6807ux65e0ux5173ux7684ux5c40ux90e8ux901aux91cfux5b9aux4e49}

取一族围住 \(P\) 的闭合小体积
\(V_{\varepsilon}\)，其直径趋于零，并保持形状不退化，例如相似缩小的球、立方体或长宽高比有界的长方体。闭合面
\(\partial V_{\varepsilon}\) 一律取外法向，定义

\[
d\mathbf S=\hat{\mathbf n}\,dS.
\]

则

\[
\boxed{
\nabla\cdot\mathbf A(P)
=\lim_{\varepsilon\to0}
\frac{1}{\operatorname{Vol}(V_{\varepsilon})}
\oint_{\partial V_{\varepsilon}}
\mathbf A\cdot d\mathbf S
}.
\]

这个极限是\textbf{单位体积净向外通量}。向外通量记正，向内通量记负；必须先按外法向计算每一片面的有向通量，再相加。

\begin{center}
\includegraphics[width=\linewidth,height=.40\textheight,keepaspectratio]{\detokenize{assets/module-01/tikz-m1-topic4-v3-01-local-net-flux.jpg}}
\par\smallskip{\small 图示：散度的正负由微小闭合面上的净流出、平衡与净流入决定}
\end{center}

\subsubsection{2.
线性与乘积法则}\label{m1:ux7ebfux6027ux4e0eux4e58ux79efux6cd5ux5219-1}

对常数 \(a,b\)，

\[
\nabla\cdot(a\mathbf A+b\mathbf B)
=a\nabla\cdot\mathbf A+b\nabla\cdot\mathbf B.
\]

对标量场 \(f\) 与矢量场 \(\mathbf A\)，

\[
\boxed{
\nabla\cdot(f\mathbf A)
=\nabla f\cdot\mathbf A
+f\nabla\cdot\mathbf A
}.
\]

第二式说明：乘积的净流出既可能来自矢量场本身的散度，也可能来自标量权重沿场方向的变化。

\par\addvspace{.4\baselineskip}
\noindent\begin{minipage}{\linewidth}

\textbf{例 M1-27：三个方向共同决定散度}

给定光滑矢量场
\(\mathbf A=2x\mathbf e_x-y\mathbf e_y+3z\mathbf e_z\)。在任意点附近取坐标轴对齐的小长方体，采用外法向计算净通量。求该场的散度，分别指出三个坐标方向对净通量的贡献，并说明仅观察
\(y\) 方向的流入能否判断总散度的符号。


\end{minipage}\par

\par\noindent\begin{minipage}{\linewidth}
\textbf{解答}

\begin{center}
\includegraphics[width=\linewidth,height=.40\textheight,keepaspectratio]{\detokenize{assets/module-01/examples/M1-27.pdf}}
\par\smallskip{\small 图 M1-27：三个方向共同决定散度}
\end{center}
\end{minipage}\par

图 M1-27中 \(x\) 向的流出与 \(y\) 向的流入并存，尚有未画在平面内的 \(z\)
向贡献。三个同指标导数必须一起记账：\(2-1+3=4\)。因此局部净通量为正；对本题线性场，闭合小盒的通量恰为
\(4\Delta V\)，没有更高阶余项。

设

\[
\mathbf A=2x\mathbf e_x-y\mathbf e_y+3z\mathbf e_z.
\]

则

\[
\nabla\cdot\mathbf A=2-1+3=4.
\]

因此任一足够小的规则体积 \(\Delta V\)，其净向外通量为

\[
\Delta\Phi=4\Delta V+o(\Delta V).
\]

某一方向上的场分量减小并不能单独决定散度符号，必须把三个法向方向的贡献相加。

\subsection{二、直角坐标公式的分量推导}\label{m1:ux4e8cux76f4ux89d2ux5750ux6807ux516cux5f0fux7684ux5206ux91cfux63a8ux5bfc}

取以 \(P(x,y,z)\) 为中心的微小长方体，三边为
\(\Delta x,\Delta y,\Delta z\)，并令三边按可比速率趋于零：

\[
\Delta V=\Delta x\Delta y\Delta z.
\]

先考察法向为 \(\pm\mathbf e_x\)
的一对面。为避免把非均匀场错误地当作面上常量，应从面积分写起：

\[
\begin{aligned}
\Delta\Phi_x
={}&\int_{-\Delta y/2}^{\Delta y/2}
\int_{-\Delta z/2}^{\Delta z/2}
\Bigl[
A_x\!\left(x+\frac{\Delta x}{2},y+\eta,z+\zeta\right)\\
&\hspace{37mm}
-A_x\!\left(x-\frac{\Delta x}{2},y+\eta,z+\zeta\right)
\Bigr]\,d\zeta\,d\eta.
\end{aligned}
\]

左面外法向为
\(-\mathbf e_x\)，所以两面的合通量自然表现为``右面法向分量减左面法向分量''。由一阶展开及偏导连续性，括号内为

\[
\frac{\partial A_x}{\partial x}(P)\Delta x+o(\Delta x)
\]

并且余项在该微小面上可一致控制。因此

\[
\Delta\Phi_x
=\frac{\partial A_x}{\partial x}(P)\Delta V
+o(\Delta V).
\]

同理，另外两对面分别给出

\[
\Delta\Phi_y
=\frac{\partial A_y}{\partial y}(P)\Delta V
+o(\Delta V),
\]

\[
\Delta\Phi_z
=\frac{\partial A_z}{\partial z}(P)\Delta V
+o(\Delta V).
\]

相加可得

\[
\Delta\Phi
=\left(
\frac{\partial A_x}{\partial x}
+\frac{\partial A_y}{\partial y}
+\frac{\partial A_z}{\partial z}
\right)_P\Delta V
+o(\Delta V).
\]

两边除以 \(\Delta V\) 并令体积趋于零，便得到直角坐标公式。

\begin{center}
\includegraphics[width=\linewidth,height=.40\textheight,keepaspectratio]{\detokenize{assets/module-01/tikz-m1-topic4-v3-02-cartesian-face-pairs.png}}
\par\smallskip{\small 图示：三对相对面分别贡献对应法向分量的偏导数，零阶通量在相对面之间抵消}
\end{center}

只有法向分量参与对应面的通量。例如在 \(x\) 面上，\(A_y\) 与 \(A_z\)
都沿面切向，对该面的点积通量没有贡献。散度中的三个偏导数正是三对相对面留下的一阶差。

\subsubsection{坐标旋转下为什么仍是标量}\label{m1:ux5750ux6807ux65cbux8f6cux4e0bux4e3aux4ec0ux4e48ux4ecdux662fux6807ux91cf}

矢量场在 \(P\) 附近的一阶变化由雅可比矩阵

\[
(D\mathbf A)_{ij}=\frac{\partial A_i}{\partial x_j}
\]

描述，而

\[
\nabla\cdot\mathbf A=\operatorname{tr}(D\mathbf A).
\]

在正交坐标轴旋转下，\(D\mathbf A\)
作相似变换，矩阵的迹不变。因此散度的数值不依赖直角坐标轴怎样旋转，这与局部净通量定义一致。

\par\addvspace{.4\baselineskip}
\noindent\begin{minipage}{\linewidth}

\textbf{例 M1-28：刚性转动场的散度}

设
\(\mathbf r=x\mathbf e_x+y\mathbf e_y+z\mathbf e_z\)，\(\boldsymbol\omega\)
为不随空间位置变化的常矢量。给定
\(\mathbf A(\mathbf r)=\boldsymbol\omega\times\mathbf r\)，可将其视为刚体转动速度场。求散度，并从任意微小体积是否有净流出的角度解释结果；判断闭合或弯曲的场线是否必然意味着存在散度。


\end{minipage}\par

\par\noindent\begin{minipage}{\linewidth}
\textbf{解答}

\begin{center}
\includegraphics[width=\linewidth,height=.40\textheight,keepaspectratio]{\detokenize{assets/module-01/examples/M1-28.pdf}}
\par\smallskip{\small 图 M1-28：刚性转动场的散度}
\end{center}
\end{minipage}\par

图
M1-28中的箭头始终沿圆周切向，对边读取的法向流入、流出相抵，微团可以转动而不膨胀。取转轴为
\(z\) 轴时，场写成
\((-\omega y,\omega x,0)\)，三个同指标导数均为零。这个图说明散度测体积变化，不测局部转动。

若 \(\mathbf A=\boldsymbol\omega\times\mathbf r\)，其中
\(\boldsymbol\omega\)
为常矢量，则该场只表示局部刚性转动。直接按分量计算可得

\[
\nabla\cdot(\boldsymbol\omega\times\mathbf r)=0.
\]

场线可以弯曲甚至闭合，但散度仍可为零；``弯曲''与``净流出''不是同一概念。

\subsection{三、高斯散度定理与局部、整体的联系}\label{m1:ux4e09ux9ad8ux65afux6563ux5ea6ux5b9aux7406ux4e0eux5c40ux90e8ux6574ux4f53ux7684ux8054ux7cfb}

若有界区域 \(V\) 的边界 \(\partial V\) 分片光滑，且 \(\mathbf A\)
在包含闭包 \(\overline V\) 的邻域内属于 \(C^1\)，则

\[
\boxed{
\oint_{\partial V}\mathbf A\cdot d\mathbf S
=\int_V\nabla\cdot\mathbf A\,dV
}.
\]

边界仍取外法向。把区域分割成许多小体积时，相邻小体积共享面的外法向相反，内部通量两两抵消，只剩最外层边界；这给出了定理的几何核心。

\begin{center}
\includegraphics[width=\linewidth,height=.40\textheight,keepaspectratio]{\detokenize{assets/module-01/tikz-m1-topic4-v3-03-gauss-cancellation-average.jpg}}
\par\smallskip{\small 图示：相邻小体积的内部通量抵消，边界总通量等于内部散度的体积分}
\end{center}

两边除以总体积，得到

\[
\frac1{\operatorname{Vol}(V)}
\oint_{\partial V}\mathbf A\cdot d\mathbf S
=\frac1{\operatorname{Vol}(V)}
\int_V\nabla\cdot\mathbf A\,dV.
\]

所以闭合面通量除以体积，是区域内散度的体积平均。若 \(V\)
连通且散度连续，则至少存在一点 \(P_*\in V\) 使

\[
\oint_{\partial V}\mathbf A\cdot d\mathbf S
=(\nabla\cdot\mathbf A)(P_*)\operatorname{Vol}(V).
\]

当 \(V\) 缩到一点时，区域平均趋于该点的局部散度。

\par\addvspace{.4\baselineskip}
\noindent\begin{minipage}{\linewidth}

\textbf{例 M1-29：常矢量场的闭合面通量}

空间中给定常矢量场
\(\mathbf A(\mathbf r)=\mathbf A_0\)。在场内任选闭合曲面
\(S=\partial V\)，所有面元取包围体积的外法向。利用通量积分或散度定理求
\(\oint_S\mathbf A_0\cdot d\mathbf S\)，说明它与体积的位置和形状有无关系，并据此解释该场的局部散度。


\end{minipage}\par

\par\noindent\begin{minipage}{\linewidth}
\textbf{解答}

\begin{center}
\includegraphics[width=\linewidth,height=.40\textheight,keepaspectratio]{\detokenize{assets/module-01/examples/M1-29.pdf}}
\par\smallskip{\small 图 M1-29：常矢量场的闭合面通量}
\end{center}
\end{minipage}\par

图
M1-29中相对两面的外法向相反，而常矢量场在两面上相同；因此一侧流入正好抵消另一侧流出。其余各对面也一样相消，闭合面总通量为零。下面将这个配对写成
\(\mathbf A_0\cdot\oint d\mathbf S=0\)，再取局部极限得到零散度。

若 \(\mathbf A=\mathbf A_0\) 为常矢量场，则
\(\nabla\cdot\mathbf A_0=0\)，所以任意闭合面总通量为零。场本身非零，但从闭合体积一侧流入的通量会从另一侧流出。

\subsubsection{定理的使用边界}\label{m1:ux5b9aux7406ux7684ux4f7fux7528ux8fb9ux754c}

若区域内含点源、线源、薄层或场的跳跃界面，则不能把``在光滑区域成立''的微分式直接跨过特殊位置积分。可采用三种办法：

\begin{enumerate}
\def\labelenumi{\arabic{enumi}.}
\tightlist
\item
  把奇点或界面从区域中挖去，在各光滑子区域分别使用定理；
\item
  用薄球、薄柱或薄盒的积分极限提取集中贡献；
\item
  使用包含狄拉克函数或面狄拉克分布的分布形式。
\end{enumerate}

\subsection{四、正交曲线坐标中的散度}\label{m1:ux56dbux6b63ux4ea4ux66f2ux7ebfux5750ux6807ux4e2dux7684ux6563ux5ea6}

设 \((q_1,q_2,q_3)\) 是正交曲线坐标，单位基矢为
\(\mathbf e_1,\mathbf e_2,\mathbf e_3\)，尺度因子为

\[
h_i=\left|\frac{\partial\mathbf r}{\partial q_i}\right|.
\]

矢量场写成

\[
\mathbf A=A_1\mathbf e_1+A_2\mathbf e_2+A_3\mathbf e_3.
\]

微小曲线坐标盒的体积为

\[
dV=h_1h_2h_3\,dq_1dq_2dq_3.
\]

法向沿 \(\mathbf e_1\) 的一对面，其面积因子为
\(h_2h_3\,dq_2dq_3\)。两面对通量的一阶净贡献为

\[
\frac{\partial}{\partial q_1}(h_2h_3A_1)
dq_1dq_2dq_3.
\]

循环加入另外两对面，再除以体积，得到

\[
\boxed{
\nabla\cdot\mathbf A
=\frac1{h_1h_2h_3}
\left[
\frac{\partial}{\partial q_1}(h_2h_3A_1)
+\frac{\partial}{\partial q_2}(h_3h_1A_2)
+\frac{\partial}{\partial q_3}(h_1h_2A_3)
\right]
}.
\]

\begin{center}
\includegraphics[width=\linewidth,height=.40\textheight,keepaspectratio]{\detokenize{assets/module-01/tikz-m1-topic4-v3-04-orthogonal-coordinate-scale.jpg}}
\par\smallskip{\small 图示：曲线坐标中的每对面既包含场分量变化，也包含面面积因子的变化}
\end{center}

\subsubsection{1. 直角坐标}\label{m1:ux76f4ux89d2ux5750ux6807}

\[
(h_x,h_y,h_z)=(1,1,1),
\]

故回到

\[
\nabla\cdot\mathbf A
=\frac{\partial A_x}{\partial x}
+\frac{\partial A_y}{\partial y}
+\frac{\partial A_z}{\partial z}.
\]

\subsubsection{\texorpdfstring{2. 柱坐标
\((\rho,\phi,z)\)}{2. 柱坐标 (\textbackslash rho,\textbackslash phi,z)}}\label{m1:ux67f1ux5750ux6807-rhophiz}

\[
(h_{\rho},h_{\phi},h_z)=(1,\rho,1),
\]

\[
\boxed{
\nabla\cdot\mathbf A
=\frac1\rho\frac{\partial(\rho A_{\rho})}{\partial\rho}
+\frac1\rho\frac{\partial A_{\phi}}{\partial\phi}
+\frac{\partial A_z}{\partial z}
},
\qquad \rho>0.
\]

\subsubsection{\texorpdfstring{3. 球坐标
\((r,\theta,\phi)\)}{3. 球坐标 (r,\textbackslash theta,\textbackslash phi)}}\label{m1:ux7403ux5750ux6807-rthetaphi}

\[
(h_r,h_{\theta},h_{\phi})=(1,r,r\sin\theta),
\]

\[
\boxed{
\nabla\cdot\mathbf A
=\frac1{r^2}\frac{\partial(r^2A_r)}{\partial r}
+\frac1{r\sin\theta}
\frac{\partial(\sin\theta A_{\theta})}{\partial\theta}
+\frac1{r\sin\theta}
\frac{\partial A_{\phi}}{\partial\phi}
}.
\]

球坐标公式在 \(r=0\)、柱坐标公式在
\(\rho=0\)、球坐标角向公式在极轴处具有坐标奇异性。坐标表达式失效不一定代表物理场发散，必须改用直角坐标或积分极限判断。

\par\addvspace{.4\baselineskip}
\noindent\begin{minipage}{\linewidth}

\textbf{例 M1-30：柱径向场的面积因子}

在圆柱坐标中给定轴对称径向场 \(\mathbf A=k\rho\mathbf e_{\rho}\)，其中
\(k\) 为常量，\(\rho\) 为到 \(z\) 轴的距离，场与方位角及 \(z\)
无关。求散度，并以同轴薄圆柱壳的内、外侧面积随 \(\rho\)
改变为依据，解释柱坐标散度中的面积因子，不能只对径向分量作普通一维求导。


\end{minipage}\par

\par\noindent\begin{minipage}{\linewidth}
\textbf{解答}

\begin{center}
\includegraphics[width=\linewidth,height=.40\textheight,keepaspectratio]{\detokenize{assets/module-01/examples/M1-30.pdf}}
\par\smallskip{\small 图 M1-30：柱径向场的面积因子}
\end{center}
\end{minipage}\par

图 M1-30中两条同心圆是同轴薄壳的内、外侧面，\(d\rho\)
连接两面。外侧不但场强随 \(\rho\) 增加，侧面积也含一个
\(\rho\)，所以要先对 \(\rho A_{\rho}=k\rho^2\) 求导，再除以体积因子
\(\rho\)。这两个图示因素合起来给出 \(2k\)。

柱坐标场

\[
\mathbf A=k\rho\,\mathbf e_{\rho}
\]

满足

\[
\nabla\cdot\mathbf A
=\frac1\rho\frac{d(k\rho^2)}{d\rho}
=2k.
\]

若漏掉外面的 \(1/\rho\) 或括号中的面积因子 \(\rho\)，都会得到错误结果。

\subsection{五、球对称径向场与点源奇点}\label{m1:ux4e94ux7403ux5bf9ux79f0ux5f84ux5411ux573aux4e0eux70b9ux6e90ux5947ux70b9}

\subsubsection{1.
一般球对称径向场}\label{m1:ux4e00ux822cux7403ux5bf9ux79f0ux5f84ux5411ux573a}

设

\[
\mathbf A=F(r)\hat{\mathbf r},
\qquad r>0.
\]

由球坐标散度公式，

\[
\boxed{
\nabla\cdot\bigl(F(r)\hat{\mathbf r}\bigr)
=\frac1{r^2}\frac{d}{dr}\bigl(r^2F(r)\bigr)
},
\qquad r>0.
\]

其物理意义直接来自球壳：半径为 \(r\) 的球面通量为

\[
\Phi(r)=4\pi r^2F(r).
\]

若相邻球面间没有体源，通量应与 \(r\) 无关，因此
\(r^2F(r)=\text{常数}\)，即 \(F(r)\propto1/r^2\)。

若写成

\[
\mathbf A=g(r)\mathbf r,
\]

则 \(F(r)=rg(r)\)，所以

\[
\boxed{
\nabla\cdot[g(r)\mathbf r]
=3g(r)+rg'(r)
},
\qquad r>0.
\]

\begin{center}
\includegraphics[width=\linewidth,height=.40\textheight,keepaspectratio]{\detokenize{assets/module-01/tikz-m1-topic4-v3-05-radial-shell-law.png}}
\par\smallskip{\small 图示：球对称径向场的散度由相邻球面的通量差决定，平方反比是无体源幂律}
\end{center}

\par\addvspace{.4\baselineskip}
\noindent\begin{minipage}{\linewidth}

\textbf{例 M1-31：球径向线性场的散度}

三维空间中给定球对称线性场 \(\mathbf A=k\mathbf r=kr\mathbf e_r\)，\(k\)
为常量。求普通散度；再取以原点为球心、半径为 \(R>0\)
的球面，使用外法向直接计算通量，与球内散度的体积分核对，并解释系数与空间维数的联系。


\end{minipage}\par

\par\noindent\begin{minipage}{\linewidth}
\textbf{解答}

\begin{center}
\includegraphics[width=\linewidth,height=.40\textheight,keepaspectratio]{\detokenize{assets/module-01/examples/M1-31.pdf}}
\par\smallskip{\small 图 M1-31：球径向线性场的散度}
\end{center}
\end{minipage}\par

图 M1-31中的球面外法向与径向场同向，球面通量为 \(4\pi r^2(kr)\)。把这个
\(r^3\) 通量对半径的变化除以球壳体积，就得到
\(3k\)。它与直角坐标下三个方向各贡献一个 \(k\)
一致，和前一题柱对称模型的面积因子不同。

若 \(\mathbf A=k\mathbf r=kr\hat{\mathbf r}\)，则

\[
\nabla\cdot\mathbf A
=\frac1{r^2}\frac{d(kr^3)}{dr}=3k.
\]

这与直角坐标下 \(\nabla\cdot(kx,ky,kz)=3k\) 完全一致。

\subsubsection{2.
平方反比场的普通函数结论}\label{m1:ux5e73ux65b9ux53cdux6bd4ux573aux7684ux666eux901aux51fdux6570ux7ed3ux8bba}

对

\[
\mathbf A=C\frac{\mathbf r}{r^3}
=\frac C{r^2}\hat{\mathbf r},
\]

在任意 \(r>0\) 的普通点，

\[
\nabla\cdot\mathbf A
=\frac1{r^2}\frac{dC}{dr}=0.
\]

但任意包围原点的球面都有

\[
\oint_{S_r}\mathbf A\cdot d\mathbf S
=4\pi C.
\]

``原点外散度为零''只说明普通体源在 \(r>0\)
区域内为零，不能据此抹去原点的集中源。

\subsubsection{3.
分布形式（拔高选读）}\label{m1:ux5206ux5e03ux5f62ux5f0fux62d4ux9ad8ux9009ux8bfb}

在整个三维空间中，平方反比场的散度应写成

\[
\boxed{
\nabla\cdot\left(C\frac{\mathbf r}{r^3}\right)
=4\pi C\,\delta^{(3)}(\mathbf r)
}.
\]

球面通量可以迅速确定右边的系数。若要严格验证分布恒等式，取任意光滑且具有紧支集的测试函数
\(\psi(\mathbf r)\)，并从空间中挖去半径为 \(\varepsilon\)
的小球。由于原点外散度为零，分部积分给出

\[
\begin{aligned}
\left\langle
\nabla\cdot\left(C\frac{\mathbf r}{r^3}\right),\psi
\right\rangle
&=-\lim_{\varepsilon\to0^+}
\int_{r>\varepsilon}
C\frac{\mathbf r}{r^3}\cdot\nabla\psi\,dV\\
&=\lim_{\varepsilon\to0^+}
\frac C{\varepsilon^2}
\int_{r=\varepsilon}\psi\,dS\\
&=4\pi C\,\psi(\mathbf0).
\end{aligned}
\]

这正是 \(4\pi C\delta^{(3)}(\mathbf r)\)
对测试函数的作用。对不包含原点的体积，两边都为零。

\begin{center}
\includegraphics[width=\linewidth,height=.40\textheight,keepaspectratio]{\detokenize{assets/module-01/tikz-m1-topic4-v3-06-inverse-square-singularity.png}}
\par\smallskip{\small 图示：平方反比场必须分成原点外普通结论、包围原点的通量和全空间分布结论}
\end{center}

在穿孔区域 \(\mathbb R^3\setminus\{0\}\)
中，场处处散度为零，但包围原点的球面不能直接作为该区域中某个完整体积的边界，因为其内部包含被排除的点。高斯定理没有失效，是使用区域的前提发生了变化。

\subsection{六、连续性方程与体积膨胀率}\label{m1:ux516dux8fdeux7eedux6027ux65b9ux7a0bux4e0eux4f53ux79efux81a8ux80c0ux7387}

\subsubsection{1.
固定控制体积中的守恒}\label{m1:ux56faux5b9aux63a7ux5236ux4f53ux79efux4e2dux7684ux5b88ux6052}

设 \(\varrho(\mathbf r,t)\) 是某种守恒量的体密度，\(\mathbf J\)
是其通量密度，\(s\) 是单位体积单位时间的产生率。对任意固定体积
\(V\)，守恒关系为

\[
\frac{d}{dt}\int_V\varrho\,dV
+\oint_{\partial V}\mathbf J\cdot d\mathbf S
=\int_V s\,dV.
\]

第二项是净向外流失率。利用散度定理并把各项移入同一体积分：

\[
\int_V
\left(
\frac{\partial\varrho}{\partial t}
+\nabla\cdot\mathbf J-s
\right)dV=0.
\]

由于 \(V\) 任意，得到局部连续性方程

\[
\boxed{
\frac{\partial\varrho}{\partial t}
+\nabla\cdot\mathbf J=s
}.
\]

无源守恒时 \(s=0\)。

\begin{center}
\includegraphics[width=\linewidth,height=.40\textheight,keepaspectratio]{\detokenize{assets/module-01/tikz-m1-topic4-v3-07-continuity-equation.jpg}}
\par\smallskip{\small 图示：固定控制体积内储量变化、边界净流出和内部产生率共同构成连续性方程}
\end{center}

\subsubsection{2.
速度场与随体体积}\label{m1:ux901fux5ea6ux573aux4e0eux968fux4f53ux4f53ux79ef}

若守恒量随流体运动，\(\mathbf J=\varrho\mathbf v\)。无源条件给出

\[
\frac{\partial\varrho}{\partial t}
+\nabla\cdot(\varrho\mathbf v)=0.
\]

展开后为

\[
\frac{D\varrho}{Dt}
+\varrho\nabla\cdot\mathbf v=0,
\qquad
\frac D{Dt}=\frac\partial{\partial t}+\mathbf v\cdot\nabla.
\]

对随体微元 \(dV\)，

\[
\boxed{
\frac1{dV}\frac{D(dV)}{Dt}
=\nabla\cdot\mathbf v
}.
\]

因此 \(\nabla\cdot\mathbf v>0\)
表示体积元膨胀，负值表示压缩。对满足质量守恒且 \(\varrho>0\) 的流体，

\[
\frac{D\varrho}{Dt}=0
\quad\Longleftrightarrow\quad
\nabla\cdot\mathbf v=0.
\]

所以不可压缩条件是随体体积不变，而不是``流体静止''或``速度处处相同''；``散度为正''也不等于物质凭空产生。

\par\addvspace{.4\baselineskip}
\noindent\begin{minipage}{\linewidth}

\textbf{例 M1-32：变截面流管的流率}

某守恒量在一根无分支流管中稳恒输运，流量密度为
\(\mathbf J\)。管内没有体源或汇，侧壁不允许该守恒量穿过。取两张与流向垂直的横截面，面积分别为
\(S_1,S_2\)；在每张截面上近似认为 \(\mathbf J\)
均匀且沿截面法向。比较两截面的总流率和 \(J_1,J_2\)，说明当 \(S_2<S_1\)
时通量密度的变化是否意味着管内产生了源。


\end{minipage}\par

\par\noindent\begin{minipage}{\linewidth}
\textbf{解答}

\begin{center}
\includegraphics[width=\linewidth,height=.40\textheight,keepaspectratio]{\detokenize{assets/module-01/examples/M1-32.pdf}}
\par\smallskip{\small 图 M1-32：变截面流管的流率}
\end{center}
\end{minipage}\par

图
M1-32中的侧壁沿流线，因此侧壁法向通量为零。将两张截面统一按流向取法向，就有流入截面的积分等于流出截面的积分。只有在每个截面上近似均匀时，才能进一步写成
\(J_1S_1=J_2S_2\)；图中小截面的箭头较长，表示通量密度增大而总流率保持不变。

稳恒流动若满足
\(\nabla\cdot\mathbf J=0\)，则任意两张横截面与侧壁围成的流管中，流入量等于流出量。截面变窄时，通量密度可以增大，但流管内部仍无源。

\subsection{七、电磁学中的散度和界面条件}\label{m1:ux4e03ux7535ux78c1ux5b66ux4e2dux7684ux6563ux5ea6ux548cux754cux9762ux6761ux4ef6}

\subsubsection{1.
电场与磁场的散度约束}\label{m1:ux7535ux573aux4e0eux78c1ux573aux7684ux6563ux5ea6ux7ea6ux675f}

真空中的麦克斯韦方程包含

\[
\boxed{
\nabla\cdot\mathbf E=\frac{\rho}{\varepsilon_0}
},
\]

\[
\boxed{
\nabla\cdot\mathbf B=0
}.
\]

前式说明电荷密度是电场的散度源；后式表示标准麦克斯韦理论中没有磁荷源项。二者都是任意时刻必须满足的约束，不只适用于静态场。

在介质中还常用

\[
\nabla\cdot\mathbf D=\rho_{\mathrm f},
\]

其中右边只计自由电荷密度。使用哪一个场，决定右边统计哪一类源，不能把
\(\mathbf E\) 与 \(\mathbf D\) 的高斯关系混写。

\begin{center}
\includegraphics[width=\linewidth,height=.40\textheight,keepaspectratio]{\detokenize{assets/module-01/figure-059-revised.pdf}}
\par\smallskip{\small 图示：电场、电位移和磁场的散度分别对应总电荷、自由电荷与零磁荷源项}
\end{center}

\subsubsection{2.
界面上的跳跃与面源}\label{m1:ux754cux9762ux4e0aux7684ux8df3ux8dc3ux4e0eux9762ux6e90}

设光滑界面 \(S\) 把空间分成负侧和正侧，单位法向 \(\hat{\mathbf n}\)
从负侧指向正侧。场的单侧极限分别为 \(\mathbf A_-\) 与 \(\mathbf A_+\)。

取跨越界面的薄柱，令厚度趋于零。侧面通量比上下底面低一阶，于是

\[
\lim_{h\to0}
\int_{V_h}\nabla\cdot\mathbf A\,dV
=\bigl(\mathbf A_+-\mathbf A_-\bigr)\cdot\hat{\mathbf n}\,\Delta S.
\]

若界面具有面源密度 \(\sigma_A\)，则

\[
\boxed{
\hat{\mathbf n}\cdot(\mathbf A_+-\mathbf A_-)=\sigma_A
}.
\]

在分布意义下可写成

\[
\nabla\cdot\mathbf A
=(\nabla\cdot\mathbf A)_{\mathrm{bulk}}
+\hat{\mathbf n}\cdot(\mathbf A_+-\mathbf A_-)
\delta_S.
\]

其中 \((\nabla\cdot\mathbf A)_{\mathrm{bulk}}\)
表示界面两侧各自的普通散度，面狄拉克分布 \(\delta_S\) 由

\[
\int_{\mathbb R^3}f\,\delta_S\,dV
=\int_S f\,dS
\]

定义。若没有面源，法向分量连续。切向分量的跳跃不进入散度的薄柱条件，而应由旋度及薄回路条件检查。

\par\addvspace{.4\baselineskip}
\noindent\begin{minipage}{\linewidth}

\textbf{例 M1-33：带面电荷界面的法向跳跃}

平面 \(z=0\) 将空间分为下侧``\(-\)区''和上侧``\(+\)区''，统一界面法向取
\(\hat{\mathbf n}=\mathbf e_z\)。界面自由面电荷密度为常量
\(\sigma_f\)，紧邻界面的下侧电位移为
\(\mathbf D_-=D_0\mathbf e_z\)，并已知上侧也只有法向分量。利用跨界薄高斯柱求
\(\mathbf D_+\)，说明哪些电荷进入电位移的法向跃变关系。


\end{minipage}\par

\par\noindent\begin{minipage}{\linewidth}
\textbf{解答}

\begin{center}
\includegraphics[width=\linewidth,height=.40\textheight,keepaspectratio]{\detokenize{assets/module-01/examples/M1-33.pdf}}
\par\smallskip{\small 图 M1-33：带面电荷界面的法向跳跃}
\end{center}
\end{minipage}\par

图 M1-33中红色薄柱的上盖、下盖区域外法向相反，分别贡献 \(D_+\Delta S\)
与 \(-D_-\Delta S\)。令柱高趋零后侧面贡献消失，柱内自由面电荷为
\(\sigma_f\Delta S\)。按图中从下侧到上侧的 \(\hat{\mathbf n}\)
取向相减，便得到 \(D_+-D_- =\sigma_f\)；反转法向时必须交换两侧标签。

平面界面的法向取 \(\hat{\mathbf n}=\mathbf e_z\)。若
\(\mathbf D_-=D_0\mathbf e_z\)，界面自由面电荷密度为
\(\sigma_{\mathrm f}\)，且场均沿法向，则

\[
\mathbf D_+=(D_0+\sigma_{\mathrm f})\mathbf e_z,
\]

因为
\(\mathbf e_z\cdot(\mathbf D_+-\mathbf D_-)=\sigma_{\mathrm f}\)。若反选法向，正负侧标签也必须同时交换，物理结论不变。

\subsubsection{3. 常见误区}\label{m1:ux5e38ux89c1ux8befux533a}

\begin{enumerate}
\def\labelenumi{\arabic{enumi}.}
\tightlist
\item
  \textbf{场线张开不等于散度必为正}：必须计算闭合面净通量的体积密度。
\item
  \textbf{场不为零不等于散度不为零}：常矢量场和刚性转动场都是反例。
\item
  \textbf{散度为零不等于场为常量}：平方反比场在源点外就是非均匀无散场。
\item
  \textbf{局部无源不等于全空间无源}：奇点、穿孔区域和界面可能携带集中源。
\item
  \textbf{不能忘记外法向}：闭合面每一片的符号都由外法向决定。
\item
  \textbf{曲线坐标不能逐分量照抄直角坐标公式}：面积与体积尺度因子必须进入散度。
\item
  \textbf{微分式不能直接写在跳跃点上}：先取单侧极限、薄柱积分或使用分布形式。
\end{enumerate}

\subsection{八、练习与解答}\label{m1:ux516bux7ec3ux4e60ux4e0eux89e3ux7b54-1}

\par\addvspace{.4\baselineskip}
\noindent\begin{minipage}{\linewidth}

\subsubsection{作业一：有心幂律场的散度与无源条件}\label{m1:ux4f5cux4e1aux4e00ux6709ux5fc3ux5e42ux5f8bux573aux7684ux6563ux5ea6ux4e0eux65e0ux6e90ux6761ux4ef6}

\textbf{题号：M1-34}

\textbf{层级：} 基础　　\textbf{方法：} 高斯定理

把 \(\mathbf A\) 视为三维数学矢量场，\(\mathbf e_r\)
沿外径向；普通散度仅在 \(r>0\) 计算，涉及原点时另按通量和分布处理。

已知球对称矢量场

\[
\mathbf A(\mathbf r)=r^n\mathbf r,
\qquad r=|\mathbf r|>0,
\]

其中 \(n\) 为常数。

\begin{enumerate}
\def\labelenumi{\arabic{enumi}.}
\tightlist
\item
  求 \(\nabla\cdot\mathbf A\)；
\item
  求使该场在 \(r>0\) 区域处处无源的 \(n\)；
\item
  说明该特殊幂律的物理意义，并判断原点是否可以直接忽略。
\end{enumerate}


\end{minipage}\par

\par\noindent\begin{minipage}{\linewidth}
\textbf{解答}

\begin{center}
\includegraphics[width=\linewidth,height=.40\textheight,keepaspectratio]{\detokenize{assets/module-01/examples/M1-34.pdf}}
\par\smallskip{\small 图 M1-34：幂律场的散度与通量}
\end{center}
\end{minipage}\par

图 M1-34中的球面是与幂律场同心的积分面，法向场分量为
\(r^{n+1}\)。因此通量先读成 \(4\pi r^{n+3}\)，再用球壳体积计算散度，得
\((n+3)r^n\)。\(n=-3\)
时图上任意球面的通量相同，却不意味着包围原点的总源为零；原点项须在最后单独补回。

利用乘积法则

\[
\nabla\cdot(f\mathbf A)
=\nabla f\cdot\mathbf A+f\nabla\cdot\mathbf A,
\]

取 \(f=r^n\)，并注意

\[
\nabla\cdot\mathbf r=3,
\qquad
\nabla(r^n)=nr^{n-2}\mathbf r.
\]

于是

\[
\begin{aligned}
\nabla\cdot(r^n\mathbf r)
&=\nabla(r^n)\cdot\mathbf r
+r^n\nabla\cdot\mathbf r\\
&=nr^{n-2}(\mathbf r\cdot\mathbf r)+3r^n\\
&=(n+3)r^n.
\end{aligned}
\]

因此

\[
\boxed{
\nabla\cdot\mathbf A=(n+3)r^n
},
\qquad r>0.
\]

要使 \(r>0\) 处散度恒为零，必须有

\[
n+3=0,
\]

所以

\[
\boxed{n=-3}.
\]

此时

\[
\mathbf A=\frac{\mathbf r}{r^3}
=\frac{\hat{\mathbf r}}{r^2},
\]

即三维平方反比径向场，是库仑场和牛顿引力场的几何幂律。结论只在 \(r>0\)
成立；全空间分布形式为

\[
\nabla\cdot\frac{\mathbf r}{r^3}
=4\pi\delta^{(3)}(\mathbf r).
\]

所以原点不能直接忽略。若场前有常系数 \(C\)，右边相应乘以 \(C\)。

\par\addvspace{.4\baselineskip}
\noindent\begin{minipage}{\linewidth}

\subsubsection{作业二：柱坐标场与闭合面通量核对}\label{m1:ux4f5cux4e1aux4e8cux67f1ux5750ux6807ux573aux4e0eux95edux5408ux9762ux901aux91cfux6838ux5bf9}

\textbf{题号：M1-35}

\textbf{层级：} 基础　　\textbf{方法：} 高斯定理

圆柱面所有通量都按所包围实体的外法向取向。轴线不属于给定普通场的定义域，使用散度定理时须先明确挖孔边界与极限。

在 \(\rho>0\) 区域给定

\[
\mathbf A
=a\rho\,\mathbf e_{\rho}
+\frac c\rho\,\mathbf e_{\phi}
+bz\,\mathbf e_z,
\]

其中 \(a,b,c\) 为常数。

\begin{enumerate}
\def\labelenumi{\arabic{enumi}.}
\tightlist
\item
  求该场的散度；
\item
  对半径为 \(R\)、高度为 \(L\)、轴线沿 \(z\) 轴且上下底面位于
  \(z=\pm L/2\)
  的闭合圆柱，把轴心奇点用挖去小圆盘后的极限处理，直接计算总通量；
\item
  正确处理轴线奇异性后与散度定理核对，并说明 \(c/\rho\)
  项为何不贡献散度型集中源与该圆柱通量。
\end{enumerate}


\end{minipage}\par

\par\noindent\begin{minipage}{\linewidth}
\textbf{解答}

\begin{center}
\includegraphics[width=\linewidth,height=.40\textheight,keepaspectratio]{\detokenize{assets/module-01/examples/M1-35.pdf}}
\par\smallskip{\small 图 M1-35：柱坐标场的闭合面通量}
\end{center}
\end{minipage}\par

图
M1-35左图的径向分量穿过圆柱侧面，轴向分量穿过上下底，方位分量沿所有边界相切而不贡献通量。于是径向与轴向通量分别是
\(2a\pi R^2L\) 和 \(b\pi R^2L\)，与散度 \(2a+b\) 的体积分相同。由于
\(c/\rho\)
在轴线上奇异，使用经典定理时先挖去细轴；其径向通量仍为零，最后可取回该极限。

柱坐标散度为

\[
\nabla\cdot\mathbf A
=\frac1\rho\frac{\partial(\rho A_{\rho})}{\partial\rho}
+\frac1\rho\frac{\partial A_{\phi}}{\partial\phi}
+\frac{\partial A_z}{\partial z}.
\]

代入

\[
A_{\rho}=a\rho,
\qquad
A_{\phi}=\frac c\rho,
\qquad
A_z=bz,
\]

得到

\[
\nabla\cdot\mathbf A
=\frac1\rho\frac{d(a\rho^2)}{d\rho}
+0+b
=2a+b.
\]

所以

\[
\boxed{\nabla\cdot\mathbf A=2a+b}.
\]

场在轴线上没有普通定义，因此以下闭合面通量按广义曲面积分理解：先从上下底面各挖去半径为
\(\varepsilon\) 的轴心圆盘，计算后再令
\(\varepsilon\to0^+\)。圆柱侧面的外法向为
\(\mathbf e_{\rho}\)，侧面通量为

\[
\Phi_{\mathrm{side}}
=A_{\rho}(R)(2\pi RL)
=aR(2\pi RL)
=2\pi aR^2L.
\]

上底面 \(z=L/2\) 的外法向为 \(+\mathbf e_z\)：

\[
\Phi_{\mathrm{top}}(\varepsilon)
=b\frac L2\,\pi\left(R^2-\varepsilon^2\right).
\]

下底面 \(z=-L/2\) 的外法向为 \(-\mathbf e_z\)，故

\[
\Phi_{\mathrm{bottom}}(\varepsilon)
=\left(-b\frac L2\mathbf e_z\right)
\cdot(-\mathbf e_z)\pi\left(R^2-\varepsilon^2\right)
=b\frac L2\pi\left(R^2-\varepsilon^2\right).
\]

令 \(\varepsilon\to0^+\)，得到总通量

\[
\Phi
=2\pi aR^2L+b\pi R^2L
=\pi R^2L(2a+b).
\]

由于 \(c/\rho\,\mathbf e_{\phi}\)
在轴线上没有普通定义，不能直接援引要求场在闭包上连续可微的散度定理。先挖去半径为
\(\varepsilon\) 的同轴细管，在环柱区域

\[
V_{\varepsilon}:\qquad
\varepsilon<\rho<R,
\qquad
-\frac L2<z<\frac L2
\]

上应用散度定理。内圆柱面的外法向为 \(-\mathbf e_{\rho}\)，故内侧面通量为

\[
\Phi_{\mathrm{\text{内}}}
=a\varepsilon\,\mathbf e_{\rho}\cdot(-\mathbf e_{\rho})
\,(2\pi\varepsilon L)
=-2\pi a\varepsilon^2L.
\]

外侧面通量仍为 \(2\pi aR^2L\)，上下两个环形底面的总通量为

\[
\Phi_{\mathrm{\text{底}}}
=b\pi\left(R^2-\varepsilon^2\right)L.
\]

于是

\[
\Phi_{\varepsilon}
=(2a+b)\pi\left(R^2-\varepsilon^2\right)L.
\]

另一方面，\(V_{\varepsilon}\) 的体积为 \(\pi(R^2-\varepsilon^2)L\)，所以

\[
\int_{V_{\varepsilon}}\nabla\cdot\mathbf A\,dV
=(2a+b)\pi\left(R^2-\varepsilon^2\right)L
=\Phi_{\varepsilon}.
\]

令 \(\varepsilon\to0^+\)，得到原闭合圆柱的总通量

\[
\Phi=\pi R^2L(2a+b).
\]

\(c/\rho\,\mathbf e_{\phi}\)
在外侧面、内侧面和上下底面上都沿切向，因此始终没有通量；它也与方位角
\(\phi\) 无关，故
\(\rho^{-1}\partial A_{\phi}/\partial\phi=0\)。上述挖孔极限还表明它不产生散度型线源；它属于环流型奇异结构，涉及旋度时必须另行检查。

\par\addvspace{.4\baselineskip}
\noindent\begin{minipage}{\linewidth}

\subsubsection{作业三：球对称稳恒流与集中源}\label{m1:ux4f5cux4e1aux4e09ux7403ux5bf9ux79f0ux7a33ux6052ux6d41ux4e0eux96c6ux4e2dux6e90}

\textbf{题号：M1-36}

\textbf{层级：} 基础　　\textbf{方法：} 高斯定理

规定 \(Q>0\) 表示向外输运，每个同心球面的外法向为
\(+\mathbf e_r\)；流率为穿过整个球面的总量，不能与单位面积流量密度混同。

某种守恒量从原点以恒定总流率 \(Q\)
向各方向均匀流出。除原点外空间没有体源，流量密度为

\[
\mathbf J=J(r)\hat{\mathbf r}.
\]

\begin{enumerate}
\def\labelenumi{\arabic{enumi}.}
\tightlist
\item
  由局部守恒求 \(J(r)\)；
\item
  写出 \(r>0\) 区域的散度；
\item
  写出包含原点的全空间分布形式。
\end{enumerate}


\end{minipage}\par

\par\noindent\begin{minipage}{\linewidth}
\textbf{解答}

\begin{center}
\includegraphics[width=\linewidth,height=.40\textheight,keepaspectratio]{\detokenize{assets/module-01/examples/M1-36.pdf}}
\par\smallskip{\small 图 M1-36：球对称稳恒流与集中源}
\end{center}
\end{minipage}\par

图 M1-36中每个同心球面都截取同一个总流率 \(Q\)，而球面面积为
\(4\pi r^2\)，因此径向密度只能是
\(Q/(4\pi r^2)\)。在两球面之间没有源，普通散度为零；将图中的中心也包括进来时，有限的总流率由
\(Q\delta^{(3)}(\mathbf r)\) 表示。

球对称且 \(r>0\) 无体源，因此

\[
\nabla\cdot\mathbf J
=\frac1{r^2}\frac{d}{dr}(r^2J)=0.
\]

于是

\[
r^2J=C,
\]

其中 \(C\) 为常数。半径为 \(r\) 的球面总流率为

\[
\oint_{S_r}\mathbf J\cdot d\mathbf S
=4\pi r^2J=Q,
\]

故

\[
C=\frac Q{4\pi}.
\]

所以

\[
\boxed{
\mathbf J
=\frac Q{4\pi r^2}\hat{\mathbf r}
=\frac Q{4\pi}\frac{\mathbf r}{r^3}
},
\qquad r>0.
\]

在原点外，

\[
\boxed{\nabla\cdot\mathbf J=0},
\qquad r>0.
\]

但任意包围原点的闭合面总流率为 \(Q\)，因此全空间分布形式为

\[
\boxed{
\nabla\cdot\mathbf J
=Q\,\delta^{(3)}(\mathbf r)
}.
\]

这清楚地区分了``原点外无体源''和``原点存在集中源''两个层次。

\section{专题五：旋度（基础必讲）}\label{m1:ux4e13ux9898ux4e94ux65cbux5ea6ux57faux7840ux5fc5ux8bb2}

旋度把矢量场在一点附近的``环绕趋势''压缩成一个矢量。旋度非零时，它的方向给出局部环量最强的有向面法向，沿任意单位法向的分量则等于该方向上的单位面积环量。法拉第电磁感应、安培--麦克斯韦定律、无旋场与势函数以及流体的局部转动，都由这一概念联系起来。

\subsection{一、对象、取向与局部定义}\label{m1:ux4e00ux5bf9ux8c61ux53d6ux5411ux4e0eux5c40ux90e8ux5b9aux4e49}

设矢量场

\[
\mathbf A:\Omega\subset\mathbb R^3\longrightarrow\mathbb R^3
\]

在点 \(P\in\Omega\) 的邻域内属于 \(C^1\)。对有向闭合曲线 \(C\)，定义环量

\[
\Gamma_C=\oint_C\mathbf A\cdot d\mathbf l.
\]

环量只读取场在曲线切向上的分量。若同时反转曲线方向，则

\[
\Gamma_{-C}=-\Gamma_C.
\]

\subsubsection{1.
有向小面上的环量密度}\label{m1:ux6709ux5411ux5c0fux9762ux4e0aux7684ux73afux91cfux5bc6ux5ea6}

给定单位矢量 \(\hat{\mathbf n}\)，取一族以 \(P\) 为内部点、法向为
\(\hat{\mathbf n}\) 的微小平面区域
\(S_{\varepsilon}\)。它们相似缩小，或至少保持长宽比有界。边界
\(\partial S_{\varepsilon}\)
的正向由右手定则确定：右手四指沿边界正向弯曲时，大拇指指向
\(\hat{\mathbf n}\)。

旋度沿 \(\hat{\mathbf n}\) 的分量定义为

\[
\boxed{
(\nabla\times\mathbf A)(P)\cdot\hat{\mathbf n}
=
\lim_{\varepsilon\to0}
\frac{1}{\operatorname{Area}(S_{\varepsilon})}
\oint_{\partial S_{\varepsilon}}
\mathbf A\cdot d\mathbf l
}.
\]

对三个线性无关方向取得这些投影，便唯一确定矢量
\(\nabla\times\mathbf A\)。旋度由各个方向的有向面环量密度共同确定，其方向和大小都受局部环量约束。

\begin{center}
\includegraphics[width=\linewidth,height=.40\textheight,keepaspectratio]{\detokenize{assets/module-01/tikz-m1-topic5-v3-01-local-circulation-density-safe-v7.png}}
\par\smallskip{\small 图示：有向小面、边界正向和旋度法向分量共同决定局部环量密度}
\end{center}

旋度的量纲为

\[
[\nabla\times\mathbf A]=\frac{[\mathbf A]}{\mathrm L}.
\]

如果把 \(\hat{\mathbf n}\)
反向，边界正向也必须同时反向。此时环量和旋度的法向分量同时变号，几何关系保持一致。

\subsubsection{2.
基本运算性质}\label{m1:ux57faux672cux8fd0ux7b97ux6027ux8d28}

对常数 \(a,b\)，

\[
\nabla\times(a\mathbf A+b\mathbf B)
=a\nabla\times\mathbf A+b\nabla\times\mathbf B.
\]

对标量场 \(f\) 与矢量场 \(\mathbf A\)，

\[
\boxed{
\nabla\times(f\mathbf A)
=\nabla f\times\mathbf A
+f\nabla\times\mathbf A
}.
\]

第二式说明，乘积场的局部环绕既可能来自 \(\mathbf A\)
本身，也可能来自标量权重在垂直于 \(\mathbf A\) 的方向上发生变化。

\par\addvspace{.4\baselineskip}
\noindent\begin{minipage}{\linewidth}

\textbf{例 M1-37：常矢量场的闭合环量}

在常矢量场 \(\mathbf A=\mathbf A_0\) 内取任意可缩小的光滑闭合回路
\(C\)，回路由连续的有向线元 \(d\mathbf l\) 构成。计算
\(\oint_C\mathbf A_0\cdot d\mathbf l\)，再从单位面积定向环量的极限说明该场的旋度；讨论回路平面取向改变是否影响结论。


\end{minipage}\par

\par\noindent\begin{minipage}{\linewidth}
\textbf{解答}

\begin{center}
\includegraphics[width=\linewidth,height=.40\textheight,keepaspectratio]{\detokenize{assets/module-01/examples/M1-37.pdf}}
\par\smallskip{\small 图 M1-37：常矢量场的闭合环量}
\end{center}
\end{minipage}\par

沿图
M1-37中的红色回路，下边与上边的切向取向相反，但常场相同，两个积分抵消；两条竖边也作同样配对。等价地，闭合回路的位移总和为零。由图上的零环量再取单位面积极限，得到零旋度，而不是把非零场误判为零。

常矢量场 \(\mathbf A=\mathbf A_0\) 沿任意闭合小回路的环量都为零，因此

\[
\nabla\times\mathbf A_0=\mathbf0.
\]

场本身非零并不意味着旋度非零。

\subsection{二、直角坐标公式的分量推导}\label{m1:ux4e8cux76f4ux89d2ux5750ux6807ux516cux5f0fux7684ux5206ux91cfux63a8ux5bfc-1}

在固定直角坐标系中写成

\[
\mathbf A
=A_x\mathbf e_x+A_y\mathbf e_y+A_z\mathbf e_z.
\]

先求旋度的 \(z\) 分量。以 \(P(x,y,z)\) 为中心，在 \(xy\) 平面内取边长为
\(\Delta x,\Delta y\) 的微小矩形，并选取法向 \(+\mathbf e_z\)；令
\(\Delta x,\Delta y\) 保持有界的长宽比并同时趋于零。从 \(+\mathbf e_z\)
方向看，边界正向为逆时针。

四条边的线积分可合并为

\[
\begin{aligned}
\Gamma_z
={}&
\int_{-\Delta x/2}^{\Delta x/2}
\left[
A_x\left(x+\xi,y-\frac{\Delta y}{2},z\right)
-A_x\left(x+\xi,y+\frac{\Delta y}{2},z\right)
\right]\,d\xi\\
&+
\int_{-\Delta y/2}^{\Delta y/2}
\left[
A_y\left(x+\frac{\Delta x}{2},y+\eta,z\right)
-A_y\left(x-\frac{\Delta x}{2},y+\eta,z\right)
\right]\,d\eta.
\end{aligned}
\]

第一组括号来自下边与上边，第二组来自右边与左边。利用一阶展开及偏导连续性，

\[
\Gamma_z
=
\left(
\frac{\partial A_y}{\partial x}
-\frac{\partial A_x}{\partial y}
\right)_P
\Delta x\Delta y
+o(\Delta x\Delta y).
\]

两边除以面积并令矩形缩到 \(P\)，得到

\[
\boxed{
(\nabla\times\mathbf A)_z
=
\frac{\partial A_y}{\partial x}
-\frac{\partial A_x}{\partial y}
}.
\]

\begin{center}
\includegraphics[width=\linewidth,height=.40\textheight,keepaspectratio]{\detokenize{assets/module-01/tikz-m1-topic5-v3-02-cartesian-loop-derivation.jpg}}
\par\smallskip{\small 图示：微小矩形四条边的零阶环量抵消，只留下切向分量的横向变化}
\end{center}

对法向分别取 \(+\mathbf e_x,+\mathbf e_y\) 的微小矩形作循环推导，得到

\[
\boxed{
\nabla\times\mathbf A
=
\left(
\frac{\partial A_z}{\partial y}
-\frac{\partial A_y}{\partial z}
\right)\mathbf e_x
+
\left(
\frac{\partial A_x}{\partial z}
-\frac{\partial A_z}{\partial x}
\right)\mathbf e_y
+
\left(
\frac{\partial A_y}{\partial x}
-\frac{\partial A_x}{\partial y}
\right)\mathbf e_z
}.
\]

常见的行列式写法

\[
\nabla\times\mathbf A
=
\begin{vmatrix}
\mathbf e_x&\mathbf e_y&\mathbf e_z\\
\partial_x&\partial_y&\partial_z\\
A_x&A_y&A_z
\end{vmatrix}
\]

只是记忆三个分量符号的形式工具，不是普通数值行列式，也不能代替有向环量定义。

\begin{center}
\includegraphics[width=\linewidth,height=.40\textheight,keepaspectratio]{\detokenize{assets/module-01/tikz-m1-topic5-v3-03-component-orientation.jpg}}
\par\smallskip{\small 图示：三个旋度分量分别由垂直于对应坐标轴的小回路测得}
\end{center}

\subsubsection{坐标旋转与轴矢量性质（拔高选读）}\label{m1:ux5750ux6807ux65cbux8f6cux4e0eux8f74ux77e2ux91cfux6027ux8d28ux62d4ux9ad8ux9009ux8bfb}

用列维--奇维塔符号可写成

\[
(\nabla\times\mathbf A)_i
=\varepsilon_{ijk}\frac{\partial A_k}{\partial x_j}.
\]

在正交坐标变换 \(\mathbf x'=Q\mathbf x\) 下，

\[
\boxed{
(\nabla'\times\mathbf A')
=\det(Q)\,Q(\nabla\times\mathbf A)
}.
\]

对保持右手性的转动，\(\det Q=1\)，旋度像普通矢量一样转动；对镜面反射，
\(\det Q=-1\)，还会多出一个符号。因此极矢量场的旋度是轴矢量或赝矢量。这与``旋转轴在镜像下的变化方式''一致，也说明使用左手坐标系时不能直接照搬右手系的叉乘符号。若原场本身是轴矢量，例如
\(\mathbf B\)，则再取一次旋度会得到极矢量；这与
\(\nabla\times\mathbf B\) 和电流密度 \(\mathbf J\) 具有相同宇称相一致。

\par\addvspace{.4\baselineskip}
\noindent\begin{minipage}{\linewidth}

\textbf{例 M1-38：二维转动场的旋度}

给定二维场
\(\mathbf A=-\Omega y\mathbf e_x+\Omega x\mathbf e_y\)，\(\Omega\)
为常量。求 \(\nabla\times\mathbf A\)；在 \(xy\)
平面取微小闭合回路，规定从 \(+z\)
方向观察逆时针为正向，利用该正向的环量判断旋度法向分量的符号，并与分量计算一致性核对。


\end{minipage}\par

\par\noindent\begin{minipage}{\linewidth}
\textbf{解答}

\begin{center}
\includegraphics[width=\linewidth,height=.40\textheight,keepaspectratio]{\detokenize{assets/module-01/examples/M1-38.pdf}}
\par\smallskip{\small 图 M1-38：二维转动场的旋度}
\end{center}
\end{minipage}\par

图
M1-38中逆时针场线对应纸外正法向。读取小回路的两组相对边，\(\partial_xA_y=\Omega\)
与 \(-\partial_yA_x=\Omega\) 同号相加，所以旋度为
\(2\Omega\mathbf e_z\)。\(\Omega<0\)
时场箭头反向，式中的法向系数也随之反号。

二维转动场

\[
\mathbf A=-\Omega y\,\mathbf e_x+\Omega x\,\mathbf e_y
\]

满足

\[
\nabla\times\mathbf A
=
\left[
\frac{\partial(\Omega x)}{\partial x}
-\frac{\partial(-\Omega y)}{\partial y}
\right]\mathbf e_z
=2\Omega\,\mathbf e_z.
\]

若 \(\Omega>0\)，从 \(+z\) 方向看场沿逆时针环绕，旋度指向
\(+\mathbf e_z\)，符号与右手定则一致。

\subsection{三、旋度与局部转动的物理图像}\label{m1:ux4e09ux65cbux5ea6ux4e0eux5c40ux90e8ux8f6cux52a8ux7684ux7269ux7406ux56feux50cf}

对速度场 \(\mathbf v(\mathbf r)\)，一点附近的速度差可写成

\[
\delta\mathbf v
=(D\mathbf v)\,\delta\mathbf r.
\]

速度梯度可唯一分解为对称部分与反对称部分：

\[
\frac{\partial v_i}{\partial x_j}
=
\frac12
\left(
\frac{\partial v_i}{\partial x_j}
+\frac{\partial v_j}{\partial x_i}
\right)
+
\frac12
\left(
\frac{\partial v_i}{\partial x_j}
-\frac{\partial v_j}{\partial x_i}
\right).
\]

对称部分描述局部伸长与剪切；反对称部分等效于刚性微小转动。局部角速度为

\[
\boxed{
\boldsymbol\omega_{\text{局部}}
=\frac12\nabla\times\mathbf v
}.
\]

因此旋度等于局部角速度的两倍，而不是角速度本身。只有在速度场讨论中，才能把旋度直接解释为局部微团的转动；对一般电场、磁场或抽象矢量场，更稳妥的定义仍是单位面积环量。

\begin{center}
\includegraphics[width=\linewidth,height=.40\textheight,keepaspectratio]{\detokenize{assets/module-01/tikz-m1-topic5-v3-04-rotation-versus-shape.jpg}}
\par\smallskip{\small 图示：匀速平移、刚体转动和径向场的场线外观不同，旋度必须由小回路环量判断}
\end{center}

\par\addvspace{.4\baselineskip}
\noindent\begin{minipage}{\linewidth}

\textbf{例 M1-39：刚体转动的局部角速度}

刚体绕通过原点的固定轴转动，其角速度矢量 \(\boldsymbol\omega\)
在空间中为常量，质点速度满足
\(\mathbf v=\boldsymbol\omega\times\mathbf r\)。求速度场的旋度，并解释由微小流体元或刚性线元观察到的局部角速度与旋度之间的倍数关系。


\end{minipage}\par

\par\noindent\begin{minipage}{\linewidth}
\textbf{解答}

\begin{center}
\includegraphics[width=\linewidth,height=.40\textheight,keepaspectratio]{\detokenize{assets/module-01/examples/M1-39.pdf}}
\par\smallskip{\small 图 M1-39：刚体转动的局部角速度}
\end{center}
\end{minipage}\par

以图 M1-39中的转轴作为 \(z\)
轴，圆周切向速度随半径线性增大。小回路两组对边各贡献一个角速度，故旋度为
\(2\boldsymbol\omega\)。下面恢复一般矢量写法后，局部角速度应读成旋度的一半，而不是把纸外旋度箭头本身当作角速度数值。

刚体以常角速度 \(\boldsymbol\omega\) 转动时

\[
\mathbf v=\boldsymbol\omega\times\mathbf r,
\]

直接计算可得

\[
\nabla\times\mathbf v=2\boldsymbol\omega.
\]

\par\addvspace{.4\baselineskip}
\noindent\begin{minipage}{\linewidth}

\textbf{例 M1-40：简单剪切中的转动与形变}

在直角坐标系中给定简单剪切速度场 \(\mathbf v=\gamma y\mathbf e_x\)，其中
\(\gamma\)
为常量。求旋度，并以旋度的一半定义局部刚性转动角速度；同时比较处于不同
\(y\)
位置的相邻质点速度，说明为什么仅有角速度仍不足以描述该速度场的局部运动。


\end{minipage}\par

\par\noindent\begin{minipage}{\linewidth}
\textbf{解答}

\begin{center}
\includegraphics[width=\linewidth,height=.40\textheight,keepaspectratio]{\detokenize{assets/module-01/examples/M1-40.pdf}}
\par\smallskip{\small 图 M1-40：简单剪切中的转动与形变}
\end{center}
\end{minipage}\par

图 M1-40中 \(y=0\)
是零速度线，上半部向右、下半部向左，形成顺时针局部转动。红色平行四边形来自
\(\Delta x=\gamma y\Delta t\)，同时表现剪切形变。按图的右手取向，旋度沿
\(-\mathbf e_z\)，局部角速度为
\(-\gamma\mathbf e_z/2\)；形变还要由速度梯度的对称部分描述。

简单剪切流

\[
\mathbf v=\gamma y\,\mathbf e_x
\]

满足

\[
\nabla\times\mathbf v
=-\gamma\,\mathbf e_z,
\qquad
\boldsymbol\omega_{\text{局部}}
=-\frac{\gamma}{2}\mathbf e_z.
\]

小正方形微团在运动中既会转动，也会剪切成平行四边形。旋度只读取速度梯度的反对称转动部分，不能代替对称部分所描述的形变。

匀速平移场虽然速度处处非零，却没有局部自转；径向场的场线虽会张开，也不因此自动具有旋度。``场线弯曲''\,``场强变化''和``局部环量''是三个不同概念。

\subsection{四、斯托克斯定理与局部、整体的联系}\label{m1:ux56dbux65afux6258ux514bux65afux5b9aux7406ux4e0eux5c40ux90e8ux6574ux4f53ux7684ux8054ux7cfb}

设 \(S\) 是分片光滑、可定向的曲面，边界为 \(\partial S\)；\(\mathbf A\)
在包含 \(S\) 的邻域内属于
\(C^1\)。选定曲面法向后，边界正向由右手定则绑定，则

\[
\boxed{
\oint_{\partial S}\mathbf A\cdot d\mathbf l
=
\int_S(\nabla\times\mathbf A)\cdot d\mathbf S
}.
\]

把曲面分割成许多小面元时，相邻面元共享边的正向相反，内部线积分两两抵消，只剩最外层边界。这正是斯托克斯定理的几何核心。

\begin{center}
\includegraphics[width=\linewidth,height=.40\textheight,keepaspectratio]{\detokenize{assets/module-01/tikz-m1-topic5-v3-05-stokes-cancellation.jpg}}
\par\smallskip{\small 图示：曲面剖分后内部公共边的线积分抵消，只留下外边界和旋度通量}
\end{center}

\subsubsection{1.
取向与曲面选择}\label{m1:ux53d6ux5411ux4e0eux66f2ux9762ux9009ux62e9}

同一条无向闭合曲线有两个绕行方向。改变曲面法向时，必须同时改变边界正向，等式两边才会一起变号。

若两个曲面 \(S_1,S_2\) 具有同一条有向边界，且 \(\mathbf A\)
在两个曲面及其间相关区域内均保持所需光滑性，则

\[
\int_{S_1}(\nabla\times\mathbf A)\cdot d\mathbf S
=
\int_{S_2}(\nabla\times\mathbf A)\cdot d\mathbf S
=
\oint_{\partial S_1}\mathbf A\cdot d\mathbf l.
\]

这里的``曲面可以任选''带有重要前提：替换后的曲面不能穿过场的奇线、跳跃层或未定义区域。

\subsubsection{2.
从积分式回到局部定义}\label{m1:ux4eceux79efux5206ux5f0fux56deux5230ux5c40ux90e8ux5b9aux4e49}

令 \(S\) 缩到普通点 \(P\)，若旋度连续，则

\[
\frac1{\operatorname{Area}(S)}
\oint_{\partial S}\mathbf A\cdot d\mathbf l
=
\frac1{\operatorname{Area}(S)}
\int_S(\nabla\times\mathbf A)\cdot\hat{\mathbf n}\,dS
\longrightarrow
(\nabla\times\mathbf A)(P)\cdot\hat{\mathbf n}.
\]

因此局部环量密度定义与斯托克斯定理完全一致。

\par\addvspace{.4\baselineskip}
\noindent\begin{minipage}{\linewidth}

\textbf{例 M1-41：单位旋度与圆盘面积}

在 \(xy\) 平面中给定
\(\mathbf A=-(y/2)\mathbf e_x+(x/2)\mathbf e_y\)。取圆盘
\(S:x^2+y^2\le R^2\)，\(R>0\)，法向为 \(+\mathbf e_z\)，边界圆周 \(C\)
按从 \(+z\) 看逆时针取向。求 \(\oint_C\mathbf A\cdot d\mathbf l\)，并与
\(S\) 上的旋度通量比较，解释环量与圆盘面积的关系。


\end{minipage}\par

\par\noindent\begin{minipage}{\linewidth}
\textbf{解答}

\begin{center}
\includegraphics[width=\linewidth,height=.40\textheight,keepaspectratio]{\detokenize{assets/module-01/examples/M1-41.pdf}}
\par\smallskip{\small 图 M1-41：单位旋度与圆盘面积}
\end{center}
\end{minipage}\par

图 M1-41中半径 \(R\) 的圆周上场沿切向，大小恒为 \(R/2\)，所以直接积分是
\(R/2\) 乘以周长 \(2\pi R\)。纸外法向的旋度恒为
\(1\)，斯托克斯路线则把它乘以圆盘面积。两种读图路线给出的同一个数都是
\(\pi R^2\)。

取

\[
\mathbf A=-\frac y2\mathbf e_x+\frac x2\mathbf e_y.
\]

在 \(xy\) 平面内半径为 \(R\) 的圆周上，

\[
\oint_C\mathbf A\cdot d\mathbf l
=\pi R^2.
\]

另一方面 \(\nabla\times\mathbf A=\mathbf e_z\)，故圆盘上的旋度通量也是
\(\pi R^2\)。

\subsubsection{3.
定理的使用边界}\label{m1:ux5b9aux7406ux7684ux4f7fux7528ux8fb9ux754c-1}

遇到以下情况时，不能把普通点上的微分式不加检查地跨区域积分：

\begin{enumerate}
\def\labelenumi{\arabic{enumi}.}
\tightlist
\item
  曲面穿过点奇点、线奇点或场的未定义位置；
\item
  场在界面两侧发生跳跃；
\item
  区域有孔，使闭合曲线不能在定义域内连续收缩为一点；
\item
  选取的曲面不可定向，或边界方向与法向不匹配。
\end{enumerate}

处理方法是挖去奇异集合、加入小边界并取极限，或使用包含狄拉克分布的旋度形式。

\subsection{五、二阶恒等式、势函数与拓扑边界}\label{m1:ux4e94ux4e8cux9636ux6052ux7b49ux5f0fux52bfux51fdux6570ux4e0eux62d3ux6251ux8fb9ux754c}

\subsubsection{1. 梯度的旋度}\label{m1:ux68afux5ea6ux7684ux65cbux5ea6}

若标量场 \(\phi\in C^2\)，则混合偏导可交换，从而

\[
\boxed{
\nabla\times(\nabla\phi)=\mathbf0
}.
\]

例如其 \(z\) 分量为

\[
\frac{\partial^2\phi}{\partial x\partial y}
-\frac{\partial^2\phi}{\partial y\partial x}
=0.
\]

因此梯度场必为无旋场。若力场写成
\(\mathbf F=-\nabla U\)，则负号不改变其无旋性。

\subsubsection{2. 旋度的散度}\label{m1:ux65cbux5ea6ux7684ux6563ux5ea6}

若 \(\mathbf A\in C^2\)，则

\[
\boxed{
\nabla\cdot(\nabla\times\mathbf A)=0
}.
\]

展开后三对混合偏导分别抵消。这个恒等式是局部光滑结论；若场含线奇点或跳跃面，必须把集中分布贡献一并计入。

若 \(\mathbf A=\nabla\times\mathbf C\)，则
\(\nabla\cdot\mathbf A=0\)。逆命题需要额外的整体条件：在星形或可缩区域中，足够光滑的无散场可以构造矢势；仅有``区域单连通''仍不充分。例如

\[
\mathbf A=\frac{\mathbf r}{r^3}
\]

在 \(\mathbb R^3\setminus\{\mathbf0\}\)
中处处无散，但任意包围原点的球面通量为
\(4\pi\)，因此不可能在整个穿孔区域写成某个处处光滑、单值矢势的旋度。

\subsubsection{3.
无旋场何时存在单值势函数}\label{m1:ux65e0ux65cbux573aux4f55ux65f6ux5b58ux5728ux5355ux503cux52bfux51fdux6570}

若区域 \(\Omega\) 开且单连通，\(\mathbf A\in C^1(\Omega)\)，并且

\[
\nabla\times\mathbf A=\mathbf0,
\]

则任意两点间的线积分与路径无关，可以定义单值势函数

\[
\Psi(\mathbf r)
=
\int_{\mathbf r_0}^{\mathbf r}\mathbf A\cdot d\mathbf l,
\qquad
\mathbf A=\nabla\Psi.
\]

严格地说，取从 \(\mathbf r_0\) 到 \(\mathbf r\) 的两条路径
\(\gamma_1,\gamma_2\)，把 \(\gamma_1\) 与反向的 \(\gamma_2\)
拼成闭合曲线 \(C\)。单连通性保证 \(C\) 可在 \(\Omega\)
内连续收缩，并可用位于定义域内的曲面或同伦扫过面应用斯托克斯定理，因此

\[
\int_{\gamma_1}\mathbf A\cdot d\mathbf l
-\int_{\gamma_2}\mathbf A\cdot d\mathbf l
=\oint_C\mathbf A\cdot d\mathbf l
=\int_S(\nabla\times\mathbf A)\cdot d\mathbf S
=0.
\]

所以 \(\Psi\) 与所选路径无关。再把终点从 \(\mathbf r\) 移到
\(\mathbf r+\Delta\mathbf r\)，沿两点间的短线段积分：

\[
\Psi(\mathbf r+\Delta\mathbf r)-\Psi(\mathbf r)
=\mathbf A(\mathbf r)\cdot\Delta\mathbf r
+o(|\Delta\mathbf r|).
\]

由梯度的局部线性定义即得 \(\nabla\Psi=\mathbf A\)。结合
\(\nabla\times\nabla\Psi=\mathbf0\)，便得到单连通光滑区域内``无旋''和``存在单值标量势''的充要关系。

在有孔区域中，``普通点处处无旋''不再足以保证全局单值势。典型反例是

\[
\mathbf A
=
\frac{\kappa}{2\pi\rho}\mathbf e_{\phi},
\qquad
\rho>0.
\]

在轴线外，

\[
\nabla\times\mathbf A=\mathbf0,
\]

但绕 \(z\) 轴一周的任意圆周都有

\[
\oint_C\mathbf A\cdot d\mathbf l=\kappa.
\]

它只能局部写成 \(\mathbf A=\nabla(\kappa\phi/2\pi)\)；方位角 \(\phi\)
绕轴一周增加 \(2\pi\)，所以这个势不是全局单值函数。

\begin{center}
\includegraphics[width=\linewidth,height=.40\textheight,keepaspectratio]{\detokenize{assets/module-01/tikz-m1-topic5-v3-06-topology-line-singularity.jpg}}
\par\smallskip{\small 图示：不包围轴线的回路可缩小，而环绕奇线的回路保持非零环量}
\end{center}

\subsubsection{4.
线奇点的分布形式（拔高选读）}\label{m1:ux7ebfux5947ux70b9ux7684ux5206ux5e03ux5f62ux5f0fux62d4ux9ad8ux9009ux8bfb}

若把定义域扩展到全空间，则上述环流场满足

\[
\boxed{
\nabla\times
\left(
\frac{\kappa}{2\pi\rho}\mathbf e_{\phi}
\right)
=
\kappa\,\delta(x)\delta(y)\,\mathbf e_z
}.
\]

对任意横截面圆盘积分，右边给出 \(\kappa\)，圆盘边界环量也为
\(\kappa\)；二者符合分布形式的斯托克斯定理。普通斯托克斯定理不能直接用于包含轴线的圆盘，因为原场在轴线上没有定义。这说明``轴外旋度为零''和``轴线上存在集中涡量''并不矛盾。

\subsection{六、正交曲线坐标中的旋度}\label{m1:ux516dux6b63ux4ea4ux66f2ux7ebfux5750ux6807ux4e2dux7684ux65cbux5ea6}

设 \((q_1,q_2,q_3)\) 为右手正交曲线坐标，单位基矢为
\(\mathbf e_1,\mathbf e_2,\mathbf e_3\)，尺度因子为

\[
h_i=\left|\frac{\partial\mathbf r}{\partial q_i}\right|,
\]

且

\[
\mathbf A=A_1\mathbf e_1+A_2\mathbf e_2+A_3\mathbf e_3.
\]

对法向为 \(\mathbf e_1\) 的微小曲线坐标面，边界环量的一阶项除以面积
\(h_2h_3\,dq_2dq_3\)，得到

\[
(\nabla\times\mathbf A)_1
=
\frac1{h_2h_3}
\left[
\frac{\partial(h_3A_3)}{\partial q_2}
-\frac{\partial(h_2A_2)}{\partial q_3}
\right].
\]

循环轮换指标可得

\[
\boxed{
\nabla\times\mathbf A
=
\sum_{\mathrm{cyc}}
\frac{\mathbf e_i}{h_jh_k}
\left[
\frac{\partial(h_kA_k)}{\partial q_j}
-\frac{\partial(h_jA_j)}{\partial q_k}
\right]
},
\]

其中循环次序为 \((i,j,k)=(1,2,3),(2,3,1),(3,1,2)\)。

\begin{center}
\includegraphics[width=\linewidth,height=.40\textheight,keepaspectratio]{\detokenize{assets/module-01/tikz-m1-topic5-v3-07-orthogonal-coordinate-curl.png}}
\par\smallskip{\small 图示：曲线坐标小面上的边长尺度因子同时进入环量和面积}
\end{center}

\subsubsection{\texorpdfstring{1. 柱坐标
\((\rho,\phi,z)\)}{1. 柱坐标 (\textbackslash rho,\textbackslash phi,z)}}\label{m1:ux67f1ux5750ux6807-rhophiz-1}

尺度因子为

\[
(h_{\rho},h_{\phi},h_z)=(1,\rho,1),
\]

所以

\[
\boxed{
\begin{aligned}
(\nabla\times\mathbf A)_{\rho}
&=\frac1\rho\frac{\partial A_z}{\partial\phi}
-\frac{\partial A_{\phi}}{\partial z},\\
(\nabla\times\mathbf A)_{\phi}
&=\frac{\partial A_{\rho}}{\partial z}
-\frac{\partial A_z}{\partial\rho},\\
(\nabla\times\mathbf A)_z
&=\frac1\rho
\left[
\frac{\partial(\rho A_{\phi})}{\partial\rho}
-\frac{\partial A_{\rho}}{\partial\phi}
\right].
\end{aligned}
}
\qquad \rho>0.
\]

\subsubsection{\texorpdfstring{2. 球坐标
\((r,\theta,\phi)\)}{2. 球坐标 (r,\textbackslash theta,\textbackslash phi)}}\label{m1:ux7403ux5750ux6807-rthetaphi-1}

尺度因子为

\[
(h_r,h_{\theta},h_{\phi})=(1,r,r\sin\theta),
\]

所以

\[
\boxed{
\begin{aligned}
(\nabla\times\mathbf A)_r
&=\frac1{r\sin\theta}
\left[
\frac{\partial(\sin\theta A_{\phi})}{\partial\theta}
-\frac{\partial A_{\theta}}{\partial\phi}
\right],\\
(\nabla\times\mathbf A)_{\theta}
&=\frac1r
\left[
\frac1{\sin\theta}\frac{\partial A_r}{\partial\phi}
-\frac{\partial(rA_{\phi})}{\partial r}
\right],\\
(\nabla\times\mathbf A)_{\phi}
&=\frac1r
\left[
\frac{\partial(rA_{\theta})}{\partial r}
-\frac{\partial A_r}{\partial\theta}
\right].
\end{aligned}
}
\]

柱坐标公式在 \(\rho=0\)、球坐标公式在 \(r=0\)
及极轴处具有坐标奇异性。不能只看公式中出现 \(1/\rho\) 或
\(1/\sin\theta\)
就断言物理场发散；应改用直角坐标、环量极限或分布方法判断。

\par\addvspace{.4\baselineskip}
\noindent\begin{minipage}{\linewidth}

\textbf{例 M1-42：柱切向场的旋度}

在圆柱坐标 \((\rho,\phi,z)\) 中给定
\(\mathbf A=k\rho\mathbf e_{\phi}\)，其中 \(k\)
为常量，\(\mathbf e_{\phi}\)
沿方位角增大的方向。求该场的旋度；将它改写为直角坐标分量，与二维刚性转动场的旋度比较，核对方向与系数。


\end{minipage}\par

\par\noindent\begin{minipage}{\linewidth}
\textbf{解答}

\begin{center}
\includegraphics[width=\linewidth,height=.40\textheight,keepaspectratio]{\detokenize{assets/module-01/examples/M1-42.pdf}}
\par\smallskip{\small 图 M1-42：柱切向场的旋度}
\end{center}
\end{minipage}\par

图 M1-42中的场沿 \(\mathbf e_{\phi}\)，因此读取垂直于 \(z\)
轴的圆周环量。坐标式要对 \(\rho A_{\phi}=k\rho^2\) 求导，再除以
\(\rho\)，得到
\(2k\)。这正是图示刚性转动场的两组边贡献，另外两个旋度分量都为零。

若

\[
\mathbf A=k\rho\,\mathbf e_{\phi},
\]

则

\[
(\nabla\times\mathbf A)_z
=
\frac1\rho\frac{d(k\rho^2)}{d\rho}
=2k,
\]

其余分量为零。这与二维刚体转动场的直角坐标结果一致。

\subsection{七、电磁学中的旋度方程与界面条件}\label{m1:ux4e03ux7535ux78c1ux5b66ux4e2dux7684ux65cbux5ea6ux65b9ux7a0bux4e0eux754cux9762ux6761ux4ef6}

\subsubsection{1. 法拉第定律}\label{m1:ux6cd5ux62c9ux7b2cux5b9aux5f8b}

对固定在空间中的回路及其张成曲面，积分形式与微分形式分别为

\[
\boxed{
\oint_{\partial S}\mathbf E\cdot d\mathbf l
=
-\frac{d}{dt}\int_S\mathbf B\cdot d\mathbf S
},
\]

\[
\boxed{
\nabla\times\mathbf E
=
-\frac{\partial\mathbf B}{\partial t}
}.
\]

负号体现楞次定律：感生电场的环量方向总是反抗磁通量的变化。这里的边界正向与磁通量法向必须满足同一右手定则。

若导线回路本身运动或变形，总电动势还含 \(\mathbf v\times\mathbf B\)
的运动电动势贡献。在通常条件下应使用

\[
\oint_{C(t)}
\left(
\mathbf E+\mathbf v\times\mathbf B
\right)\cdot d\mathbf l
=-\frac d{dt}
\int_{S(t)}\mathbf B\cdot d\mathbf S.
\]

因此不能把固定回路的电场环量公式直接搬到运动回路上而忽略
\(\mathbf v\times\mathbf B\)。

\subsubsection{2.
安培--麦克斯韦定律}\label{m1:ux5b89ux57f9ux9ea6ux514bux65afux97e6ux5b9aux5f8b}

真空中

\[
\boxed{
\oint_{\partial S}\mathbf B\cdot d\mathbf l
=
\mu_0\int_S\mathbf J\cdot d\mathbf S
+\mu_0\varepsilon_0
\frac{d}{dt}\int_S\mathbf E\cdot d\mathbf S
},
\]

\[
\boxed{
\nabla\times\mathbf B
=
\mu_0\mathbf J
+\mu_0\varepsilon_0
\frac{\partial\mathbf E}{\partial t}
}.
\]

在介质中常写成

\[
\boxed{
\nabla\times\mathbf H
=
\mathbf J_{\mathrm f}
+\frac{\partial\mathbf D}{\partial t}
}.
\]

只有在稳恒情形下，位移电流项才可消失。把
\(\nabla\times\mathbf B=\mu_0\mathbf J\)
当作任意时变过程中的完整方程，会破坏电荷守恒。

\begin{center}
\includegraphics[width=\linewidth,height=.40\textheight,keepaspectratio]{\detokenize{assets/module-01/figure-071-revised.pdf}}
\par\smallskip{\small 图示：变化磁通产生电场环量，传导电流与变化电通量共同产生磁场环量}
\end{center}

\par\addvspace{.4\baselineskip}
\noindent\begin{minipage}{\linewidth}

\textbf{例 M1-43：轴向变化磁场产生的环形电场}

某轴对称时变磁场沿 \(z\) 轴，考察的半径为 \(\rho>0\)
的圆周所围圆盘内，磁场可近似写为
\(\mathbf B=B_z(t)\mathbf e_z\)，并在该圆盘上均匀。圆周的正向规定为从
\(+z\) 看逆时针。利用法拉第电磁感应定律求圆周上的方位电场分量
\(E_{\phi}(\rho,t)\)，分别说明 \(dB_z/dt\) 为正或为负时的绕行方向。


\end{minipage}\par

\par\noindent\begin{minipage}{\linewidth}
\textbf{解答}

\begin{center}
\includegraphics[width=\linewidth,height=.40\textheight,keepaspectratio]{\detokenize{assets/module-01/examples/M1-43.pdf}}
\par\smallskip{\small 图 M1-43：轴向变化磁场产生的环形电场}
\end{center}
\end{minipage}\par

图
M1-43中的圆盘是积分面而非磁场的物理边界，外侧磁场未示出。取纸外为正法向、逆时针为正边界后，左边是
\(2\pi\rho E_{\phi}\)，右边是
\(-\pi\rho^2dB_z/dt\)。图示磁通量增加时电场顺时针，正对应最后系数的负号；若磁通量减少，电场方向随之反转。

若半径 \(\rho\) 的圆周所围区域内，轴向磁场近似均匀，则

\[
2\pi\rho E_{\phi}
=
-\pi\rho^2\frac{dB_z}{dt},
\]

所以

\[
E_{\phi}=-\frac{\rho}{2}\frac{dB_z}{dt}.
\]

方向由 \(dB_z/dt\) 的符号和楞次定律共同决定，不能只凭图中箭头猜测。

\subsubsection{3.
界面上的切向跳跃}\label{m1:ux754cux9762ux4e0aux7684ux5207ux5411ux8df3ux8dc3}

设光滑界面 \(S\) 的单位法向 \(\hat{\mathbf n}\)
从负侧指向正侧，场的单侧极限为
\(\mathbf A_-,\mathbf A_+\)。取跨界窄矩形回路并令其高度趋于零。若旋度方程右边含面分布
\(\mathbf K_A\delta_S\)，则

\[
\boxed{
\hat{\mathbf n}\times
(\mathbf A_+-\mathbf A_-)
=
\mathbf K_A
}.
\]

在分布意义下，

\[
\nabla\times\mathbf A
=
(\nabla\times\mathbf A)_{\mathrm{bulk}}
+
\hat{\mathbf n}\times
(\mathbf A_+-\mathbf A_-)\delta_S.
\]

其中 \(\delta_S\)
是面狄拉克分布。若不存在旋度型面源，则场的切向分量连续；法向分量是否连续，应由散度方程和跨界薄柱判断。

对常见电磁边界，在场量有限且没有冲击型磁通时，

\[
\boxed{
\hat{\mathbf n}\times(\mathbf E_+-\mathbf E_-)=\mathbf0
},
\]

\[
\boxed{
\hat{\mathbf n}\times(\mathbf H_+-\mathbf H_-)
=\mathbf K_{\mathrm f}
},
\]

其中 \(\mathbf K_{\mathrm f}\) 为自由面电流密度。

\par\addvspace{.4\baselineskip}
\noindent\begin{minipage}{\linewidth}

\textbf{例 M1-44：面电流与磁场的切向跳跃}

平面 \(z=0\) 上分布自由面电流 \(\mathbf K_f=K\mathbf e_y\)，\(K\)
为常量；令 \(\mathbf H_{\pm}\) 表示从 \(z\gtrless0\)
逼近界面的场，统一法向取 \(\hat{\mathbf n}=\mathbf e_z\)。已知
\(\mathbf H_-=0\)。构造一个纯切向的 \(\mathbf H_+\)
使其满足磁场切向跃变关系，并直接核对叉乘的次序和方向；不要求由切向条件确定未给出的法向分量。


\end{minipage}\par

\par\noindent\begin{minipage}{\linewidth}
\textbf{解答}

\begin{center}
\includegraphics[width=\linewidth,height=.40\textheight,keepaspectratio]{\detokenize{assets/module-01/examples/M1-44.pdf}}
\par\smallskip{\small 图 M1-44：面电流与磁场的切向跳跃}
\end{center}
\end{minipage}\par

图 M1-44中 \(x\) 向右、\(z\) 向上，因此右手坐标的 \(+y\)
指向纸内，面电流用 \(\otimes\)
标出。上侧磁场向右、下侧为零，故法向叉乘跳跃量为
\(\mathbf e_z\times K\mathbf e_x=K\mathbf e_y\)。这里必须同时核对空间坐标的手性与边界公式，不能让跳跃箭头与面电流箭头平行。

取平面界面 \(z=0\)，令 \(\hat{\mathbf n}=\mathbf e_z\)，自由面电流为

\[
\mathbf K_{\mathrm f}=K\mathbf e_y.
\]

若 \(\mathbf H_-=\mathbf0\)，则可取

\[
\mathbf H_+=K\mathbf e_x,
\]

因为

\[
\mathbf e_z\times
(K\mathbf e_x-\mathbf0)
=K\mathbf e_y
=\mathbf K_{\mathrm f}.
\]

这个例子也说明，\(\mathbf H\) 的跳跃方向一般不与面电流同向，而是由
\(\hat{\mathbf n}\times(\mathbf H_+-\mathbf H_-)\) 决定。

\subsection{八、常见误区与使用条件}\label{m1:ux516bux5e38ux89c1ux8befux533aux4e0eux4f7fux7528ux6761ux4ef6}

\begin{enumerate}
\def\labelenumi{\arabic{enumi}.}
\tightlist
\item
  \textbf{旋度方向不是沿回路切线}：它沿有向小面的法向，回路只负责测量该法向分量。
\item
  \textbf{场线弯曲不等于旋度非零}：必须计算局部闭合环量除以面积的极限。
\item
  \textbf{旋度为零不等于场为常量}：非均匀梯度场也可处处无旋。
\item
  \textbf{局部无旋不自动推出全局保守}：还需检查区域是否单连通以及场是否在整个区域有定义。
\item
  \textbf{斯托克斯曲面不能任意穿过奇线}：更换曲面前必须检查场的定义域和光滑性。
\item
  \textbf{法向反向时边界方向也要反向}：只改一边会造成符号错误。
\item
  \textbf{行列式只是记忆工具}：曲线坐标下必须加入尺度因子。
\item
  \textbf{时变电磁问题使用安培定律时不能漏掉位移电流}：完整方程是安培--麦克斯韦定律。
\item
  \textbf{普通旋度与集中旋度要分层判断}：先检查普通区域，再用小回路或分布形式读取奇线贡献。
\end{enumerate}

\subsection{例题}\label{m1:ux4f8bux9898-1}

\par\addvspace{.4\baselineskip}
\noindent\begin{minipage}{\linewidth}

\subsubsection{由面内环量判定法向旋度（M1-45）}\label{m1:ux7531ux9762ux5185ux73afux91cfux5224ux5b9aux6cd5ux5411ux65cbux5ea6m1-45}

在光滑矢量场 \(\mathbf A\) 中取单位法向为 \(\mathbf n\) 的微小平面区域
\(\Delta S\)，边界 \(\partial S\) 的正向由右手定则与 \(\mathbf n\)
配对。已知当区域缩小到观察点时，正向环量与面积之比趋向一个正的有限数。判断该点旋度在
\(\mathbf n\)
上的投影符号，并说明同时反转面法向与边界正向后哪些量改变符号。


\end{minipage}\par

\par\noindent\begin{minipage}{\linewidth}
\textbf{解答}

\begin{center}
\includegraphics[width=\linewidth,height=.40\textheight,keepaspectratio]{\detokenize{assets/module-01/examples/M1-45.pdf}}
\par\smallskip{\small 图 M1-45：环量与法向旋度的符号}
\end{center}
\end{minipage}\par

图 M1-45 的正视图用蓝色圆点圆圈符号表示指向观察者的法向
\(\mathbf n\)；红色边界沿逆时针绕行，二者满足右手定则。侧视图把曲面画成一条边，向上的蓝色箭头与该边垂直，直角标记明确给出
\(\mathbf n\perp S\)。正向环量大于零，就意味着旋度在蓝色法向上的投影为正。若把法向翻转，绕行方向也翻转，下面的分子和投影同时变号，符号约定仍保持一致。

取单位法向为 \(\mathbf n\)
的微小平面区域，边界正向按右手定则确定。若沿该正向的环量为正，则

\[
(\nabla\times\mathbf A)\cdot\mathbf n
=\lim_{\Delta S\to0}\frac{1}{\Delta S}
\oint_{\partial S}\mathbf A\cdot d\mathbf l>0.
\]

反转法向时边界正向也必须反转；只反转其中一个会制造错误负号。

\par\addvspace{.4\baselineskip}
\noindent\begin{minipage}{\linewidth}

\subsubsection{刚体旋转与势涡不能只看流线（M1-46）}\label{m1:ux521aux4f53ux65cbux8f6cux4e0eux52bfux6da1ux4e0dux80fdux53eaux770bux6d41ux7ebfm1-46}

在垂直于 \(z\) 轴的截面中比较两种轴对称速度场：刚体转动场
\(\mathbf v=\boldsymbol\Omega\times\mathbf r\)，其中
\(\boldsymbol\Omega=\Omega\mathbf e_z\)；势涡场
\(\mathbf v=\Gamma(2\pi r)^{-1}\mathbf e_{\phi}\)，其中 \(r>0\)
为距轴的柱半径，\(\Omega,\Gamma\)
为常量。求两场的普通旋度及绕轴圆周的环量，指出势涡轴线的特殊性，并解释为什么圆形流线不能单独判定局部旋度。


\end{minipage}\par

\par\noindent\begin{minipage}{\linewidth}
\textbf{解答}

\begin{center}
\includegraphics[width=\linewidth,height=.40\textheight,keepaspectratio]{\detokenize{assets/module-01/examples/M1-46.pdf}}
\par\smallskip{\small 图 M1-46：刚体旋转与势涡的区别}
\end{center}
\end{minipage}\par

图 M1-46 示意取
\(\Omega>0,\Gamma>0\)，参数反号时相应场与涡量方向反转。图中的两幅图都有圆形流线，所以只看流线形状不能区分它们。左图速度随半径增加，局部小回路处处有正环量密度；右图速度按
\(1/r\)
减小，轴外局部环量密度为零，却有绕轴的有限总环量。计算时分别保留图中纸外的普通旋度与轴线上的集中涡量。

刚体旋转场 \(\mathbf v=\boldsymbol\Omega\times\mathbf r\) 满足

\[
\nabla\times\mathbf v=2\boldsymbol\Omega.
\]

对半径 \(R>0\)、从 \(+z\) 方向看逆时针为正向的圆周 \(C_R\)，图 M1-46
左侧刚体场给出

\[
\oint_{C_R}\mathbf v_{\rm rigid}\cdot d\mathbf l
=\int_0^{2\pi}\Omega R^2\,d\phi=2\pi\Omega R^2.
\]

势涡 \(\mathbf v=\Gamma(2\pi r)^{-1}\mathbf e_{\phi}\) 在 \(r>0\)
处普通旋度为零，而同一正向圆周的环量为

\[
\oint_{C_R}\mathbf v_{\rm vortex}\cdot d\mathbf l
=\int_0^{2\pi}\frac{\Gamma}{2\pi}\,d\phi=\Gamma.
\]

刚体场的环量随圆盘面积变化；势涡的绕轴环量与半径无关，其涡量集中在被排除的轴线。若回路不包围轴线且可在轴外缩回，势涡环量为零；把回路正向反转时，两种环量均反号。

\subsection{九、练习与解答}\label{m1:ux4e5dux7ec3ux4e60ux4e0eux89e3ux7b54}

\par\addvspace{.4\baselineskip}
\noindent\begin{minipage}{\linewidth}

\subsubsection{作业一：刚体定轴转动的速度场旋度}\label{m1:ux4f5cux4e1aux4e00ux521aux4f53ux5b9aux8f74ux8f6cux52a8ux7684ux901fux5ea6ux573aux65cbux5ea6}

\textbf{题号：M1-47}

\textbf{层级：} 基础　　\textbf{方法：} 旋度与环量、机电耦合

取右手单位正交基，\(\mathbf r\)
是质点相对于轴上原点的位置矢量；所有空间导数作用于
\(\mathbf r\)，而不作用于常角速度。

刚体以不随空间位置变化的角速度 \(\boldsymbol\omega\)
绕过原点的固定轴转动，速度场为

\[
\mathbf v=\boldsymbol\omega\times\mathbf r.
\]

证明

\[
\nabla\times\mathbf v=2\boldsymbol\omega,
\]

并解释系数 \(2\) 的物理意义。


\end{minipage}\par

\par\noindent\begin{minipage}{\linewidth}
\textbf{解答}

\begin{center}
\includegraphics[width=\linewidth,height=.40\textheight,keepaspectratio]{\detokenize{assets/module-01/examples/M1-47.pdf}}
\par\smallskip{\small 图 M1-47：刚体转动速度场的旋度}
\end{center}
\end{minipage}\par

图 M1-47中的速度箭头沿半径 \(\rho\) 的圆周切向，大小为
\(\omega\rho\)。若将整圆盘作为曲面，总环量是
\(2\pi\omega\rho^2\)，除以面积 \(\pi\rho^2\) 就读出法向旋度
\(2\omega\)；随后把转轴方向恢复为
\(\boldsymbol\omega\)。散度则来自法向通量，不能由这张环量图直接替代。

写成直角坐标分量：

\[
\boldsymbol\omega
=\omega_x\mathbf e_x+\omega_y\mathbf e_y+\omega_z\mathbf e_z,
\qquad
\mathbf r=x\mathbf e_x+y\mathbf e_y+z\mathbf e_z,
\]

\[
\begin{aligned}
v_x&=\omega_yz-\omega_zy,\\
v_y&=\omega_zx-\omega_xz,\\
v_z&=\omega_xy-\omega_yx.
\end{aligned}
\]

于是

\[
\begin{aligned}
(\nabla\times\mathbf v)_x
&=\frac{\partial v_z}{\partial y}
-\frac{\partial v_y}{\partial z}
=\omega_x-(-\omega_x)=2\omega_x,\\
(\nabla\times\mathbf v)_y
&=\frac{\partial v_x}{\partial z}
-\frac{\partial v_z}{\partial x}
=\omega_y-(-\omega_y)=2\omega_y,\\
(\nabla\times\mathbf v)_z
&=\frac{\partial v_y}{\partial x}
-\frac{\partial v_x}{\partial y}
=\omega_z-(-\omega_z)=2\omega_z.
\end{aligned}
\]

因此

\[
\boxed{
\nabla\times\mathbf v
=2\omega_x\mathbf e_x
+2\omega_y\mathbf e_y
+2\omega_z\mathbf e_z
=2\boldsymbol\omega
}.
\]

速度梯度的反对称部分对应局部刚性转动，而局部角速度定义为

\[
\boldsymbol\omega_{\text{局部}}
=\frac12\nabla\times\mathbf v.
\]

对刚体转动，\(\boldsymbol\omega_{\text{局部}}=\boldsymbol\omega\)，所以旋度恰为刚体角速度的两倍。

\par\addvspace{.4\baselineskip}
\noindent\begin{minipage}{\linewidth}

\subsubsection{作业二：普通旋度、环量与轴线集中涡量}\label{m1:ux4f5cux4e1aux4e8cux666eux901aux65cbux5ea6ux73afux91cfux4e0eux8f74ux7ebfux96c6ux4e2dux6da1ux91cf}

\textbf{题号：M1-48}

\textbf{层级：} 基础　　\textbf{方法：} 旋度与环量

圆柱坐标中的 \(\rho\) 为距轴的距离；圆周位于垂直于 \(z\) 轴的平面，从
\(+z\) 看逆时针为正向，配对圆盘法向为 \(+\mathbf e_z\)。

在 \(\rho>0\) 区域给定

\[
\mathbf A
=
\left(
a\rho+\frac{\kappa}{2\pi\rho}
\right)\mathbf e_{\phi},
\]

其中 \(a,\kappa\) 为常数。

\begin{enumerate}
\def\labelenumi{\arabic{enumi}.}
\tightlist
\item
  求 \(\rho>0\) 区域的普通旋度；
\item
  求半径为 \(R\)、绕 \(z\) 轴一周的正向圆周环量；
\item
  写出全空间的分布旋度，并与斯托克斯定理核对。
\end{enumerate}


\end{minipage}\par

\par\noindent\begin{minipage}{\linewidth}
\textbf{解答}

\begin{center}
\includegraphics[width=\linewidth,height=.40\textheight,keepaspectratio]{\detokenize{assets/module-01/examples/M1-48.pdf}}
\par\smallskip{\small 图 M1-48：轴线奇点的集中涡量}
\end{center}
\end{minipage}\par

图 M1-48中的白色小孔排除了轴线奇点，环形区域内可以先使用普通旋度
\(2a\mathbf e_z\)。内边界对环形区域取顺时针，和外边界的逆时针方向相反；这个符号使小孔上的有限环量不能被遗漏。最后把轴线补回时，剩余的
\(\kappa\) 正是集中涡量的系数。

柱坐标下只有方位分量，且与 \(\phi,z\) 无关，因此

\[
\nabla\times\mathbf A
=
\frac1\rho
\frac{d}{d\rho}
\left(
\rho A_{\phi}
\right)\mathbf e_z.
\]

代入

\[
\rho A_{\phi}
=a\rho^2+\frac{\kappa}{2\pi},
\]

得到

\[
\boxed{
\nabla\times\mathbf A
=2a\,\mathbf e_z
},
\qquad \rho>0.
\]

圆周取从 \(+z\) 方向看逆时针为正向。圆周上
\(d\mathbf l=R\,d\phi\,\mathbf e_{\phi}\)，故

\[
\begin{aligned}
\Gamma_R
&=\int_0^{2\pi}
\left(
aR+\frac{\kappa}{2\pi R}
\right)R\,d\phi\\
&=2\pi aR^2+\kappa.
\end{aligned}
\]

普通旋度 \(2a\mathbf e_z\) 在圆盘上的通量只有

\[
\int_{S_R}2a\,dS=2\pi aR^2,
\]

比环量少 \(\kappa\)。缺少的部分集中在轴线上，因此全空间分布形式为

\[
\boxed{
\nabla\times\mathbf A
=
2a\,\mathbf e_z
+
\kappa\,\delta(x)\delta(y)\mathbf e_z
}.
\]

对圆盘积分后恰得

\[
2\pi aR^2+\kappa=\Gamma_R,
\]

与斯托克斯定理一致。若只在穿孔圆盘 \(\varepsilon<\rho<R\)
上使用普通旋度，正法向取
\(+\mathbf e_z\)，则外圆周为逆时针，内圆周为顺时针。内边界环量为

\[
\Gamma_{\mathrm{in}}
=-\left(2\pi a\varepsilon^2+\kappa\right),
\]

而环形区域上的普通旋度通量为

\[
\int_{\varepsilon<\rho<R}2a\,dS
=2\pi a\left(R^2-\varepsilon^2\right).
\]

普通斯托克斯定理给出

\[
\Gamma_R+\Gamma_{\mathrm{in}}
=2\pi a\left(R^2-\varepsilon^2\right),
\]

故

\[
\Gamma_R=2\pi aR^2+\kappa.
\]

这表明穿孔计算并没有丢失集中涡量：内边界按顺时针取向产生
\(-\kappa\)，移到等式另一边后正好恢复外圆周环量中的 \(+\kappa\)。

\par\addvspace{.4\baselineskip}
\noindent\begin{minipage}{\linewidth}

\subsubsection{作业三：有限半径时变磁场产生的感生电场}\label{m1:ux4f5cux4e1aux4e09ux6709ux9650ux534aux5f84ux65f6ux53d8ux78c1ux573aux4ea7ux751fux7684ux611fux751fux7535ux573a}

\textbf{题号：M1-49}

\textbf{层级：} 基础　　\textbf{方法：} 旋度与环量

这里 \(R>0\) 为磁场区域半径，\(\rho\) 为距同一轴线的距离，\(B(t)\)
可正可负；\(\mathbf e_{\phi}\) 取从 \(+z\)
看逆时针方向。只求具有所给轴对称性的感生方位电场。

无限长圆柱区域内存在轴向磁场

\[
\mathbf B(\rho,t)
=
\begin{cases}
B(t)\mathbf e_z,&0\le \rho<R,\\
\mathbf0,&\rho>R.
\end{cases}
\]

假设系统具有轴对称性，求感生电场
\(\mathbf E=E_{\phi}(\rho,t)\mathbf e_{\phi}\)，并完成以下检查：

\begin{enumerate}
\def\labelenumi{\arabic{enumi}.}
\tightlist
\item
  分别求 \(\rho<R\) 与 \(\rho>R\) 的 \(E_{\phi}\)；
\item
  在两个普通区域内核对法拉第定律的微分形式；
\item
  检查 \(\rho=R\) 处的切向连续性。
\end{enumerate}


\end{minipage}\par

\par\noindent\begin{minipage}{\linewidth}
\textbf{解答}

\begin{center}
\includegraphics[width=\linewidth,height=.40\textheight,keepaspectratio]{\detokenize{assets/module-01/examples/M1-49.pdf}}
\par\smallskip{\small 图 M1-49：有限半径磁场的感生电场}
\end{center}
\end{minipage}\par

图 M1-49左图的半径 \(R\) 是磁场分布的物理边界，求半径 \(\rho\)
的环量时要先判断积分圆位于内区还是外区。内区磁通面积为
\(\pi\rho^2\)，外区只保留 \(\pi R^2\)；两者再除以相同的周长
\(2\pi\rho\)，分别给出正比于 \(\rho\) 和反比于 \(\rho\)
的两段曲线。右图在 \(R\) 处连续，正是两段解必须通过的检验。

取半径为 \(\rho\) 的圆周，从 \(+z\) 方向看逆时针为正向。法拉第定律为

\[
2\pi\rho E_{\phi}
=
-\frac{d\Phi_B}{dt}.
\]

当 \(\rho<R\) 时，圆盘全部处于磁场区，

\[
\Phi_B=\pi\rho^2B(t),
\]

所以

\[
\boxed{
E_{\phi}(\rho,t)
=
-\frac{\rho}{2}\dot B(t)
},
\qquad \rho<R.
\]

当 \(\rho>R\) 时，只有半径 \(R\) 内有磁通，

\[
\Phi_B=\pi R^2B(t),
\]

所以

\[
\boxed{
E_{\phi}(\rho,t)
=
-\frac{R^2}{2\rho}\dot B(t)
},
\qquad \rho>R.
\]

在轴对称情形，

\[
(\nabla\times\mathbf E)_z
=
\frac1\rho\frac{\partial(\rho E_{\phi})}{\partial\rho}.
\]

内区有

\[
(\nabla\times\mathbf E)_z
=
\frac1\rho\frac{\partial}{\partial\rho}
\left(
-\frac{\rho^2}{2}\dot B
\right)
=-\dot B,
\]

正好等于 \(-\partial B_z/\partial t\)。外区有

\[
(\nabla\times\mathbf E)_z
=
\frac1\rho\frac{\partial}{\partial\rho}
\left(
-\frac{R^2}{2}\dot B
\right)
=0,
\]

也与外区 \(\mathbf B=\mathbf0\) 相符。

在界面两侧取极限：

\[
E_{\phi}(R^-)
=-\frac R2\dot B,
\qquad
E_{\phi}(R^+)
=-\frac{R^2}{2R}\dot B
=-\frac R2\dot B.
\]

因此

\[
\boxed{
E_{\phi}(R^-)=E_{\phi}(R^+)
}.
\]

切向电场连续，没有额外的面狄拉克旋度。虽然理想化磁场本身在 \(\rho=R\)
跳跃，但其时间导数只具有阶跃而不含空间狄拉克项，所得电场仍满足上述边界条件。

\section{专题六：矢量积分定理：线积分、通量、高斯定理与斯托克斯定理（基础必讲）}\label{m1:ux4e13ux9898ux516dux77e2ux91cfux79efux5206ux5b9aux7406ux7ebfux79efux5206ux901aux91cfux9ad8ux65afux5b9aux7406ux4e0eux65afux6258ux514bux65afux5b9aux7406ux57faux7840ux5fc5ux8bb2}

\begin{quote}
阅读准备：积分对象必须带方向 线积分中的 \(d\mathbf l\)
沿路径正向，通量积分中的 \(d\mathbf S=\mathbf n\,dS\)
沿选定法向。开曲面的法向与边界绕向由右手定则成对确定。``单连通''可先理解为区域内任意闭合曲线都能连续缩成一点，过程中不跨越孔洞或奇线。
\end{quote}

\par\addvspace{.4\baselineskip}
\noindent\begin{minipage}{\linewidth}

\textbf{例 M1-50：斯托克斯取向的同时反转}

设光滑开曲面 \(S\) 的单位法向为 \(\mathbf n\)，边界 \(C=\partial S\)
的方向与它按右手定则配对；矢量场 \(\mathbf A\) 在包含 \(S\)
的邻域内一阶连续可微。保持曲面的几何位置不变，只将法向改为
\(-\mathbf n\)。说明斯托克斯公式中的边界正向、曲面通量及闭合环量各怎样改变，判断仅反转一方是否仍保持原等式。


\end{minipage}\par

\par\noindent\begin{minipage}{\linewidth}
\textbf{解答}

\begin{center}
\includegraphics[width=\linewidth,height=.40\textheight,keepaspectratio]{\detokenize{assets/module-01/examples/M1-50.pdf}}
\par\smallskip{\small 图 M1-50：斯托克斯取向的同时反转}
\end{center}
\end{minipage}\par

图
M1-50左图的纸外法向配逆时针边界，右图的纸内法向配顺时针边界。右图不仅改了曲面方向，也将沿边界的所有切向微元反号，因此环量与旋度通量同时反号。将两幅取向成对比较，斯托克斯等式保持成立。

同一开曲面若把法向反向，边界正向也必须反向，因而斯托克斯定理两边同时变号，等式不变。

梯度、散度和旋度描述场在一点附近的局部变化；积分定理则把这些局部量累积成端点差、闭合面通量或边界环量。本专题讨论积分对象的定向、定理成立的条件、积分式局部化为微分式的条件，以及奇点和界面带来的额外边界项。

\subsection{一、三类积分对象与方向约定}\label{m1:ux4e00ux4e09ux7c7bux79efux5206ux5bf9ux8c61ux4e0eux65b9ux5411ux7ea6ux5b9a}

\subsubsection{1.
有向曲线与线积分}\label{m1:ux6709ux5411ux66f2ux7ebfux4e0eux7ebfux79efux5206}

设分片光滑曲线 \(C\) 的参数方程为

\[
\mathbf r=\mathbf r(t),
\qquad
a\le t\le b.
\]

参数增大的方向定义曲线正向，线元为

\[
d\mathbf l=d\mathbf r
=\mathbf r'(t)\,dt.
\]

矢量场 \(\mathbf A\) 沿 \(C\) 的线积分定义为

\[
\boxed{
\int_C\mathbf A\cdot d\mathbf l
=
\int_a^b
\mathbf A(\mathbf r(t))\cdot\mathbf r'(t)\,dt
}.
\]

它只累加 \(\mathbf A\)
的切向分量。保持方向的光滑重新参数化不改变积分值；反转曲线方向则使积分变号：

\[
\int_{-C}\mathbf A\cdot d\mathbf l
=
-\int_C\mathbf A\cdot d\mathbf l.
\]

闭合曲线的线积分写作

\[
\oint_C\mathbf A\cdot d\mathbf l,
\]

称为环量。

\subsubsection{2.
有向曲面与通量积分}\label{m1:ux6709ux5411ux66f2ux9762ux4e0eux901aux91cfux79efux5206}

设正则曲面片 \(S\) 参数化为

\[
\mathbf r=\mathbf r(u,v),
\qquad
(u,v)\in D,
\]

其中 \(\mathbf r_u\times\mathbf r_v\neq\mathbf0\)。参数次序 \((u,v)\)
选定有向面积元

\[
d\mathbf S
=
(\mathbf r_u\times\mathbf r_v)\,du\,dv
=
\hat{\mathbf n}\,dS.
\]

矢量场穿过曲面的通量定义为

\[
\boxed{
\int_S\mathbf A\cdot d\mathbf S
=
\iint_D
\mathbf A(\mathbf r(u,v))
\cdot
(\mathbf r_u\times\mathbf r_v)\,du\,dv
}.
\]

它只累加 \(\mathbf A\)
的法向分量。交换参数次序或反转法向会使通量变号。若 \(S=\partial V\)
是闭合区域的边界，则统一取外法向。

\subsubsection{3.
开曲面的边界方向}\label{m1:ux5f00ux66f2ux9762ux7684ux8fb9ux754cux65b9ux5411}

对可定向开放曲面 \(S\)，选定法向 \(\hat{\mathbf n}\) 后，边界
\(\partial S\)
的正向由右手定则确定：右手四指沿边界正向弯曲时，大拇指指向
\(\hat{\mathbf n}\)。法向反向时，边界正向必须同时反向。

``可定向''不是多余条件：只有能在整个曲面上连续选取法向，才能定义全局有向通量。莫比乌斯带不存在这样的连续法向，因此不能直接对整个莫比乌斯带使用上述通量积分或斯托克斯定理。

\begin{center}
\includegraphics[width=\linewidth,height=.40\textheight,keepaspectratio]{\detokenize{assets/module-01/tikz-m1-topic6-v4-01-integral-objects-orientation.png}}
\par\smallskip{\small 图示：线积分读取切向分量；曲面侧视图明确显示法向垂直于曲面，通量只读取法向分量；开曲面法向与边界方向成对绑定}
\end{center}

\par\addvspace{.4\baselineskip}
\noindent\begin{minipage}{\linewidth}

\textbf{例 M1-51：常矢量场的路径积分}

空间中给定常矢量场 \(\mathbf A=\mathbf A_0\)。设 \(P,Q\)
的位置矢量分别为 \(\mathbf r_P,\mathbf r_Q\)，\(L\) 为从 \(P\) 到 \(Q\)
的任意分段光滑有向路径。求
\(\int_L\mathbf A_0\cdot d\mathbf l\)，比较直线与弯曲路径的结果，并说明反向从
\(Q\) 移到 \(P\) 时积分如何变化。


\end{minipage}\par

\par\noindent\begin{minipage}{\linewidth}
\textbf{解答}

\begin{center}
\includegraphics[width=\linewidth,height=.40\textheight,keepaspectratio]{\detokenize{assets/module-01/examples/M1-51.pdf}}
\par\smallskip{\small 图 M1-51：常矢量场的路径积分}
\end{center}
\end{minipage}\par

图 M1-51中绿色、橙色曲线的局部切向不同，但从 \(P\) 到 \(Q\)
的位移总和相同。常矢量 \(\mathbf A_0\) 可以移到积分外，积分只累加各段
\(d\mathbf l\)，得到图中的直线位移。将它与 \(\mathbf A_0\)
点乘即可，不需要逐段参数化弯曲路径。

若 \(\mathbf A=\mathbf A_0\) 为常矢量，则任意从 \(P\) 到 \(Q\)
的曲线都有

\[
\int_P^Q\mathbf A_0\cdot d\mathbf l
=
\mathbf A_0\cdot(\mathbf r_Q-\mathbf r_P).
\]

路径形状可以不同，但端点位移相同；若曲线闭合，结果为零。

\subsection{二、三条基本积分定理}\label{m1:ux4e8cux4e09ux6761ux57faux672cux79efux5206ux5b9aux7406}

\subsubsection{1. 梯度定理}\label{m1:ux68afux5ea6ux5b9aux7406}

若标量场 \(\phi\) 在包含曲线 \(C\) 的开邻域内属于 \(C^1\)，且 \(C\) 从
\(P\) 指向 \(Q\)，则

\[
\boxed{
\int_C\nabla\phi\cdot d\mathbf l
=
\phi(Q)-\phi(P)
}.
\]

它把一维曲线上的梯度积分化为零维边界，即两个端点的函数值之差。特别地，

\[
\oint_C\nabla\phi\cdot d\mathbf l=0.
\]

\par\addvspace{.4\baselineskip}
\noindent\begin{minipage}{\linewidth}

\textbf{例 M1-52：梯度场积分的路径无关性}

平面上给定 \(\phi(x,y)=x^2+y^2\) 及其梯度场 \(\nabla\phi\)。两固定端点为
\(P=(0,0)\)、\(Q=(a,b)\)，其中 \(a,b\) 为常量。计算沿任意分段光滑路径从
\(P\) 到 \(Q\) 的
\(\int_P^Q\nabla\phi\cdot d\mathbf l\)，并说明结果依赖端点还是路径的具体形状。


\end{minipage}\par

\par\noindent\begin{minipage}{\linewidth}
\textbf{解答}

\begin{center}
\includegraphics[width=\linewidth,height=.40\textheight,keepaspectratio]{\detokenize{assets/module-01/examples/M1-52.pdf}}
\par\smallskip{\small 图 M1-52：梯度场积分的路径无关性}
\end{center}
\end{minipage}\par

图
M1-52中的三条路径具有相同的起终点与积分方向。因为被积场是同一个单值势函数的梯度，各段切向变化在积分后只留下
\(\phi(Q)-\phi(P)\)。用图中 \(P=(0,0),Q=(a,b)\) 代入端点值，便得到
\(a^2+b^2\)。

取 \(\phi=x^2+y^2\)。无论从 \(P=(0,0)\) 到 \(Q=(a,b)\)
的分片光滑路径怎样弯曲，都有

\[
\int_P^Q\nabla\phi\cdot d\mathbf l
=
\phi(Q)-\phi(P)
=a^2+b^2.
\]

\subsubsection{2.
高斯散度定理}\label{m1:ux9ad8ux65afux6563ux5ea6ux5b9aux7406}

若有界区域 \(V\) 的边界 \(\partial V\) 分片光滑，且 \(\mathbf A\)
在包含闭包 \(\overline V\) 的邻域内属于 \(C^1\)，则

\[
\boxed{
\oint_{\partial V}\mathbf A\cdot d\mathbf S
=
\int_V\nabla\cdot\mathbf A\,dV
}.
\]

闭合曲面取外法向。它把三维体积中的散度积分化为二维边界通量。

\par\addvspace{.4\baselineskip}
\noindent\begin{minipage}{\linewidth}

\textbf{例 M1-53：径向场的高斯定理核对}

在三维空间中取 \(\mathbf A(\mathbf r)=\mathbf r\)。设 \(V\)
为任意有界、边界分片光滑的体积，闭合面 \(S=\partial V\)
一律取体积外法向。分别求 \(\nabla\cdot\mathbf A\) 和
\(\oint_S\mathbf A\cdot d\mathbf S\)，用高斯定理说明后者与 \(V\)
的体积有什么关系，并用同心球面作核对。


\end{minipage}\par

\par\noindent\begin{minipage}{\linewidth}
\textbf{解答}

\begin{center}
\includegraphics[width=\linewidth,height=.40\textheight,keepaspectratio]{\detokenize{assets/module-01/examples/M1-53.pdf}}
\par\smallskip{\small 图 M1-53：径向场的高斯定理核对}
\end{center}
\end{minipage}\par

图 M1-53以球面作为高斯面的具体检验，外法向与 \(\mathbf r\)
平行，所以面上被积量就是 \(R\)。球面法给 \(R\times4\pi R^2\)，体积分法给
\(3\times4\pi R^3/3\)，两者相同。对一般闭区域，下述
\(\nabla\cdot\mathbf r=3\) 直接把这条图示规律推广为通量 \(3V\)。

取 \(\mathbf A=\mathbf r=x\mathbf e_x+y\mathbf e_y+z\mathbf e_z\)，
\(S\) 为半径 \(R\) 的球面。由于 \(\nabla\cdot\mathbf A=3\)，

\[
\oint_S\mathbf A\cdot d\mathbf S
=
3\cdot\frac{4\pi R^3}{3}
=4\pi R^3.
\]

直接计算时球面上 \(\mathbf A\cdot\hat{\mathbf n}=R\)，乘以面积
\(4\pi R^2\) 也得到同一结果。

\subsubsection{3.
斯托克斯定理}\label{m1:ux65afux6258ux514bux65afux5b9aux7406}

若 \(S\) 是分片光滑、可定向的曲面，边界 \(\partial S\) 分片光滑，且
\(\mathbf A\) 在包含 \(S\) 的邻域内属于 \(C^1\)，则

\[
\boxed{
\oint_{\partial S}\mathbf A\cdot d\mathbf l
=
\int_S(\nabla\times\mathbf A)\cdot d\mathbf S
}.
\]

边界正向与曲面法向满足右手定则。它把二维曲面上的旋度通量化为一维边界环量。

\par\addvspace{.4\baselineskip}
\noindent\begin{minipage}{\linewidth}

\textbf{例 M1-54：线性场的斯托克斯核对}

给定光滑场 \(\mathbf A=x\mathbf e_y\)。圆盘 \(S\) 位于 \(xy\)
平面，圆心为原点、半径为 \(R>0\)，法向为 \(+\mathbf e_z\)；边界 \(C\) 从
\(+z\) 看按逆时针方向绕行。先求旋度，再使用斯托克斯定理计算
\(\oint_C\mathbf A\cdot d\mathbf l\)，写清取向进入结果符号的方式。


\end{minipage}\par

\par\noindent\begin{minipage}{\linewidth}
\textbf{解答}

\begin{center}
\includegraphics[width=\linewidth,height=.40\textheight,keepaspectratio]{\detokenize{assets/module-01/examples/M1-54.pdf}}
\par\smallskip{\small 图 M1-54：线性场的斯托克斯核对}
\end{center}
\end{minipage}\par

图 M1-54中的绿色箭头在左半盘向下、右半盘向上，正对应
\(\mathbf A=x\mathbf e_y\)。固定逆时针边界后，法向取纸外，计算得到的旋度投影为
\(1\)，所以斯托克斯积分就是图中圆盘的面积 \(\pi R^2\)。圆心上方的
\(\odot\) 只标记法向，不是额外的点源。

取 \(\mathbf A=x\mathbf e_y\)，\(C\) 为 \(xy\) 平面内半径 \(R\)
的逆时针圆周，\(S\) 取法向 \(+\mathbf e_z\) 的圆盘。因为
\(\nabla\times\mathbf A=\mathbf e_z\)，所以

\[
\oint_C\mathbf A\cdot d\mathbf l
=
\int_S\mathbf e_z\cdot d\mathbf S
=\pi R^2.
\]

这三条定理都具有同一结构：

\[
\boxed{
\text{区域内部的导数积分}
=
\text{区域边界上的原量积分}
}.
\]

梯度定理中的边界是两个端点；高斯定理中的边界是闭合曲面；斯托克斯定理中的边界是开放曲面的闭合边界曲线。

\subsection{三、三条定理的证明机制}\label{m1:ux4e09ux4e09ux6761ux5b9aux7406ux7684ux8bc1ux660eux673aux5236}

\subsubsection{1.
梯度定理：链式法则与一元微积分基本定理}\label{m1:ux68afux5ea6ux5b9aux7406ux94feux5f0fux6cd5ux5219ux4e0eux4e00ux5143ux5faeux79efux5206ux57faux672cux5b9aux7406}

沿曲线 \(\mathbf r(t)\)，

\[
\frac d{dt}\phi(\mathbf r(t))
=
\nabla\phi(\mathbf r(t))\cdot\mathbf r'(t).
\]

两边从 \(a\) 积分到 \(b\)：

\[
\begin{aligned}
\int_C\nabla\phi\cdot d\mathbf l
&=
\int_a^b
\nabla\phi(\mathbf r(t))\cdot\mathbf r'(t)\,dt\\
&=
\int_a^b
\frac d{dt}\phi(\mathbf r(t))\,dt\\
&=
\phi(\mathbf r(b))-\phi(\mathbf r(a)).
\end{aligned}
\]

证明只使用链式法则和一元微积分基本定理，因此路径可以弯曲，结果仍只由端点决定。

\subsubsection{2.
高斯定理：一维基本定理加内部面抵消}\label{m1:ux9ad8ux65afux5b9aux7406ux4e00ux7ef4ux57faux672cux5b9aux7406ux52a0ux5185ux90e8ux9762ux62b5ux6d88}

先取直角长方体

\[
V=[x_0,x_1]\times[y_0,y_1]\times[z_0,z_1].
\]

法向为 \(\pm\mathbf e_x\) 的两面总通量可精确写成

\[
\begin{aligned}
\Phi_x
&=
\int_{y_0}^{y_1}\int_{z_0}^{z_1}
\left[
A_x(x_1,y,z)-A_x(x_0,y,z)
\right]\,dz\,dy\\
&=
\int_{y_0}^{y_1}\int_{z_0}^{z_1}
\int_{x_0}^{x_1}
\frac{\partial A_x}{\partial x}\,dx\,dz\,dy\\
&=
\int_V\frac{\partial A_x}{\partial x}\,dV.
\end{aligned}
\]

同理，

\[
\Phi_y=\int_V\frac{\partial A_y}{\partial y}\,dV,
\qquad
\Phi_z=\int_V\frac{\partial A_z}{\partial z}\,dV.
\]

相加便得到长方体上的高斯定理，而且没有把非均匀场近似为面上常量。

把一般区域剖分成许多小体积时，相邻小体积在公共面上使用相反外法向，内部通量两两抵消，只留下整体外边界。对分片光滑边界取剖分极限，即得到一般形式。

\begin{center}
\includegraphics[width=\linewidth,height=.40\textheight,keepaspectratio]{\detokenize{assets/module-01/tikz-m1-topic6-v4-03-gauss-exact-box-cancellation.jpg}}
\par\smallskip{\small 图示：高斯定理在小长方体上由三次一维基本定理得到，拼接时内部公共面通量抵消}
\end{center}

\subsubsection{3.
斯托克斯定理：参数平面上的格林公式与内部边抵消}\label{m1:ux65afux6258ux514bux65afux5b9aux7406ux53c2ux6570ux5e73ux9762ux4e0aux7684ux683cux6797ux516cux5f0fux4e0eux5185ux90e8ux8fb9ux62b5ux6d88}

取曲面片

\[
\mathbf r=\mathbf r(u,v),
\qquad
(u,v)\in D,
\]

并令

\[
a(u,v)=\mathbf A(\mathbf r)\cdot\mathbf r_u,
\qquad
b(u,v)=\mathbf A(\mathbf r)\cdot\mathbf r_v.
\]

边界线积分在参数平面中写成

\[
\oint_{\partial S}\mathbf A\cdot d\mathbf l
=
\oint_{\partial D}
\left(a\,du+b\,dv\right).
\]

二维格林公式给出

\[
\oint_{\partial D}
\left(a\,du+b\,dv\right)
=
\iint_D
\left(
\frac{\partial b}{\partial u}
-\frac{\partial a}{\partial v}
\right)\,du\,dv.
\]

链式法则和混合偏导抵消可得

\[
\frac{\partial b}{\partial u}
-\frac{\partial a}{\partial v}
=
(\nabla\times\mathbf A)
\cdot
(\mathbf r_u\times\mathbf r_v).
\]

所以单个参数曲面片满足斯托克斯定理。一般曲面可剖分成多个曲面片；相邻曲面片在公共边上的正向相反，内部线积分抵消，只剩整体边界。

\begin{center}
\includegraphics[width=\linewidth,height=.40\textheight,keepaspectratio]{\detokenize{assets/module-01/tikz-m1-topic6-v4-04-stokes-patch-cancellation.jpg}}
\par\smallskip{\small 图示：斯托克斯定理由参数平面的格林公式得到，拼接曲面片时内部公共边环量抵消}
\end{center}

二维格林公式本身可先在矩形上用两次一元基本定理证明，再通过平面区域剖分推广。因此三条基本积分定理最终都建立在``一维基本定理加内部边界抵消''之上。

\subsection{四、积分式与微分式之间的局部化}\label{m1:ux56dbux79efux5206ux5f0fux4e0eux5faeux5206ux5f0fux4e4bux95f4ux7684ux5c40ux90e8ux5316}

\subsubsection{1.
微分式推出积分式}\label{m1:ux5faeux5206ux5f0fux63a8ux51faux79efux5206ux5f0f}

若微分方程在区域中逐点成立，且场具有积分定理要求的光滑性，则直接积分并使用相应定理即可得到积分式。例如

\[
\nabla\cdot\mathbf A=f
\]

在 \(V\) 中成立时，

\[
\oint_{\partial V}\mathbf A\cdot d\mathbf S
=
\int_Vf\,dV.
\]

\subsubsection{2.
体积分的局部化}\label{m1:ux4f53ux79efux5206ux7684ux5c40ux90e8ux5316}

反过来，若对任意足够小且闭包包含在 \(\Omega\) 内的体积 \(V\) 都有

\[
\oint_{\partial V}\mathbf A\cdot d\mathbf S
=
\int_Vf\,dV,
\]

则高斯定理给出

\[
\int_V
\left(
\nabla\cdot\mathbf A-f
\right)dV=0.
\]

若被积函数连续，反证法可得

\[
\boxed{
\nabla\cdot\mathbf A=f
}
\]

逐点成立：若某点差值为正或负，连续性会保证其在该点附近保持同号，从而与任意小体积上的积分为零矛盾。

\subsubsection{3.
曲面积分的局部化}\label{m1:ux66f2ux9762ux79efux5206ux7684ux5c40ux90e8ux5316}

若对任意足够小的有向曲面片 \(S\) 都有

\[
\oint_{\partial S}\mathbf A\cdot d\mathbf l
=
\int_S\mathbf B\cdot d\mathbf S,
\]

则

\[
\int_S
\left(
\nabla\times\mathbf A-\mathbf B
\right)\cdot\hat{\mathbf n}\,dS=0.
\]

在连续条件下，可以把小面法向依次取为任意方向，得到

\[
\boxed{
\nabla\times\mathbf A=\mathbf B
}
\]

逐点成立。

``某一个球''\,``某一族同方向圆盘''或``某一个特殊高斯面''上的积分等式都不足以推出完整微分方程。关键是区域能够任意缩小并改变位置、形状或法向。

\subsubsection{4.
弱形式与分布形式（拔高选读）}\label{m1:ux5f31ux5f62ux5f0fux4e0eux5206ux5e03ux5f62ux5f0fux62d4ux9ad8ux9009ux8bfb}

若场只有局部可积性而不连续，上述结论通常只能理解为几乎处处成立，或在分布意义下成立。点源、线源和面源会表现为狄拉克分布；跳跃场的普通导数还会产生界面面分布。

\begin{center}
\includegraphics[width=\linewidth,height=.40\textheight,keepaspectratio]{\detokenize{assets/module-01/tikz-m1-topic6-v4-05-localization-regular-singular-interface.png}}
\par\smallskip{\small 图示：从积分式局部化前必须区分光滑普通点、集中奇点和跳跃界面}
\end{center}

因此从积分式转成微分式前，应依次检查：

\begin{enumerate}
\def\labelenumi{\arabic{enumi}.}
\tightlist
\item
  等式是否对任意足够小的区域成立；
\item
  场与源是否在该区域内连续或足够光滑；
\item
  区域是否穿过奇点、奇线或跳跃界面；
\item
  时间导数能否与固定区域上的积分交换；
\item
  若条件不满足，是否应保留积分式、挖去奇异集合或改用分布形式。
\end{enumerate}

\subsection{五、由积分定理导出的分部积分公式}\label{m1:ux4e94ux7531ux79efux5206ux5b9aux7406ux5bfcux51faux7684ux5206ux90e8ux79efux5206ux516cux5f0f}

以下设 \(V\) 为边界分片光滑的有界区域，矢量场在包含 \(\overline V\)
的邻域内具有所需的一阶连续偏导，标量场在需要使用拉普拉斯算符时具有二阶连续偏导。较弱的正则性条件下，以下公式仍可在弱意义中成立，但边界迹与法向导数需另行定义。

\subsubsection{1.
标量---矢量分部积分}\label{m1:ux6807ux91cfux77e2ux91cfux5206ux90e8ux79efux5206}

由乘积法则

\[
\nabla\cdot(f\mathbf A)
=
\nabla f\cdot\mathbf A
+
f\nabla\cdot\mathbf A
\]

和高斯定理，得到

\[
\boxed{
\int_V f\nabla\cdot\mathbf A\,dV
=
\oint_{\partial V}f\mathbf A\cdot d\mathbf S
-
\int_V\nabla f\cdot\mathbf A\,dV
}.
\]

这是多维分部积分的基本形式。取 \(\mathbf A=g\mathbf e_i\) 可得分量形式

\[
\boxed{
\int_V f\frac{\partial g}{\partial x_i}\,dV
=
\oint_{\partial V}fg\,n_i\,dS
-
\int_Vg\frac{\partial f}{\partial x_i}\,dV
}.
\]

\subsubsection{2.
旋度的分部积分}\label{m1:ux65cbux5ea6ux7684ux5206ux90e8ux79efux5206}

由恒等式

\[
\nabla\cdot(\mathbf A\times\mathbf B)
=
\mathbf B\cdot(\nabla\times\mathbf A)
-
\mathbf A\cdot(\nabla\times\mathbf B)
\]

和高斯定理，

\[
\boxed{
\int_V
\left[
\mathbf B\cdot(\nabla\times\mathbf A)
-
\mathbf A\cdot(\nabla\times\mathbf B)
\right]dV
=
\oint_{\partial V}
(\mathbf A\times\mathbf B)\cdot d\mathbf S
}.
\]

这个公式常用于电磁场算符的对称性、互易关系和边界项分析。边界项能否消失取决于具体边界条件，不能在推导中无理由删去。

\subsubsection{3.
格林第一恒等式}\label{m1:ux683cux6797ux7b2cux4e00ux6052ux7b49ux5f0f}

在标量---矢量分部积分中取 \(\mathbf A=\nabla\psi\)、\(f=\phi\)，得到

\[
\boxed{
\int_V
\left(
\phi\nabla^2\psi
+
\nabla\phi\cdot\nabla\psi
\right)dV
=
\oint_{\partial V}
\phi\frac{\partial\psi}{\partial n}\,dS
},
\]

其中

\[
\frac{\partial\psi}{\partial n}
=
\nabla\psi\cdot\hat{\mathbf n}.
\]

\subsubsection{4.
格林第二恒等式}\label{m1:ux683cux6797ux7b2cux4e8cux6052ux7b49ux5f0f}

交换 \(\phi,\psi\) 后相减，得到

\[
\boxed{
\int_V
\left(
\phi\nabla^2\psi
-
\psi\nabla^2\phi
\right)dV
=
\oint_{\partial V}
\left(
\phi\frac{\partial\psi}{\partial n}
-
\psi\frac{\partial\phi}{\partial n}
\right)dS
}.
\]

\begin{center}
\includegraphics[width=\linewidth,height=.40\textheight,keepaspectratio]{\detokenize{assets/module-01/figure-083-revised.pdf}}
\par\smallskip{\small 图示：高斯定理结合乘积法则依次产生多维分部积分、格林第一恒等式和格林第二恒等式}
\end{center}

\par\addvspace{.4\baselineskip}
\noindent\begin{minipage}{\linewidth}

\textbf{例 M1-55：零边界调和解的唯一性}

设 \(V\) 为有界连通区域，边界分片光滑；实函数 \(u\)
在闭包上连续并具有使用格林第一恒等式所需的光滑性。在区域内有
\(\nabla^2u=0\)，整个边界上有 \(u=0\)。利用格林第一恒等式证明 \(u\) 在
\(V\) 内恒为零，并指出边界条件和连通性各在推理中的作用。


\end{minipage}\par

\par\noindent\begin{minipage}{\linewidth}
\textbf{解答}

\begin{center}
\includegraphics[width=\linewidth,height=.40\textheight,keepaspectratio]{\detokenize{assets/module-01/examples/M1-55.pdf}}
\par\smallskip{\small 图 M1-55：零边界调和解的唯一性}
\end{center}
\end{minipage}\par

图 M1-55中的封闭边界规定
\(u=0\)，因此格林第一恒等式的边界项为零；区域内部又有
\(\nabla^2u=0\)，只剩非负的梯度平方积分。它为零迫使梯度处处为零。沿图中连接任意两点的区域内路径积分，函数值相同，再用边界上的零值确定该常数为零。

若连通区域内 \(\nabla^2u=0\) 且
\(u|_{\partial V}=0\)，在格林第一恒等式中令 \(\phi=\psi=u\)，则

\[
\int_V|\nabla u|^2\,dV=0.
\]

因此 \(\nabla u=\mathbf0\)，\(u\) 在连通区域内为常数；再由边界值为零得到
\(u\equiv0\)。这正是静电边值问题唯一性证明的核心步骤。

\subsection{六、拓扑、奇点与额外边界}\label{m1:ux516dux62d3ux6251ux5947ux70b9ux4e0eux989dux5916ux8fb9ux754c}

\subsubsection{1.
``边界的边界为空''}\label{m1:ux8fb9ux754cux7684ux8fb9ux754cux4e3aux7a7a}

开放曲面的边界 \(\partial S\) 是闭合曲线，因此

\[
\partial(\partial S)=\varnothing.
\]

这条几何事实解释了内部边抵消，也对应两个二阶恒等式：

\[
\nabla\times(\nabla\phi)=\mathbf0,
\qquad
\nabla\cdot(\nabla\times\mathbf A)=0.
\]

这些恒等式在分布意义下仍成立；若场有奇异性，集中项必须与普通部分一起计入。

\subsubsection{2.
局部无旋与全局势函数}\label{m1:ux5c40ux90e8ux65e0ux65cbux4e0eux5168ux5c40ux52bfux51fdux6570}

设 \(\Omega\) 路径连通，\(\mathbf A\) 在其中连续。下列三个命题等价：

\begin{enumerate}
\def\labelenumi{\arabic{enumi}.}
\tightlist
\item
  \(\mathbf A\) 的线积分只由端点决定；
\item
  对 \(\Omega\) 内任意分片光滑闭合曲线 \(C\)，都有
  \(\displaystyle\oint_C\mathbf A\cdot d\mathbf l=0\)；
\item
  存在单值标量场 \(\phi\)，使 \(\mathbf A=\nabla\phi\)。
\end{enumerate}

例如，若闭合曲线积分恒为零，固定参考点 \(P_0\) 并定义

\[
\phi(P)
=
\int_{P_0}^{P}\mathbf A\cdot d\mathbf l,
\]

积分与路径无关，沿坐标方向考察增量即可得到 \(\nabla\phi=\mathbf A\)。

若 \(\mathbf A\in C^1\)，存在势函数必然推出
\(\nabla\times\mathbf A=\mathbf0\)。反过来，
\(\nabla\times\mathbf A=\mathbf0\)
只是局部条件；当定义域单连通时，它才足以推出全局单值势函数。

\textbf{反例}：在去掉 \(z\) 轴的区域
\(\Omega=\mathbb R^3\setminus\{z\text{ 轴}\}\) 中，令

\[
\mathbf A
=
\frac{-y\mathbf e_x+x\mathbf e_y}{x^2+y^2}
=
\frac1\rho\mathbf e_{\phi}.
\]

此处 \(\rho=\sqrt{x^2+y^2}\) 表示到 \(z\) 轴的距离，不表示电荷密度。

在 \(\Omega\) 内逐点有 \(\nabla\times\mathbf A=\mathbf0\)，但对环绕
\(z\) 轴一周的圆周 \(C\)，

\[
\oint_C\mathbf A\cdot d\mathbf l=2\pi\neq0.
\]

因此它只有局部势函数 \(\phi\)；形式上的 \(\phi=\varphi\)
是多值方位角，不能作为整个 \(\Omega\)
上的单值势函数。缺失的轴线正是拓扑障碍所在。

\subsubsection{3.
更换积分曲面或体积的条件}\label{m1:ux66f4ux6362ux79efux5206ux66f2ux9762ux6216ux4f53ux79efux7684ux6761ux4ef6}

同一条有向闭合曲线可张成不同曲面。若 \(\mathbf A\)
在这些曲面上均光滑且曲面不穿过未定义集合，则斯托克斯定理保证旋度通量相同。若替换曲面穿过奇线，必须挖去奇线并保留新增的小边界。

同样，平方反比径向场在原点外散度为零，但包围原点的闭合面具有固定通量。普通高斯定理不能跨过未定义原点直接使用。

\subsubsection{4. 挖孔法}\label{m1:ux6316ux5b54ux6cd5}

对点奇点，挖去半径为 \(\varepsilon\) 的小球；对线奇点，挖去半径为
\(\varepsilon\)
的细管；对跳跃界面，用跨界薄柱或窄回路。先在剩余光滑区域使用积分定理，再令
\(\varepsilon\to0^+\)。

\begin{center}
\includegraphics[width=\linewidth,height=.40\textheight,keepaspectratio]{\detokenize{assets/module-01/tikz-m1-topic6-v4-07-puncture-extra-boundaries.png}}
\par\smallskip{\small 图示：挖去点奇点后电场只画在球壳区域内，内球面法向指向球心；挖去线奇点后内外边界方向相反}
\end{center}

典型分布恒等式为

\[
\nabla\cdot\frac{\mathbf r}{r^3}
=
4\pi\delta^{(3)}(\mathbf r),
\]

\[
\nabla\times
\left(
\frac{\kappa}{2\pi\rho}\mathbf e_{\phi}
\right)
=
\kappa\delta(x)\delta(y)\mathbf e_z.
\]

第一个由包围点源的小球面通量确定，第二个由环绕奇线的小圆周环量确定。

\subsection{七、麦克斯韦方程的积分形式与微分形式}\label{m1:ux4e03ux9ea6ux514bux65afux97e6ux65b9ux7a0bux7684ux79efux5206ux5f62ux5f0fux4e0eux5faeux5206ux5f62ux5f0f}

在 SI 单位制下，对固定在空间中的体积、曲面和边界回路，含电荷密度
\(\rho\) 与电流密度 \(\mathbf J\) 的微观麦克斯韦方程可成对写成

\[
\boxed{
\oint_{\partial V}\mathbf E\cdot d\mathbf S
=
\frac1{\varepsilon_0}\int_V\rho\,dV
\quad\Longleftrightarrow\quad
\nabla\cdot\mathbf E=\frac{\rho}{\varepsilon_0}
},
\]

\[
\boxed{
\oint_{\partial V}\mathbf B\cdot d\mathbf S=0
\quad\Longleftrightarrow\quad
\nabla\cdot\mathbf B=0
},
\]

\[
\boxed{
\oint_{\partial S}\mathbf E\cdot d\mathbf l
=
-\frac d{dt}\int_S\mathbf B\cdot d\mathbf S
\quad\Longleftrightarrow\quad
\nabla\times\mathbf E
=
-\frac{\partial\mathbf B}{\partial t}
},
\]

\[
\boxed{
\oint_{\partial S}\mathbf B\cdot d\mathbf l
=
\mu_0\int_S\mathbf J\cdot d\mathbf S
+
\mu_0\varepsilon_0
\frac d{dt}\int_S\mathbf E\cdot d\mathbf S
\quad\Longleftrightarrow\quad
\nabla\times\mathbf B
=
\mu_0\mathbf J
+
\mu_0\varepsilon_0
\frac{\partial\mathbf E}{\partial t}
}.
\]

双向箭头包含光滑性和任意区域条件；在奇点或界面上，应把右边理解为分布或先保留积分形式。

\begin{center}
\includegraphics[width=\linewidth,height=.40\textheight,keepaspectratio]{\detokenize{assets/module-01/figure-085-revised.pdf}}
\par\smallskip{\small 图示：四条麦克斯韦积分方程分别通过高斯定理或斯托克斯定理局部化为四条微分方程}
\end{center}

对安培--麦克斯韦方程取散度：

\[
0
=
\mu_0\nabla\cdot\mathbf J
+
\mu_0\varepsilon_0
\frac{\partial}{\partial t}
(\nabla\cdot\mathbf E).
\]

再用电场高斯定律，得到连续性方程

\[
\boxed{
\frac{\partial\rho}{\partial t}
+
\nabla\cdot\mathbf J=0
}.
\]

位移电流项保证了方程与局部电荷守恒相容。

对固定体积 \(V\) 积分并使用高斯定理，还得到等价的整体守恒式

\[
\boxed{
\frac d{dt}\int_V\rho\,dV
=
-\oint_{\partial V}\mathbf J\cdot d\mathbf S
}.
\]

闭合面上的净外向电流为正时，体积内的总电荷以同样速率减少。

若材料回路 \(C(t)\) 以局部速度 \(\mathbf v\)
运动，取与它共同运动的张成曲面 \(S(t)\)，则磁通定律写成

\[
\boxed{
\oint_{C(t)}
\left(
\mathbf E+\mathbf v\times\mathbf B
\right)\cdot d\mathbf l
=
-\frac d{dt}
\int_{S(t)}\mathbf B\cdot d\mathbf S
}.
\]

左边第二项是运动电动势。它说明运动曲面上的总磁通导数不能直接代替固定曲面形式中的
\(\displaystyle\int_S\partial\mathbf B/\partial t\cdot d\mathbf S\)。

\par\addvspace{.4\baselineskip}
\noindent\begin{minipage}{\linewidth}

\textbf{例 M1-56：导体棒的动生电动势}

一根长度为 \(L\) 的直导体棒以恒定速度 \(\mathbf v\) 在匀强磁场
\(\mathbf B\) 中作平移，设
\(\mathbf v\perp\mathbf B\)，棒的轴线同时垂直于 \(\mathbf v\) 与
\(\mathbf B\)，故棒沿 \(\mathbf v\times\mathbf B\)
所在直线放置。忽略端部效应，求两端之间运动电动势的大小；说明若规定的棒长积分方向反转，有向电动势的符号应如何处理。


\end{minipage}\par

\par\noindent\begin{minipage}{\linewidth}
\textbf{解答}

\begin{center}
\includegraphics[width=\linewidth,height=.40\textheight,keepaspectratio]{\detokenize{assets/module-01/examples/M1-56.pdf}}
\par\smallskip{\small 图 M1-56：导体棒的动生电动势}
\end{center}
\end{minipage}\par

图 M1-56中速度向右、磁场纸外，所以正电荷所受磁力沿棒向下；按该方向取
\(d\mathbf l\)，点乘是正的 \(Bv\,dl\)。沿图示棒长从 \(0\) 积到
\(L\)，得到有向电动势
\(BLv\)；反选从下到上的积分方向时结果反号，但题目所求大小不变。

长为 \(L\) 的导体棒在均匀磁场中以速度 \(v\) 垂直切割磁感线，且
\(\mathbf v\perp\mathbf B\)、导体棒同时垂直于 \(\mathbf v\) 与
\(\mathbf B\)，则运动电动势大小为

\[
\mathcal E
=
\int(\mathbf v\times\mathbf B)\cdot d\mathbf l
=BLv.
\]

\subsection{八、解题时如何选择积分定理}\label{m1:ux516bux89e3ux9898ux65f6ux5982ux4f55ux9009ux62e9ux79efux5206ux5b9aux7406}

积分定理始终成立并不意味着一定能直接求出场。高斯定律或安培环路定律只有在对称性足以使未知场从积分号中提出时，才成为高效求场工具。

\begin{center}
\includegraphics[width=\linewidth,height=.40\textheight,keepaspectratio]{\detokenize{assets/module-01/tikz-m1-topic6-v4-09-method-selection-flow.png}}
\par\smallskip{\small 图示：选择积分定理前依次检查目标量、对称性、取向、光滑性和区域类型}
\end{center}

可按以下顺序判断：

\begin{enumerate}
\def\labelenumi{\arabic{enumi}.}
\tightlist
\item
  \textbf{目标是端点差、通量还是环量}：分别考虑梯度定理、高斯定理或斯托克斯定理。
\item
  \textbf{对称性是否足够}：场的方向和大小能否在所选路径或曲面上简化。
\item
  \textbf{方向约定是否完整}：闭合面取外法向；开放曲面法向与边界正向成对选择。
\item
  \textbf{区域是否光滑}：是否含点源、线源、界面或场的未定义位置。
\item
  \textbf{区域是否固定}：时间导数能否移入积分号。
\item
  \textbf{是否要推出微分式}：积分关系是否对任意足够小区域成立。
\item
  \textbf{是否存在拓扑障碍}：闭合曲线能否在场的定义域内收缩，替换曲面是否会穿过奇线。
\end{enumerate}

\subsubsection{常见误区}\label{m1:ux5e38ux89c1ux8befux533a-1}

\begin{enumerate}
\def\labelenumi{\arabic{enumi}.}
\tightlist
\item
  只在一个特殊区域上验证积分等式，就宣称得到逐点微分方程。
\item
  把高斯面或安培回路选得很漂亮，却没有足够对称性把未知场提出积分号。
\item
  闭合曲面漏取某一片边界，或内边界仍误用普通外侧方向。
\item
  改变斯托克斯曲面时让新曲面穿过奇线。
\item
  把曲线切向分量、曲面法向分量和场的完整大小混为一谈。
\item
  对运动曲面直接交换时间导数与积分。
\item
  在奇点处使用普通微分式，却遗漏狄拉克分布贡献。
\end{enumerate}

\subsection{例题}\label{m1:ux4f8bux9898-2}

\par\addvspace{.4\baselineskip}
\noindent\begin{minipage}{\linewidth}

\subsubsection{同一曲面反转法向时如何保持斯托克斯定理（M1-57）}\label{m1:ux540cux4e00ux66f2ux9762ux53cdux8f6cux6cd5ux5411ux65f6ux5982ux4f55ux4fddux6301ux65afux6258ux514bux65afux5b9aux7406m1-57}

设 \(S\) 为有向开曲面，边界为 \(C\)，矢量场在邻域内一阶连续可微。最初
\(d\mathbf S=\mathbf n\,dS\)，\(C\)
的正向按右手定则选取。保持几何曲面不动，将其法向反转为
\(-\mathbf n\)，重新写出配对的斯托克斯等式，并解释曲线积分与曲面积分的符号为什么必须同时改变。


\end{minipage}\par

\par\noindent\begin{minipage}{\linewidth}
\textbf{解答}

\begin{center}
\includegraphics[width=\linewidth,height=.40\textheight,keepaspectratio]{\detokenize{assets/module-01/examples/M1-57.pdf}}
\par\smallskip{\small 图 M1-57：反转法向后的边界取向}
\end{center}
\end{minipage}\par

图 M1-57的左右两图分别给出同一几何曲面的两种定向。将左图的 \(\mathbf n\)
改为右图的 \(-\mathbf n\)
后，面积矢量与边界切向都反转，所以斯托克斯等式两边一起乘以
\(-1\)。解题时应先选图中的一对取向，再代入场，而不要在积分算完后只补一个负号。

若 \(S\) 的法向取 \(\mathbf n\)，边界 \(\partial S\) 取对应右手方向，则

\[
\oint_{\partial S}\mathbf A\cdot d\mathbf l
=\iint_S(\nabla\times\mathbf A)\cdot d\mathbf S.
\]

将法向改为 \(-\mathbf n\) 后，\(d\mathbf S\)
变号，边界方向也反转，因此等式两边同时变号。

\par\addvspace{.4\baselineskip}
\noindent\begin{minipage}{\linewidth}

\subsubsection{挖孔区域中的内边界不能丢掉（M1-58）}\label{m1:ux6316ux5b54ux533aux57dfux4e2dux7684ux5185ux8fb9ux754cux4e0dux80fdux4e22ux6389m1-58}

给定径向场 \(\mathbf F=C\mathbf e_r/r^2\)，\(r>0\)，其中 \(C\)
为常量。在球形区域 \(r<R\) 中挖去以原点为球心、半径为 \(\varepsilon\)
的小球，取
\(0<\varepsilon<R\)；所得球壳的内、外两边界均使用球壳实体的外法向。分别计算两边界通量，验证球壳上的高斯定理，并讨论
\(\varepsilon\to0\) 时内边界通量是否消失及其点源含义。


\end{minipage}\par

\par\noindent\begin{minipage}{\linewidth}
\textbf{解答}

\begin{center}
\includegraphics[width=\linewidth,height=.40\textheight,keepaspectratio]{\detokenize{assets/module-01/examples/M1-58.pdf}}
\par\smallskip{\small 图 M1-58：挖孔区域的内边界通量}
\end{center}
\end{minipage}\par

图
M1-58中的内边界属于挖孔区域，它的区域外法向指向球心，与外球面的外法向相反。两面上的平方反比场强分别乘以各自面积后，得到
\(+4\pi C\) 与
\(-4\pi C\)，在光滑挖孔区域内相加为零。小孔半径趋零并不会使后一通量消失，保留它才能补回原点源。

设 \(\mathbf F=C\mathbf e_r/r^2\)。在
\(V_{\varepsilon}=V\setminus B_{\varepsilon}\) 中有
\(\nabla\cdot\mathbf F=0\)，故

\[
\iint_{S_R}\mathbf F\cdot d\mathbf S
+\iint_{S_{\varepsilon}}^{\text{区域外法向}}
\mathbf F\cdot d\mathbf S=0.
\]

内边界对挖孔区域的外法向指向球心。令 \(\varepsilon\to0\)
时，该通量保持有限，正是原点集中源的积分表现。

\subsection{九、练习与解答}\label{m1:ux4e5dux7ec3ux4e60ux4e0eux89e3ux7b54-1}

\par\addvspace{.4\baselineskip}
\noindent\begin{minipage}{\linewidth}

\subsubsection{作业一：二维旋转场的环量与斯托克斯核对}\label{m1:ux4f5cux4e1aux4e00ux4e8cux7ef4ux65cbux8f6cux573aux7684ux73afux91cfux4e0eux65afux6258ux514bux65afux6838ux5bf9}

\textbf{题号：M1-59}

\textbf{层级：} 基础　　\textbf{方法：} 高斯定理、旋度与环量

采用右手直角坐标系，\(R>0\)；配对圆盘法向固定为
\(+\mathbf e_z\)，直接积分与斯托克斯计算须使用同一取向。

设

\[
\mathbf A=-y\,\mathbf e_x+x\,\mathbf e_y.
\]

求半径为 \(R\)、以原点为圆心、从 \(+z\) 方向看逆时针绕行的圆周 \(C\)
上的环量，并用斯托克斯定理核对。


\end{minipage}\par

\par\noindent\begin{minipage}{\linewidth}
\textbf{解答}

\begin{center}
\includegraphics[width=\linewidth,height=.40\textheight,keepaspectratio]{\detokenize{assets/module-01/examples/M1-59.pdf}}
\par\smallskip{\small 图 M1-59：圆周环量的直接积分}
\end{center}
\end{minipage}\par

图 M1-59中的圆周切向场与逆时针边界同向，面上纸外法向与
\(\nabla\times\mathbf A=2\mathbf e_z\)
同向。直接路线读取周长乘以切向场强 \(R\)，斯托克斯路线读取圆盘面积乘以
\(2\)，均得到 \(2\pi R^2\)。以下参数化应保持与图中边界相同的方向。

参数化圆周：

\[
\mathbf r(\theta)
=
R\cos\theta\,\mathbf e_x
+
R\sin\theta\,\mathbf e_y,
\qquad
0\le\theta\le2\pi.
\]

于是

\[
d\mathbf l
=
\left(
-R\sin\theta\,\mathbf e_x
+
R\cos\theta\,\mathbf e_y
\right)d\theta,
\]

\[
\mathbf A|_C
=
-R\sin\theta\,\mathbf e_x
+
R\cos\theta\,\mathbf e_y.
\]

因此

\[
\begin{aligned}
\oint_C\mathbf A\cdot d\mathbf l
&=
\int_0^{2\pi}
R^2(\sin^2\theta+\cos^2\theta)\,d\theta\\
&=
\boxed{2\pi R^2}.
\end{aligned}
\]

另一方面，

\[
\nabla\times\mathbf A
=
\left[
\frac{\partial x}{\partial x}
-
\frac{\partial(-y)}{\partial y}
\right]\mathbf e_z
=
2\mathbf e_z.
\]

逆时针边界对应圆盘法向 \(+\mathbf e_z\)，所以

\[
\int_S
(\nabla\times\mathbf A)\cdot d\mathbf S
=
\int_S2\,dS
=
2\pi R^2.
\]

两种方法一致。若把圆周方向改为顺时针，则法向必须改为
\(-\mathbf e_z\)，两边同时变为 \(-2\pi R^2\)。

\par\addvspace{.4\baselineskip}
\noindent\begin{minipage}{\linewidth}

\subsubsection{作业二：平方反比场、挖孔区域与点源通量}\label{m1:ux4f5cux4e1aux4e8cux5e73ux65b9ux53cdux6bd4ux573aux6316ux5b54ux533aux57dfux4e0eux70b9ux6e90ux901aux91cf}

\textbf{题号：M1-60}

\textbf{层级：} 基础　　\textbf{方法：} 分布与奇性、高斯定理、旋度与环量

在三维空间中取 \(\mathbf r\)
为原点到观察点的位移，\(\mathbf e_r=\mathbf r/r\)。涉及球壳内边界时，外法向按球壳实体而非球心径向确定。

设

\[
\mathbf A
=
\frac{\mathbf r}{r^3},
\qquad
r>0.
\]

\begin{enumerate}
\def\labelenumi{\arabic{enumi}.}
\tightlist
\item
  求半径为 \(R\) 的球面外向通量；
\item
  求 \(r>0\) 区域的普通散度；
\item
  在球壳 \(\varepsilon<r<R\) 上使用高斯定理，解释非零通量来自何处；
\item
  写出全空间分布形式。
\end{enumerate}


\end{minipage}\par

\par\noindent\begin{minipage}{\linewidth}
\textbf{解答}

\begin{center}
\includegraphics[width=\linewidth,height=.40\textheight,keepaspectratio]{\detokenize{assets/module-01/examples/M1-60.pdf}}
\par\smallskip{\small 图 M1-60：点源通量与挖孔法}
\end{center}
\end{minipage}\par

图 M1-60中外球面法向向外，内球面法向朝向被挖去的球心，因此内外通量分别为
\(4\pi\) 和
\(-4\pi\)。在球壳里普通散度为零，两者正好抵消；将小球补回全空间时，要用
\(4\pi\delta^{(3)}(\mathbf r)\) 记录那个不能随孔径消失的通量。

在球面 \(r=R\) 上，

\[
\mathbf A=\frac{\hat{\mathbf r}}{R^2},
\qquad
d\mathbf S
=\hat{\mathbf r}R^2\sin\theta\,d\theta\,d\phi.
\]

故

\[
\begin{aligned}
\Phi_R
&=
\oint_{r=R}\mathbf A\cdot d\mathbf S\\
&=
\int_0^{2\pi}\int_0^{\pi}
\sin\theta\,d\theta\,d\phi\\
&=
\boxed{4\pi}.
\end{aligned}
\]

在 \(r>0\)，

\[
\nabla\cdot\mathbf A
=
\frac1{r^2}\frac d{dr}
\left(
r^2\frac1{r^2}
\right)
=0.
\]

取球壳

\[
V_{\varepsilon}
=
\{\mathbf r:\varepsilon<r<R\}.
\]

其边界包括外球面和内球面。对球壳区域而言，内球面的外法向指向球心，即
\(-\hat{\mathbf r}\)，所以

\[
\Phi_{\mathrm{\text{外}}}=4\pi,
\qquad
\Phi_{\mathrm{\text{内}}}=-4\pi.
\]

因此

\[
\oint_{\partial V_{\varepsilon}}
\mathbf A\cdot d\mathbf S
=
4\pi-4\pi=0
=
\int_{V_{\varepsilon}}
\nabla\cdot\mathbf A\,dV.
\]

高斯定理在穿孔区域中完全成立。令 \(\varepsilon\to0^+\)
时，内球面并没有``变得可以忽略''，其通量始终为
\(-4\pi\)；把原点重新放回区域，必须用集中源补回这一贡献：

\[
\boxed{
\nabla\cdot\frac{\mathbf r}{r^3}
=
4\pi\delta^{(3)}(\mathbf r)
}.
\]

\par\addvspace{.4\baselineskip}
\noindent\begin{minipage}{\linewidth}

\subsubsection{作业三：格林第二恒等式与静电互易关系}\label{m1:ux4f5cux4e1aux4e09ux683cux6797ux7b2cux4e8cux6052ux7b49ux5f0fux4e0eux9759ux7535ux4e92ux6613ux5173ux7cfb}

\textbf{题号：M1-61}

\textbf{层级：} 决赛　　\textbf{方法：} 积分定理、静电互易

两组电荷均位于同一真空区域内，并使用同一个齐次狄利克雷边界；假设各函数足够光滑，使格林第二恒等式及下列积分成立。

在有界区域 \(V\) 中，两组电荷密度 \(\rho_1,\rho_2\) 分别产生电势
\(\varphi_1,\varphi_2\)，满足

\[
\nabla^2\varphi_1=-\frac{\rho_1}{\varepsilon_0},
\qquad
\nabla^2\varphi_2=-\frac{\rho_2}{\varepsilon_0},
\]

并且

\[
\varphi_1|_{\partial V}
=
\varphi_2|_{\partial V}
=0.
\]

证明互易关系

\[
\boxed{
\int_V\rho_1\varphi_2\,dV
=
\int_V\rho_2\varphi_1\,dV
}.
\]


\end{minipage}\par

\par\noindent\begin{minipage}{\linewidth}
\textbf{解答}

\begin{center}
\includegraphics[width=\linewidth,height=.40\textheight,keepaspectratio]{\detokenize{assets/module-01/examples/M1-61.pdf}}
\par\smallskip{\small 图 M1-61：同一区域的静电互易}
\end{center}
\end{minipage}\par

图 M1-61中两图的区域 \(V\)
与边界形状完全相同，交换的只是两组源密度及其电势在被积式中的角色。图上边界条件使格林第二恒等式的两项边界贡献都为零。随后把区域内的两个拉普拉斯分别换成泊松源项，移项就得到两图积分相等；这一步不需要把一般密度假定成点电荷。

对 \(\varphi_1,\varphi_2\) 使用格林第二恒等式：

\[
\int_V
\left(
\varphi_1\nabla^2\varphi_2
-
\varphi_2\nabla^2\varphi_1
\right)dV
=
\oint_{\partial V}
\left(
\varphi_1\frac{\partial\varphi_2}{\partial n}
-
\varphi_2\frac{\partial\varphi_1}{\partial n}
\right)dS.
\]

因为图 M1-61 的同一边界上两电势均为零，右边为零。代入泊松方程：

\[
\int_V
\left(
-\frac{\varphi_1\rho_2}{\varepsilon_0}
+
\frac{\varphi_2\rho_1}{\varepsilon_0}
\right)dV
=0.
\]

乘以 \(\varepsilon_0\) 并移项，得到

\[
\boxed{
\int_V\rho_1\varphi_2\,dV
=
\int_V\rho_2\varphi_1\,dV
}.
\]

这个结论说明：第一组电荷在第二组电势中的加权作用，与第二组电荷在第一组电势中的加权作用相同。它是静电互易定理、格林函数对称性和多种电容系数对称关系的数学基础。

\section{\texorpdfstring{专题七：狄拉克 \(\delta\)
分布及其物理应用（基础必讲与拔高选讲）}{专题七：狄拉克 \textbackslash delta 分布及其物理应用（基础必讲与拔高选讲）}}\label{m1:ux4e13ux9898ux4e03ux72c4ux62c9ux514b-delta-ux5206ux5e03ux53caux5176ux7269ux7406ux5e94ux7528ux57faux7840ux5fc5ux8bb2ux4e0eux62d4ux9ad8ux9009ux8bb2}

\begin{quote}
阅读准备：先积分，再谈 \(\delta\) 的意义 \(\delta\)
不是在原点取``无穷大''的普通函数。它通过对光滑函数的作用定义：\(\int f(x)\delta(x-x_0)\,dx=f(x_0)\)。因此判断系数和量纲时，应把它放回积分中检查总量。
\end{quote}

\par\addvspace{.4\baselineskip}
\noindent\begin{minipage}{\linewidth}

\textbf{例 M1-62：一维点电荷的支撑点}

用一维线密度 \(\rho(x)=q\delta(x-x_0)\) 表示位于 \(x=x_0\)
的点电荷，\(q\) 为有符号电荷量。取有限区间 \([a,b]\)，\(a<b\)，并排除
\(x_0\) 恰在端点的情形。分别计算 \(x_0\in(a,b)\) 和 \(x_0\notin[a,b]\)
时区间内的总电荷，解释结果与支撑点的位置关系。


\end{minipage}\par

\par\noindent\begin{minipage}{\linewidth}
\textbf{解答}

\begin{center}
\includegraphics[width=\linewidth,height=.40\textheight,keepaspectratio]{\detokenize{assets/module-01/examples/M1-62.pdf}}
\par\smallskip{\small 图 M1-62：一维点电荷的支撑点}
\end{center}
\end{minipage}\par

图 M1-62中的红色箭头只标出位于 \(x_0\)
的集中权重，不表示普通函数在该点的高度。绿色区间覆盖支撑点，故积分读到全部
\(q\)；橙色区间与支撑点分离，积分为零。这里若端点恰好落在
\(x_0\)，还需说明区间或正则化约定，图中两区间都避开了这种端点情形。

一维点电荷写成 \(\rho(x)=q\delta(x-x_0)\)，则任意包含 \(x_0\)
的区间内总电荷为 \(\int\rho(x)dx=q\)；不包含 \(x_0\) 时积分为零。

点电荷、理想细导线、面电荷层和瞬时冲量都集中在低维集合上。若只使用普通函数，它们要么表现为``无穷大''，要么在几乎处处为零后丢失有限总量。狄拉克
\(\delta\)
分布的作用，是用积分作用保留这些集中量，并把奇点、跳跃和理想化薄层统一写进场方程。

竞赛计算主要使用取样、缩放和坐标变换规则。测试函数、分布导数与奇点恒等式用于说明这些规则的来源，列为拔高内容；点、线、面源的表示及量纲检查则是后续计算的基础。

\subsection{\texorpdfstring{一、点源为何需要 \(\delta\)
分布（基础必讲）}{一、点源为何需要 \textbackslash delta 分布（基础必讲）}}\label{m1:ux4e00ux70b9ux6e90ux4e3aux4f55ux9700ux8981-delta-ux5206ux5e03ux57faux7840ux5fc5ux8bb2}

原点处点电荷 \(q\) 的库仑场为

\[
\mathbf E(\mathbf r)
=
\frac{q}{4\pi\varepsilon_0}
\frac{\mathbf r}{r^3}.
\]

在 \(r>0\) 的普通点上，

\[
\nabla\cdot\mathbf E=0,
\]

但任意包围原点的闭合曲面都有

\[
\oint_{\partial V}\mathbf E\cdot d\mathbf S
=
\frac q{\varepsilon_0}.
\]

两者并不矛盾。第一式只适用于场的光滑定义域，第二式检测到了被排除的原点。若希望高斯定律在整个空间中统一写成微分形式，就必须令

\[
\rho(\mathbf r)
=
q\,\delta^{(3)}(\mathbf r),
\]

使

\[
\nabla\cdot\mathbf E
=
\frac{q}{\varepsilon_0}
\delta^{(3)}(\mathbf r)
\]

在分布意义下成立。

\begin{center}
\includegraphics[width=\linewidth,height=.40\textheight,keepaspectratio]{\detokenize{assets/module-01/tikz-m1-topic7-v1-01-point-source-need-delta.jpg}}
\par\smallskip{\small 图示：普通散度只描述原点外的光滑区域，三维 delta
分布把有限总电荷集中在原点}
\end{center}

\(\delta\)
是通过积分读取测试函数局部信息的线性泛函，不能按普通函数在某一点取值。

\subsection{二、定义、近似恒等族与量纲}\label{m1:ux4e8cux5b9aux4e49ux8fd1ux4f3cux6052ux7b49ux65cfux4e0eux91cfux7eb2}

\subsubsection{1.
分布定义（拔高选讲）}\label{m1:ux5206ux5e03ux5b9aux4e49ux62d4ux9ad8ux9009ux8bb2}

记 \(C_c^{\infty}(\mathbb R)\)
为所有无穷次可微且具有紧支集的测试函数集合。位于 \(a\) 的狄拉克分布
\(\delta_a\) 定义为

\[
\boxed{
\langle\delta_a,f\rangle=f(a),
\qquad
f\in C_c^{\infty}(\mathbb R)
}.
\]

物理中常把这条配对关系写成积分形式

\[
\boxed{
\int_{-\infty}^{\infty}
f(x)\delta(x-a)\,dx
=f(a)
}.
\]

这就是筛选性质。式中的积分是分布对测试函数的作用，不应理解为先给
\(\delta(a)\)
一个普通数值再积分。因此，把``\(\delta(0)=\infty\)''当作普通函数等式不是合法定义；``宽度趋零、高度增大、总面积保持为
\(1\)''只能作为近似图像。

\subsubsection{2.
高斯近似恒等族（基础必讲）}\label{m1:ux9ad8ux65afux8fd1ux4f3cux6052ux7b49ux65cfux57faux7840ux5fc5ux8bb2}

取

\[
\delta_{\varepsilon}(x)
=
\frac1{\sqrt\pi\,\varepsilon}
\exp\left(-\frac{x^2}{\varepsilon^2}\right),
\qquad
\varepsilon>0.
\]

每个 \(\delta_{\varepsilon}\) 都是普通光滑函数，并满足

\[
\int_{-\infty}^{\infty}\delta_{\varepsilon}(x)\,dx=1.
\]

令 \(x=\varepsilon u\)，则

\[
\begin{aligned}
\int_{-\infty}^{\infty}
f(x)\delta_{\varepsilon}(x)\,dx
&=
\frac1{\sqrt\pi}
\int_{-\infty}^{\infty}
f(\varepsilon u)e^{-u^2}\,du\\
&\xrightarrow[\varepsilon\to0^+]{}
f(0).
\end{aligned}
\]

\(\delta_{\varepsilon}\) 的极限应在积分作用下理解；它收敛到 \(\delta\)
分布，而非某个逐点极限函数。

\begin{center}
\includegraphics[width=\linewidth,height=.40\textheight,keepaspectratio]{\detokenize{assets/module-01/tikz-m1-topic7-v1-02-distribution-approximate-identity-margin12-v11.png}}
\par\smallskip{\small 图示：高斯近似恒等族的宽度趋零、峰值增大、总面积保持为一，积分只保留取样点的函数值}
\end{center}

\par\addvspace{.4\baselineskip}
\noindent\begin{minipage}{\linewidth}

\textbf{例 M1-63：δ 对测试函数的筛选}

设 \(f(x)=1+x^2\)，\(a\)
为任意实常数，积分范围为整个实轴。按狄拉克分布的筛选定义计算
\(\int_{-\infty}^{\infty}f(x)\delta(x-a)\,dx\)，并解释为什么读取的是支撑点
\(x=a\) 处的函数值，而不是邻域平均。


\end{minipage}\par

\par\noindent\begin{minipage}{\linewidth}
\textbf{解答}

\begin{center}
\includegraphics[width=\linewidth,height=.40\textheight,keepaspectratio]{\detokenize{assets/module-01/examples/M1-63.pdf}}
\par\smallskip{\small 图 M1-63：δ 对测试函数的筛选}
\end{center}
\end{minipage}\par

图 M1-63中的绿色点位于测试函数曲线的 \(x=a\) 处，δ
的作用是读取该点函数值，而不是沿整条曲线累加。把测试函数代入 \(a\) 得
\(1+a^2\)，即右图筛选积分的结果。红色箭头只表示单位权重，与曲线上的纵坐标不能作峰高相乘。

对连续函数 \(f(x)=1+x^2\)，

\[
\int_{-\infty}^{\infty}
(1+x^2)\delta(x-a)\,dx
=1+a^2.
\]

\subsubsection{3.
量纲（基础必讲）}\label{m1:ux91cfux7eb2ux57faux7840ux5fc5ux8bb2}

若 \(x\) 的量纲为长度 \(L\)，由

\[
\int\delta(x-a)\,dx=1
\]

可知

\[
[\delta(x-a)]=L^{-1}.
\]

同理，

\[
[\delta^{(2)}(\mathbf r)]=L^{-2},
\qquad
[\delta^{(3)}(\mathbf r)]=L^{-3}.
\]

例如 \(q\delta^{(3)}(\mathbf r)\)
的量纲是电荷除以体积，正好是体电荷密度。量纲检查往往能立即发现漏写雅可比因子的问题。

\subsection{\texorpdfstring{三、一维 \(\delta\)
分布的运算规则}{三、一维 \textbackslash delta 分布的运算规则}}\label{m1:ux4e09ux4e00ux7ef4-delta-ux5206ux5e03ux7684ux8fd0ux7b97ux89c4ux5219}

\subsubsection{1.
平移、偶性与乘法}\label{m1:ux5e73ux79fbux5076ux6027ux4e0eux4e58ux6cd5}

平移后的分布满足

\[
\int f(x)\delta(x-a)\,dx=f(a).
\]

由变量反射可得偶性

\[
\boxed{\delta(-x)=\delta(x)}.
\]

若 \(f\) 在 \(a\) 附近光滑，则

\[
\boxed{
f(x)\delta(x-a)
=
f(a)\delta(x-a)
}.
\]

这个等式必须理解为两边对任意测试函数的作用相同。

\subsubsection{2.
尺度变换与绝对值}\label{m1:ux5c3aux5ea6ux53d8ux6362ux4e0eux7eddux5bf9ux503c}

设 \(a\neq0\)。对任意测试函数 \(f\)，令
\(y=ax\)，并同时处理积分方向，可得

\[
\begin{aligned}
\int_{-\infty}^{\infty}f(x)\delta(ax)\,dx
&=
\frac1{|a|}
\int_{-\infty}^{\infty}
f\left(\frac y a\right)\delta(y)\,dy\\
&=
\frac1{|a|}f(0).
\end{aligned}
\]

所以

\[
\boxed{
\delta(ax)=\frac1{|a|}\delta(x)
}.
\]

\par\noindent\begin{minipage}{\linewidth}
更一般地，若 \(a\neq0\)，

\[
\boxed{
\delta(ax-b)
=
\frac1{|a|}
\delta\left(x-\frac b a\right)
}.
\]
\end{minipage}\par

\par\addvspace{.4\baselineskip}
\noindent\begin{minipage}{\linewidth}

\textbf{例 M1-64：δ 的尺度变换}

将 \(\delta(3x-6)\) 视为实轴上的分布。把它改写为以 \(x=2\) 为支撑点的
\(\delta(x-2)\) 的倍数，并通过与任意光滑测试函数 \(f(x)\)
的积分配对说明比例系数的来源，注意变量尺度改变对积分测度的影响。


\end{minipage}\par

\par\noindent\begin{minipage}{\linewidth}
\textbf{解答}

\begin{center}
\includegraphics[width=\linewidth,height=.40\textheight,keepaspectratio]{\detokenize{assets/module-01/examples/M1-64.pdf}}
\par\smallskip{\small 图 M1-64：δ 的尺度变换}
\end{center}
\end{minipage}\par

图 M1-64中两根箭头表示同一个支撑点
\(x=2\)，为便于比较才略作横向错开。变换 \(u=3(x-2)\)
把局部横向尺度放大三倍，δ
的权重相应缩为三分之一，所以与测试函数配对后读取
\(f(2)/3\)。支撑位置与权重是这张图中需要分别核对的两件事。

\[
\delta(3x-6)=\frac13\delta(x-2).
\]

绝对值同时保证分布非负并消除变量代换方向反转带来的符号。

\subsubsection{3.
复合函数与单根条件}\label{m1:ux590dux5408ux51fdux6570ux4e0eux5355ux6839ux6761ux4ef6}

设 \(g\in C^1\)，在积分区间内只有有限个零点
\(x_i\)，并且每个零点都是单根：

\[
g(x_i)=0,
\qquad
g'(x_i)\neq0.
\]

在每个根附近，\(g\) 都可作为局部坐标。分区后逐段作变量代换，得到

\[
\boxed{
\delta(g(x))
=
\sum_i
\frac{\delta(x-x_i)}{|g'(x_i)|}
}.
\]

相应的积分公式为

\[
\boxed{
\int f(x)\delta(g(x))\,dx
=
\sum_i\frac{f(x_i)}{|g'(x_i)|}
}.
\]

\begin{center}
\includegraphics[width=\linewidth,height=.40\textheight,keepaspectratio]{\detokenize{assets/module-01/tikz-m1-topic7-v1-03-scaling-composite-roots.png}}
\par\smallskip{\small 图示：尺度变换由长度缩放产生绝对值，复合 delta
分布中每个单根按局部斜率倒数加权}
\end{center}

若 \(g'(x_i)=0\)，上述公式失效。例如 \(g(x)=x^2\) 在原点有重根，不能把
\(\delta(x^2)\) 机械写成含 \(1/|2x|\)
的表达式；对应的正则化积分一般发散，必须回到具体物理模型重新处理。

\subsubsection{4.
分布导数（拔高选讲）}\label{m1:ux5206ux5e03ux5bfcux6570ux62d4ux9ad8ux9009ux8bb2}

普通函数的分部积分提示我们定义

\[
\boxed{
\langle\delta_a',f\rangle
=
-f'(a)
}.
\]

因此

\[
\boxed{
\int f(x)\delta'(x-a)\,dx
=
-f'(a)
}.
\]

更一般地，

\[
\langle\delta_a^{(n)},f\rangle
=
(-1)^n f^{(n)}(a).
\]

乘积恒等式为

\[
\boxed{
f(x)\delta'(x-a)
=
f(a)\delta'(x-a)
-f'(a)\delta(x-a)
}.
\]

\par\addvspace{.4\baselineskip}
\noindent\begin{minipage}{\linewidth}

\textbf{例 M1-65：δ′ 的负导数作用}

设 \(a\) 为实常数，\(\delta'(x-a)\) 表示对其自变量 \(x-a\)
求导所得的狄拉克分布导数。积分在整个实轴上进行，按分布配对理解。计算
\(\int_{-\infty}^{\infty}e^x\delta'(x-a)\,dx\)，解释作用到测试函数导数时负号的来源。


\end{minipage}\par

\par\noindent\begin{minipage}{\linewidth}
\textbf{解答}

\begin{center}
\includegraphics[width=\linewidth,height=.40\textheight,keepaspectratio]{\detokenize{assets/module-01/examples/M1-65.pdf}}
\par\smallskip{\small 图 M1-65：δ′ 的负导数作用}
\end{center}
\end{minipage}\par

图 M1-65中绿色切线强调 \(f=e^x\) 在 \(a\) 处的局部斜率；δ′
读取的不是红点的函数值，而是该斜率的负值。将分布导数转移到测试函数上产生一个负号，再代入
\(f^{\prime}(a)=e^a\)，结果为
\(-e^a\)。图的曲线按比例缩放，分布恒等式不受绘图尺度影响。

\[
\int e^x\delta'(x-a)\,dx=-e^a.
\]

\subsection{四、阶跃、跳跃与界面分布}\label{m1:ux56dbux9636ux8dc3ux8df3ux8dc3ux4e0eux754cux9762ux5206ux5e03}

\subsubsection{1.
海维赛德阶跃函数}\label{m1:ux6d77ux7ef4ux8d5bux5fb7ux9636ux8dc3ux51fdux6570}

定义

\[
H(x-a)
=
\begin{cases}
0,&x<a,\\
1,&x>a.
\end{cases}
\]

\(H(a)\) 的单点取值不影响分布。对测试函数分部积分：

\[
\begin{aligned}
\langle H',f\rangle
&=-\langle H,f'\rangle\\
&=-\int_a^{\infty} f'(x)\,dx\\
&=f(a),
\end{aligned}
\]

所以

\[
\boxed{
\frac d{dx}H(x-a)=\delta(x-a)
}.
\]

\par\addvspace{.4\baselineskip}
\noindent\begin{minipage}{\linewidth}

\textbf{例 M1-66：冲量与动量跃变}

一个质点沿固定直线运动，在 \(t=t_0\) 前的动量为 \(p_-\)。它在 \(t_0\)
受到有符号冲量 \(I\)，除这一瞬时冲击外不受其他力。用阶跃函数
\(H(t-t_0)\) 写出各时刻的动量 \(p(t)\)，在分布意义下求冲击力
\(F=dp/dt\)，并用跨越 \(t_0\) 的短时间积分核对总冲量。


\end{minipage}\par

\par\noindent\begin{minipage}{\linewidth}
\textbf{解答}

\begin{center}
\includegraphics[width=\linewidth,height=.40\textheight,keepaspectratio]{\detokenize{assets/module-01/examples/M1-66.pdf}}
\par\smallskip{\small 图 M1-66：冲量与动量跃变}
\end{center}
\end{minipage}\par

图 M1-66中的上下两条水平段分别代表冲量前后的动量，竖向差值是
\(I\)。先写成阶跃形式
\(p_-+I H(t-t_0)\)，再对时间求分布导数，阶跃的跃变量成为 δ
的系数。将力在 \(t_0\)
附近积分，读回同一个动量差，这也检验了解答中的权重。

若一次瞬时冲量 \(I\) 使动量在 \(t=t_0\) 发生跃变，可写成

\[
p(t)=p_-+I H(t-t_0).
\]

于是冲击力在分布意义下为

\[
F(t)=\frac{dp}{dt}=I\delta(t-t_0),
\]

并且

\[
\int_{t_0-\eta}^{t_0+\eta}F(t)\,dt=I
\]

对任意 \(\eta>0\) 都成立。

\subsubsection{2.
分段光滑函数的导数}\label{m1:ux5206ux6bb5ux5149ux6ed1ux51fdux6570ux7684ux5bfcux6570}

设 \(F\) 在 \(x=a\) 两侧分段光滑，跳跃量为

\[
[F]_a=F(a^+)-F(a^-).
\]

则分布导数为

\[
\boxed{
\frac{dF}{dx}
=
\left(\frac{dF}{dx}\right)_{\mathrm{\text{普通}}}
+[F]_a\delta(x-a)
}.
\]

普通导数只描述两侧内部变化，\(\delta\) 项记录跨越界面时的有限跳跃。

\subsubsection{\texorpdfstring{3. 曲面上的 \(\delta\)
分布}{3. 曲面上的 \textbackslash delta 分布}}\label{m1:ux66f2ux9762ux4e0aux7684-delta-ux5206ux5e03}

设光滑曲面由

\[
g(\mathbf r)=0
\]

给出，并且曲面上 \(\nabla g\neq\mathbf0\)。定义曲面分布

\[
\boxed{
\delta_S(\mathbf r)
=
\delta(g(\mathbf r))|\nabla g(\mathbf r)|
}.
\]

它满足

\[
\int_{\mathbb R^3}
f(\mathbf r)\delta_S(\mathbf r)\,dV
=
\int_S f\,dS.
\]

若

\[
\hat{\mathbf n}
=
\frac{\nabla g}{|\nabla g|},
\]

则

\[
\boxed{
\nabla H(g)
=
\delta(g)\nabla g
=
\hat{\mathbf n}\,\delta_S
}.
\]

\begin{center}
\includegraphics[width=\linewidth,height=.40\textheight,keepaspectratio]{\detokenize{assets/module-01/figure-095-revised.pdf}}
\par\smallskip{\small 图示：阶跃函数的分布导数是 delta
分布，界面跳跃在导数中产生集中在曲面上的附加项}
\end{center}

这条关系是推导电磁场边界条件、面电荷和面电流分布形式的直接工具。

\subsection{\texorpdfstring{五、多维 \(\delta\)
分布与坐标变换}{五、多维 \textbackslash delta 分布与坐标变换}}\label{m1:ux4e94ux591aux7ef4-delta-ux5206ux5e03ux4e0eux5750ux6807ux53d8ux6362}

\subsubsection{1.
直角坐标中的定义}\label{m1:ux76f4ux89d2ux5750ux6807ux4e2dux7684ux5b9aux4e49}

三维狄拉克分布定义为

\[
\boxed{
\delta^{(3)}(\mathbf r-\mathbf r_0)
=
\delta(x-x_0)\delta(y-y_0)\delta(z-z_0)
}.
\]

它满足

\[
\int_{\mathbb R^3}
f(\mathbf r)
\delta^{(3)}(\mathbf r-\mathbf r_0)
\,dV
=f(\mathbf r_0).
\]

若 \(A\) 是可逆的 \(n\times n\) 常矩阵，则

\[
\boxed{
\delta^{(n)}(A\mathbf x)
=
\frac1{|\det A|}
\delta^{(n)}(\mathbf x)
}.
\]

一维公式中的 \(|a|\) 正是多维雅可比行列式的特例。

\subsubsection{2.
一般曲线坐标中的雅可比因子}\label{m1:ux4e00ux822cux66f2ux7ebfux5750ux6807ux4e2dux7684ux96c5ux53efux6bd4ux56e0ux5b50}

设正则坐标变换

\[
\mathbf r=\mathbf r(q_1,q_2,q_3)
\]

的体积元为

\[
dV=J(q_1,q_2,q_3)\,dq_1dq_2dq_3,
\qquad J>0.
\]

若 \(\mathbf r_0\) 是坐标正则点，则

\[
\boxed{
\delta^{(3)}(\mathbf r-\mathbf r_0)
=
\frac{
\delta(q_1-q_{10})
\delta(q_2-q_{20})
\delta(q_3-q_{30})
}{J(q_{10},q_{20},q_{30})}
}.
\]

在圆柱坐标的正则点 \(\rho_0>0\)，

\[
\delta^{(3)}(\mathbf r-\mathbf r_0)
=
\frac{
\delta(\rho-\rho_0)
\delta(\phi-\phi_0)
\delta(z-z_0)
}{\rho_0}.
\]

在球坐标的正则点 \(r_0>0\)、\(0<\theta_0<\pi\)，

\[
\delta^{(3)}(\mathbf r-\mathbf r_0)
=
\frac{
\delta(r-r_0)
\delta(\theta-\theta_0)
\delta(\phi-\phi_0)
}{r_0^2\sin\theta_0}.
\]

坐标轴、球心和极点处角坐标不唯一，不能机械套用上式；此时应返回直角坐标表达或直接使用积分定义。

方位角具有 \(2\pi\) 周期。若积分区间跨越方位角分支端点，应使用周期分布

\[
\delta_{2\pi}(\phi-\phi_0)
=
\sum_{k\in\mathbb Z}
\delta(\phi-\phi_0-2\pi k),
\]

而不是把 \(\phi=0\) 与 \(\phi=2\pi\) 当成两个不同方向。

\begin{center}
\includegraphics[width=\linewidth,height=.40\textheight,keepaspectratio]{\detokenize{assets/module-01/tikz-m1-topic7-v1-05-multidimensional-delta-jacobian-safe-v7.png}}
\par\smallskip{\small 图示：一维、二维和三维 delta
分布的量纲随支撑维数变化，曲线坐标表达必须除以体积元雅可比}
\end{center}

\subsection{六、点源、线源、面源与运动点电荷}\label{m1:ux516dux70b9ux6e90ux7ebfux6e90ux9762ux6e90ux4e0eux8fd0ux52a8ux70b9ux7535ux8377}

\subsubsection{1. 点电荷}\label{m1:ux70b9ux7535ux8377}

位于 \(\mathbf r_0\) 的点电荷 \(q\) 写成

\[
\boxed{
\rho(\mathbf r)
=
q\,\delta^{(3)}(\mathbf r-\mathbf r_0)
}.
\]

若 \(\mathbf r_0\) 位于区域 \(V\) 的内部，则

\[
\int_V\rho\,dV=q.
\]

\subsubsection{2. 曲线电荷}\label{m1:ux66f2ux7ebfux7535ux8377}

设曲线以弧长 \(s\) 参数化为 \(\mathbf R(s)\)，线电荷密度为
\(\lambda(s)\)，则等效体密度为

\[
\boxed{
\rho_{\mathrm{\text{线}}}(\mathbf r)
=
\int_C
\lambda(s)
\delta^{(3)}(\mathbf r-\mathbf R(s))\,ds
}.
\]

它满足

\[
\int_{\mathbb R^3}
f(\mathbf r)\rho_{\mathrm{\text{线}}}(\mathbf r)\,dV
=
\int_C\lambda(s)f(\mathbf R(s))\,ds.
\]

特别地，沿 \(z\) 轴的均匀无限长线电荷可写为

\[
\boxed{
\rho(\mathbf r)=\lambda\delta(x)\delta(y)
}.
\]

\subsubsection{3. 曲面电荷}\label{m1:ux66f2ux9762ux7535ux8377}

设曲面电荷密度为 \(\sigma(\mathbf r')\)，则

\[
\boxed{
\rho_{\mathrm{\text{面}}}(\mathbf r)
=
\int_S
\sigma(\mathbf r')
\delta^{(3)}(\mathbf r-\mathbf r')\,dS'
}.
\]

若 \(S\) 由 \(g(\mathbf r)=0\) 给出，也可写成

\[
\boxed{
\rho_{\mathrm{\text{面}}}(\mathbf r)
=
\sigma(\mathbf r)\delta(g)|\nabla g|
}.
\]

这里可把只定义在曲面上的 \(\sigma\)
任意光滑延拓到曲面邻域；\(\delta(g)\)
最终只读取其曲面值，因此结果与延拓方式无关。

例如，\(z=0\) 平面上的均匀面电荷为

\[
\rho(\mathbf r)=\sigma\delta(z).
\]

\subsubsection{4.
运动点电荷与连续性方程}\label{m1:ux8fd0ux52a8ux70b9ux7535ux8377ux4e0eux8fdeux7eedux6027ux65b9ux7a0b}

若点电荷轨迹为 \(\mathbf r_q(t)\)，速度为 \(\mathbf v_q(t)\)，则

\[
\boxed{
\rho(\mathbf r,t)
=q\delta^{(3)}(\mathbf r-\mathbf r_q(t)),
\qquad
\mathbf J(\mathbf r,t)
=q\mathbf v_q(t)
\delta^{(3)}(\mathbf r-\mathbf r_q(t))
}.
\]

利用链式法则的分布形式，

\[
\frac{\partial\rho}{\partial t}
=
-q\mathbf v_q\cdot\nabla
\delta^{(3)}(\mathbf r-\mathbf r_q),
\]

而

\[
\nabla\cdot\mathbf J
=
q\mathbf v_q\cdot\nabla
\delta^{(3)}(\mathbf r-\mathbf r_q),
\]

故

\[
\boxed{
\frac{\partial\rho}{\partial t}
+\nabla\cdot\mathbf J=0
}.
\]

\begin{center}
\includegraphics[width=\linewidth,height=.40\textheight,keepaspectratio]{\detokenize{assets/module-01/tikz-m1-topic7-v1-06-point-line-surface-sources.jpg}}
\par\smallskip{\small 图示：点电荷、线电荷和面电荷分别由三维 delta
分布集中到点、曲线和曲面，积分后恢复原有总量}
\end{center}

\subsection{七、奇异场的分布恒等式（拔高选讲）}\label{m1:ux4e03ux5947ux5f02ux573aux7684ux5206ux5e03ux6052ux7b49ux5f0fux62d4ux9ad8ux9009ux8bb2}

\subsubsection{1. 点源散度}\label{m1:ux70b9ux6e90ux6563ux5ea6}

令

\[
\mathbf F(\mathbf r)=\frac{\mathbf r}{r^3},
\qquad r>0.
\]

在原点外，\(\nabla\cdot\mathbf F=0\)。要确定全空间中的分布散度，取任意测试函数
\(f\in C_c^{\infty}(\mathbb R^3)\)。选择足够大的 \(R\)，使
\(\operatorname{supp}f\subset B_R\) 且 \(f\) 在 \(\partial B_R\)
附近为零；再挖去半径为 \(\varepsilon\) 的小球，记

\[
\Omega_{\varepsilon}
=
B_R\setminus\overline{B_{\varepsilon}}.
\]

按分布导数定义，

\[
\left\langle
\nabla\cdot\mathbf F,f
\right\rangle
=
-\lim_{\varepsilon\to0^+}
\int_{\Omega_{\varepsilon}}
\mathbf F\cdot\nabla f\,dV.
\]

由于 \(\nabla\cdot\mathbf F=0\)，在光滑穿孔区域中使用高斯定理：

\[
\int_{\Omega_{\varepsilon}}
\mathbf F\cdot\nabla f\,dV
=
\oint_{\partial\Omega_{\varepsilon}}
f\mathbf F\cdot d\mathbf S.
\]

测试函数在外边界附近为零，只有内球面贡献。穿孔区域的内球面外法向为
\(-\hat{\mathbf r}\)，所以

\[
\begin{aligned}
\oint_{r=\varepsilon}
f\mathbf F\cdot d\mathbf S
&=
-\int_{S^2}f(\varepsilon\hat{\mathbf r})\,d\Omega\\
&\xrightarrow[\varepsilon\to0^+]{}
-4\pi f(\mathbf0).
\end{aligned}
\]

代回分布定义得到

\[
\left\langle
\nabla\cdot\frac{\mathbf r}{r^3},f
\right\rangle
=4\pi f(\mathbf0),
\]

即

\[
\boxed{
\nabla\cdot\frac{\mathbf r}{r^3}
=
4\pi\delta^{(3)}(\mathbf r)
}.
\]

又因为

\[
\nabla\left(\frac1r\right)
=
-\frac{\mathbf r}{r^3},
\]

所以

\[
\boxed{
\nabla^2\left(\frac1r\right)
=
-4\pi\delta^{(3)}(\mathbf r)
}.
\]

\begin{center}
\includegraphics[width=\linewidth,height=.40\textheight,keepaspectratio]{\detokenize{assets/module-01/tikz-m1-topic7-v1-07-punctured-ball-test-function.jpg}}
\par\smallskip{\small 图示：严格证明平方反比场的分布散度时先挖去原点小球，测试函数使外边界消失，内边界极限保留
\(4\pi\) 贡献}
\end{center}

这个推导不能改成``直接在含原点的球体上使用普通高斯定理''，因为原场在原点不满足定理的光滑性条件。

\subsubsection{2.
线奇点的旋度}\label{m1:ux7ebfux5947ux70b9ux7684ux65cbux5ea6}

在去掉 \(z\) 轴的区域内令

\[
\mathbf A
=
\frac1{2\pi\rho}\mathbf e_{\phi}.
\]

轴外逐点有

\[
\nabla\times\mathbf A=\mathbf0,
\]

但任意绕 \(z\) 轴一周的正向圆周 \(C\) 都有

\[
\oint_C\mathbf A\cdot d\mathbf l=1.
\]

更一般地，若闭合曲线在不穿过轴线的变形中绕轴 \(N\) 周，则

\[
\oint_C\mathbf A\cdot d\mathbf l=N.
\]

取法向为 \(+\mathbf e_z\) 且横截轴线一次的圆盘
\(S\)。先挖去轴线附近的小圆盘，在剩余圆环上使用斯托克斯定理，再令孔径趋于零，外边界环量保留为
\(1\)。另一方面，若分布旋度写成

\[
\nabla\times\mathbf A
=
c\,\delta(x)\delta(y)\mathbf e_z,
\]

则它穿过 \(S\) 的通量为 \(c\)。与边界环量比较可得 \(c=1\)。

因此全空间的分布旋度为

\[
\boxed{
\nabla\times
\left(
\frac1{2\pi\rho}\mathbf e_{\phi}
\right)
=
\delta(x)\delta(y)\mathbf e_z
}.
\]

这与点源散度完全平行：普通导数描述奇线外的光滑区域，\(\delta\)
项补回绕奇线的有限环量。

\begin{center}
\includegraphics[width=\linewidth,height=.40\textheight,keepaspectratio]{\detokenize{assets/module-01/tikz-m1-topic7-v1-08-line-singularity-circulation.jpg}}
\par\smallskip{\small 图示：绕轴场在轴外旋度为零，但任意环绕轴线的闭合曲线都有单位环量，分布旋度集中在轴线上}
\end{center}

\subsection{\texorpdfstring{八、\(\delta\)
分布的使用条件与常见错误}{八、\textbackslash delta 分布的使用条件与常见错误}}\label{m1:ux516bdelta-ux5206ux5e03ux7684ux4f7fux7528ux6761ux4ef6ux4e0eux5e38ux89c1ux9519ux8bef}

\begin{enumerate}
\def\labelenumi{\arabic{enumi}.}
\tightlist
\item
  \textbf{先确定集中集合的维数}：点源用三维
  \(\delta\)，直线源常用两个横向 \(\delta\)，平面源常用一个法向
  \(\delta\)。
\item
  \textbf{再做量纲检查}：\(\delta\)
  的量纲必须补足从点、线或面密度到体密度所缺的长度幂次。
\item
  \textbf{变量变换必须带绝对雅可比}：一维是 \(|g'(x_i)|\)，多维是
  \(|\det A|\) 或曲线坐标体积因子。
\item
  \textbf{复合公式只适用于单根}：出现 \(g'(x_i)=0\) 时不能直接套公式。
\item
  \textbf{分布导数按分部积分定义}：每求一次导数都会给测试函数导数带来一个负号。
\item
  \textbf{奇点处不能直接使用普通积分定理}：先挖孔，在光滑区域使用定理，再取极限。
\item
  \textbf{界面跳跃会产生面分布}：只计算两侧普通导数会漏掉 \(\delta_S\)
  项。
\item
  \textbf{一般不能随意相乘分布}：\(\delta^2\)、\(\delta(x^2)\)
  等表达式通常没有自然的经典分布意义。
\item
  \textbf{支撑点位于积分边界时要说明约定}：端点贡献取决于所用正则化或边界分布定义，不能未经说明直接按内部点处理。
\end{enumerate}

\begin{center}
\includegraphics[width=\linewidth,height=.40\textheight,keepaspectratio]{\detokenize{assets/module-01/tikz-m1-topic7-v1-09-delta-method-flow.png}}
\par\smallskip{\small 图示：处理集中源时依次判断支撑维数、量纲、雅可比、根的类型以及是否需要挖孔或加入界面分布}
\end{center}

\subsection{例题}\label{m1:ux4f8bux9898-3}

\par\addvspace{.4\baselineskip}
\noindent\begin{minipage}{\linewidth}

\subsubsection{\texorpdfstring{复合 \(\delta\)
分布为什么出现绝对值（M1-67）}{复合 \textbackslash delta 分布为什么出现绝对值（M1-67）}}\label{m1:ux590dux5408-delta-ux5206ux5e03ux4e3aux4ec0ux4e48ux51faux73b0ux7eddux5bf9ux503cm1-67}

设实函数 \(g\) 在其零点邻域光滑，所有零点 \(x_i\) 都是单零点，即
\(g'(x_i)\ne0\)。从局部变量伸缩出发，说明将 \(\delta(g(x))\)
展成各支撑点的狄拉克分布时，为何系数要含导数的绝对值。以
\(g(x)=x^2-a^2\)、\(a>0\) 为例写出展开式，并解释两个支撑点的权重。


\end{minipage}\par

\par\noindent\begin{minipage}{\linewidth}
\textbf{解答}

\begin{center}
\includegraphics[width=\linewidth,height=.40\textheight,keepaspectratio]{\detokenize{assets/module-01/examples/M1-67.pdf}}
\par\smallskip{\small 图 M1-67：复合 δ 的单根权重}
\end{center}
\end{minipage}\par

图 M1-67中的两个零点为 \(-a\) 与 \(a\)，曲线在两根处的斜率分别是 \(-2a\)
和 \(2a\)。把每个根附近的小横向区间换到 \(g\)
坐标时，权重使用长度缩放率的绝对值，所以两个根的贡献都带
\(1/(2a)\)。不能让左根的负斜率生成负权重；下面先作局部缩放，再把两个筛选贡献相加。

若 \(g(x)\) 的单零点为 \(x_i\)，在零点附近
\(g(x)\simeq g'(x_i)(x-x_i)\)。变量伸缩给出

\[
\delta(g(x))=\sum_i\frac{\delta(x-x_i)}{|g'(x_i)|}.
\]

例如 \(g(x)=x^2-a^2\)、\(a>0\) 时，

\[
\delta(x^2-a^2)=\frac{\delta(x-a)+\delta(x+a)}{2a}.
\]

绝对值保证坐标方向反转不会改变分布的筛选结果。

\par\addvspace{.4\baselineskip}
\noindent\begin{minipage}{\linewidth}

\subsubsection{\texorpdfstring{点源、线源和面源各需要几个
\(\delta\)（M1-68）}{点源、线源和面源各需要几个 \textbackslash delta（M1-68）}}\label{m1:ux70b9ux6e90ux7ebfux6e90ux548cux9762ux6e90ux5404ux9700ux8981ux51e0ux4e2a-deltam1-68}

在三维直角坐标系中分别考虑三个互不叠加的理想电荷模型：原点处总电荷为
\(Q\) 的点电荷；沿整条 \(z\) 轴、线密度为常量 \(\lambda\)
的线电荷；位于整个 \(z=0\) 平面、面密度为常量 \(\sigma\)
的面电荷。将三者各自写成三维体电荷密度分布，并通过空间受约束的方向数解释所需
\(\delta\) 因子的个数。


\end{minipage}\par

\par\noindent\begin{minipage}{\linewidth}
\textbf{解答}

\begin{center}
\includegraphics[width=\linewidth,height=.40\textheight,keepaspectratio]{\detokenize{assets/module-01/examples/M1-68.pdf}}
\par\smallskip{\small 图 M1-68：点、线、面源的约束维数}
\end{center}
\end{minipage}\par

图 M1-68 的三个坐标系分别标出各源的支撑：左图只有原点 \(O\)，要求
\(x=y=z=0\)；中图沿整条 \(z\) 轴延伸，只要求 \(x=y=0\)，\(z\)
自由；右图沿整个 \(z=0\) 面延伸，只约束 \(z\)，而 \(x,y\)
自由。图中的有限线段和面片只表示无限源的局部。每个被约束的独立方向对应一个
δ 因子，未约束方向仍需保留线积分或面积分。

三维空间中，原点点电荷、\(z\) 轴线电荷和 \(z=0\) 面电荷可分别写成

\[
\rho_Q=Q\delta(x)\delta(y)\delta(z),
\]

\[
\rho_{\lambda}=\lambda\delta(x)\delta(y),
\qquad
\rho_{\sigma}=\sigma\delta(z).
\]

每约束一个空间方向，就需要一个相应的 \(\delta\)
因子；系数的量纲随源的维数改变。

在左图中对包含原点的体积积分得到 \(Q\)；在中图中对垂直于 \(z\)
轴且包含轴线的横截面积分得到 \(\lambda\)；在右图中沿穿过 \(z=0\)
的法向积分得到 \(\sigma\)。因此三者虽分别含三个、两个和一个 δ
因子，体密度的量纲都为 \(\mathrm{C\,m^{-3}}\)，与图下方的积分检验一致。

\subsection{九、练习与解答}\label{m1:ux4e5dux7ec3ux4e60ux4e0eux89e3ux7b54-2}

\par\addvspace{.4\baselineskip}
\noindent\begin{minipage}{\linewidth}

\subsubsection{\texorpdfstring{作业一：复合 \(\delta\)
分布的代数积分}{作业一：复合 \textbackslash delta 分布的代数积分}}\label{m1:ux4f5cux4e1aux4e00ux590dux5408-delta-ux5206ux5e03ux7684ux4ee3ux6570ux79efux5206}

\textbf{题号：M1-69}

\textbf{层级：} 基础　　\textbf{方法：} 分布与奇性

积分按实轴上的分布积分理解，\(\delta\)
为一维狄拉克分布；须分别计入其复合自变量的每个单零点及相应权重。

计算

\[
I
=
\int_{-\infty}^{\infty}
(x^2+\sin x)
\delta(x^2-4x+3)\,dx.
\]


\end{minipage}\par

\par\noindent\begin{minipage}{\linewidth}
\textbf{解答}

\begin{center}
\includegraphics[width=\linewidth,height=.40\textheight,keepaspectratio]{\detokenize{assets/module-01/examples/M1-69.pdf}}
\par\smallskip{\small 图 M1-69：二次复合 δ 的代数积分}
\end{center}
\end{minipage}\par

图 M1-69中抛物线与横轴交于 \(1,3\)，两处切线斜率大小均为
\(2\)，故每个支撑点的权重都是 \(1/2\)。先按这两个根把复合 δ
拆开，再分别读取 \(x^2+\sin x\) 在两点的值。最终常数项 \(5\) 来自
\((1+9)/2\)，两个正弦项也必须分别保留，不能只取图中的一个交点。

令

\[
g(x)=x^2-4x+3=(x-1)(x-3).
\]

图 M1-69 的两个零点 \(x_1=1\)、\(x_2=3\) 都是单根，而且

\[
g'(x)=2x-4,
\qquad
|g'(1)|=|g'(3)|=2.
\]

所以

\[
\delta(g(x))
=
\frac12\delta(x-1)
+
\frac12\delta(x-3).
\]

代入筛选公式：

\[
\begin{aligned}
I
&=
\frac12(1+\sin 1)
+
\frac12(9+\sin 3)\\
&=
\boxed{
5+\frac12(\sin 1+\sin 3)
}.
\end{aligned}
\]

\par\addvspace{.4\baselineskip}
\noindent\begin{minipage}{\linewidth}

\subsubsection{\texorpdfstring{作业二：\(\delta'\)
的筛选作用与乘积恒等式}{作业二：\textbackslash delta\textquotesingle{} 的筛选作用与乘积恒等式}}\label{m1:ux4f5cux4e1aux4e8cdelta-ux7684ux7b5bux9009ux4f5cux7528ux4e0eux4e58ux79efux6052ux7b49ux5f0f}

\textbf{题号：M1-70}

\textbf{层级：} 基础　　\textbf{方法：} 分布与奇性

在整个实轴上按分布与光滑函数的配对进行计算；\(\delta'\)
表示对自变量求导，证明恒等式时应对任意测试函数检验，而非仅比较普通函数值。

\begin{enumerate}
\def\labelenumi{\arabic{enumi}.}
\tightlist
\item
  计算
  \(\displaystyle I(a)=\int_{-\infty}^{\infty}(x^2+\sin x)\delta'(x-a)\,dx\)；
\item
  证明分布恒等式 \(\displaystyle x\delta'(x)=-\delta(x)\)。
\end{enumerate}


\end{minipage}\par

\par\noindent\begin{minipage}{\linewidth}
\textbf{解答}

\begin{center}
\includegraphics[width=\linewidth,height=.40\textheight,keepaspectratio]{\detokenize{assets/module-01/examples/M1-70.pdf}}
\par\smallskip{\small 图 M1-70：δ′ 的筛选与乘积恒等式}
\end{center}
\end{minipage}\par

图 M1-70中 δ′
的正、负两瓣只是正则化示意，表示它对测试函数的负斜率敏感。第一问把导数移到
\(x^2+\sin x\)，得到 \(-(2a+\cos a)\)；第二问则把测试函数改成
\(xf\)，取支撑点为原点，按右下公式卡读取
\((xf)^{\prime}(0)=f(0)\)，因此得到
\(-\delta\)。两问的负号来自同一次分布导数转移。

由 \(\delta'\) 的定义，

\[
\begin{aligned}
I(a)
&=
-\left.
\frac d{dx}(x^2+\sin x)
\right|_{x=a}\\
&=
\boxed{-(2a+\cos a)}.
\end{aligned}
\]

对任意测试函数 \(f\)，

\[
\begin{aligned}
\langle x\delta',f\rangle
&=
\langle\delta',xf\rangle\\
&=
-(xf)'|_{x=0}\\
&=-f(0)\\
&=
\langle-\delta,f\rangle.
\end{aligned}
\]

因此

\[
\boxed{x\delta'(x)=-\delta(x)}.
\]

\par\addvspace{.4\baselineskip}
\noindent\begin{minipage}{\linewidth}

\subsubsection{作业三：均匀带电细圆环的体密度表示}\label{m1:ux4f5cux4e1aux4e09ux5747ux5300ux5e26ux7535ux7ec6ux5706ux73afux7684ux4f53ux5bc6ux5ea6ux8868ux793a}

\textbf{题号：M1-71}

\textbf{层级：} 基础　　\textbf{方法：} 分布与奇性

圆环视为厚度趋于零的理想线源，电荷在方位角方向均匀分布；\(a>0\)，圆柱坐标体积元采用
\(dV=\rho\,d\rho\,d\phi\,dz\)。

半径为 \(a\)、位于 \(z=0\) 平面且圆心在原点的均匀细圆环带总电荷
\(Q\)。在圆柱坐标中写出等效体电荷密度，并验证总电荷。


\end{minipage}\par

\par\noindent\begin{minipage}{\linewidth}
\textbf{解答}

\begin{center}
\includegraphics[width=\linewidth,height=.40\textheight,keepaspectratio]{\detokenize{assets/module-01/examples/M1-71.pdf}}
\par\smallskip{\small 图 M1-71：均匀细圆环的体密度}
\end{center}
\end{minipage}\par

图 M1-71左图是均匀细圆环，子午截面把它约束为
\(\rho=a,z=0\)，而方位角仍可自由变化。因而密度需要两个
δ；与右图的柱坐标体积元相乘后，径向积分额外读出一个 \(a\)，方位积分读出
\(2\pi\)。归一化系数必须消去这两个因子，才使总电荷回到 \(Q\)。

圆环线电荷密度为

\[
\lambda=\frac{Q}{2\pi a}.
\]

圆环由 \(\rho=a\) 与 \(z=0\) 两个条件共同确定，因此

\[
\boxed{
\rho_{\mathrm v}(\rho,\phi,z)
=
\frac{Q}{2\pi a}
\delta(\rho-a)\delta(z)
}.
\]

量纲为

\[
\frac{Q}{L}\cdot\frac1L\cdot\frac1L
=
\frac{Q}{L^3},
\]

与体电荷密度一致。利用 \(dV=\rho\,d\rho\,d\phi\,dz\)，

\[
\begin{aligned}
\int\rho_{\mathrm v}\,dV
&=
\frac{Q}{2\pi a}
\int_0^{\infty}\delta(\rho-a)\rho\,d\rho
\int_0^{2\pi}d\phi
\int_{-\infty}^{\infty}\delta(z)\,dz\\
&=
\frac{Q}{2\pi a}\cdot a\cdot2\pi\\
&=
\boxed{Q}.
\end{aligned}
\]

这里的因子 \(a\) 来自图 M1-71 中的柱坐标体积元，不能遗漏。

\par\addvspace{.4\baselineskip}
\noindent\begin{minipage}{\linewidth}

\subsubsection{作业四：点电荷电势与泊松方程}\label{m1:ux4f5cux4e1aux56dbux70b9ux7535ux8377ux7535ux52bfux4e0eux6ccaux677eux65b9ux7a0b}

\textbf{题号：M1-72}

\textbf{层级：} 基础　　\textbf{方法：} 分布与奇性、边值问题

介质为真空，\(r=|\mathbf r|\)；原点处应按三维分布理解，不能仅凭原点以外的普通拉普拉斯计算判断整个空间的源项。

原点处点电荷的电势为

\[
\varphi(\mathbf r)
=
\frac{q}{4\pi\varepsilon_0r}.
\]

严格证明其满足

\[
\nabla^2\varphi
=
-\frac q{\varepsilon_0}
\delta^{(3)}(\mathbf r).
\]


\end{minipage}\par

\par\noindent\begin{minipage}{\linewidth}
\textbf{解答}

\begin{center}
\includegraphics[width=\linewidth,height=.40\textheight,keepaspectratio]{\detokenize{assets/module-01/examples/M1-72.pdf}}
\par\smallskip{\small 图 M1-72：点电荷电势的分布拉普拉斯}
\end{center}
\end{minipage}\par

图 M1-72中的白色小球排除了原点，穿孔区域内的 \(1/r\)
是调和函数。外球面靠近测试函数支集外，边界项为零；内球面的区域外法向指向球心，使
\(\mathbf r/r^3\) 的通量贡献为负。令孔径趋零后读到 \(-4\pi f(0)\)，再乘
\(q/(4\pi\varepsilon_0)\)，就得到题设电势的集中拉普拉斯项。

在 \(r>0\)，球坐标拉普拉斯算符给出

\[
\nabla^2\left(\frac1r\right)
=
\frac1{r^2}\frac d{dr}
\left[
r^2\frac d{dr}\left(\frac1r\right)
\right]
=0.
\]

但这只说明原点外的普通函数部分为零。取任意
\(f\in C_c^{\infty}(\mathbb R^3)\)，选取 \(R\) 使
\(\operatorname{supp}f\subset B_R\)，并令

\[
\Omega_{\varepsilon}
=
B_R\setminus\overline{B_{\varepsilon}}.
\]

由

\[
\nabla\left(\frac1r\right)
=
-\frac{\mathbf r}{r^3}
\]

以及分布导数定义，

\[
\begin{aligned}
\left\langle
\nabla^2\left(\frac1r\right),f
\right\rangle
&=
-\left\langle
\nabla\left(\frac1r\right),\nabla f
\right\rangle\\
&=
\lim_{\varepsilon\to0^+}
\int_{\Omega_{\varepsilon}}
\frac{\mathbf r}{r^3}\cdot\nabla f\,dV.
\end{aligned}
\]

在 \(\Omega_{\varepsilon}\) 内
\(\nabla\cdot(\mathbf r/r^3)=0\)，故高斯定理给出

\[
\int_{\Omega_{\varepsilon}}
\frac{\mathbf r}{r^3}\cdot\nabla f\,dV
=
\oint_{\partial\Omega_{\varepsilon}}
f\frac{\mathbf r}{r^3}\cdot d\mathbf S.
\]

外球面附近 \(f=0\)。图 M1-72 的内球面是穿孔区域的内边界，其外法向为
\(-\hat{\mathbf r}\)，所以

\[
\begin{aligned}
\oint_{r=\varepsilon}
f\frac{\mathbf r}{r^3}\cdot d\mathbf S
&=
-\int_{S^2}
f(\varepsilon\hat{\mathbf r})\,d\Omega\\
&\xrightarrow[\varepsilon\to0^+]{}
-4\pi f(\mathbf0).
\end{aligned}
\]

因此

\[
\left\langle
\nabla^2\left(\frac1r\right),f
\right\rangle
=
-4\pi f(\mathbf0).
\]

因此

\[
\nabla^2\left(\frac1r\right)
=
-4\pi\delta^{(3)}(\mathbf r).
\]

乘以 \(q/(4\pi\varepsilon_0)\)，得到

\[
\boxed{
\nabla^2\varphi
=
-\frac q{\varepsilon_0}
\delta^{(3)}(\mathbf r)
}.
\]

这与泊松方程 \(\nabla^2\varphi=-\rho/\varepsilon_0\) 及
\(\rho=q\delta^{(3)}(\mathbf r)\) 完全一致。

\section{模块一附录：常用正交曲线坐标系下的微元与算符表（基础必讲与查表选讲）}\label{m1:ux6a21ux5757ux4e00ux9644ux5f55ux5e38ux7528ux6b63ux4ea4ux66f2ux7ebfux5750ux6807ux7cfbux4e0bux7684ux5faeux5143ux4e0eux7b97ux7b26ux8868ux57faux7840ux5fc5ux8bb2ux4e0eux67e5ux8868ux9009ux8bb2}

电磁学中常见柱对称、球对称和平面极坐标对称。使用与边界和对称性相配的坐标，可以省去大量代数计算。模块一最后集中列出常用坐标系的微元和算符，供静电场、多极展开、分离变量和高斯定理计算查用。

坐标变换、尺度因子、线元、面积元和体积元属于基础内容；完整的梯度、散度、旋度与拉普拉斯算符表可在具体问题中按需调用。

本讲义约定：圆柱坐标为 \((\rho,\phi,z)\)；球坐标为
\((r,\theta,\phi)\)，其中 \(\theta\) 是从 \(+z\) 轴量起的极角，\(\phi\)
是绕 \(z\) 轴的方位角。

\subsubsection{图示：柱坐标与球坐标的微元来源}\label{m1:ux56feux793aux67f1ux5750ux6807ux4e0eux7403ux5750ux6807ux7684ux5faeux5143ux6765ux6e90}

\begin{center}
\includegraphics[width=\linewidth,height=.40\textheight,keepaspectratio]{\detokenize{assets/module-01/tikz-m1m2-08-coordinate-elements-v2-safe-v7.png}}
\par\smallskip{\small 图示：柱坐标与球坐标局部微元的三条正交边长及体积元}
\end{center}

柱坐标中弧长来自 \(\rho\,d\phi\)；球坐标中纬向弧长来自
\(r\sin\theta\,d\phi\)。记住这两个弧长，面积元和体积元就不再是死记硬背。

\subsection{一、一般正交曲线坐标的公式}\label{m1:ux4e00ux4e00ux822cux6b63ux4ea4ux66f2ux7ebfux5750ux6807ux7684ux516cux5f0f}

\begin{center}
\includegraphics[width=\linewidth,height=.40\textheight,keepaspectratio]{\detokenize{assets/module-01/figure-108-revised.pdf}}
\par\smallskip{\small 图示：正交曲线坐标的尺度因子把坐标增量转化为真实长度}
\end{center}

设正交曲线坐标为 \((q_1,q_2,q_3)\)，三个单位矢量为
\(\mathbf{e}_1,\mathbf{e}_2,\mathbf{e}_3\)，对应的尺度因子为
\(h_1,h_2,h_3\)。微小位移为

\[
d\mathbf{l}
=
h_1dq_1\,\mathbf{e}_1
+
h_2dq_2\,\mathbf{e}_2
+
h_3dq_3\,\mathbf{e}_3.
\]

三个坐标面上的有向面积元分别为

\[
d\mathbf{S}_{q_1}=h_2h_3\,dq_2dq_3\,\mathbf{e}_1,
\]

\[
d\mathbf{S}_{q_2}=h_3h_1\,dq_3dq_1\,\mathbf{e}_2,
\]

\[
d\mathbf{S}_{q_3}=h_1h_2\,dq_1dq_2\,\mathbf{e}_3.
\]

体积元为

\[
dV=h_1h_2h_3\,dq_1dq_2dq_3.
\]

若标量场为 \(u\)，矢量场为

\[
\mathbf{A}=A_1\mathbf{e}_1+A_2\mathbf{e}_2+A_3\mathbf{e}_3,
\]

则梯度为

\[
\nabla u
=
\frac{1}{h_1}\frac{\partial u}{\partial q_1}\mathbf{e}_1
+
\frac{1}{h_2}\frac{\partial u}{\partial q_2}\mathbf{e}_2
+
\frac{1}{h_3}\frac{\partial u}{\partial q_3}\mathbf{e}_3.
\]

散度为

\[
\nabla\cdot\mathbf{A}
=
\frac{1}{h_1h_2h_3}
\left[
\frac{\partial}{\partial q_1}(h_2h_3A_1)
+
\frac{\partial}{\partial q_2}(h_3h_1A_2)
+
\frac{\partial}{\partial q_3}(h_1h_2A_3)
\right].
\]

旋度可写成行列式形式：

\[
\nabla\times\mathbf{A}
=
\frac{1}{h_1h_2h_3}
\begin{vmatrix}
h_1\mathbf{e}_1 & h_2\mathbf{e}_2 & h_3\mathbf{e}_3\\
\frac{\partial}{\partial q_1} & \frac{\partial}{\partial q_2} & \frac{\partial}{\partial q_3}\\
h_1A_1 & h_2A_2 & h_3A_3
\end{vmatrix}.
\]

标量拉普拉斯算符为

\[
\nabla^2u
=
\frac{1}{h_1h_2h_3}
\left[
\frac{\partial}{\partial q_1}
\left(
\frac{h_2h_3}{h_1}\frac{\partial u}{\partial q_1}
\right)
+
\frac{\partial}{\partial q_2}
\left(
\frac{h_3h_1}{h_2}\frac{\partial u}{\partial q_2}
\right)
+
\frac{\partial}{\partial q_3}
\left(
\frac{h_1h_2}{h_3}\frac{\partial u}{\partial q_3}
\right)
\right].
\]

这组公式是后面所有坐标系算符的来源。尺度因子 \(h_i\) 表示坐标增量
\(dq_i\) 与真实长度的换算关系：真实弧长为 \(h_i dq_i\)，一般不等于
\(dq_i\)。

\subsection{\texorpdfstring{二、 直角坐标
\((x,y,z)\)}{二、 直角坐标 (x,y,z)}}\label{m1:ux4e8c-ux76f4ux89d2ux5750ux6807-xyz}

尺度因子为

\[
h_x=h_y=h_z=1.
\]

微元为

\[
d\mathbf{l}=dx\,\mathbf{i}+dy\,\mathbf{j}+dz\,\mathbf{k},
\qquad
dV=dx\,dy\,dz.
\]

三个坐标面的有向面积元为

\[
d\mathbf{S}_x=dy\,dz\,\mathbf{i},
\qquad
d\mathbf{S}_y=dz\,dx\,\mathbf{j},
\qquad
d\mathbf{S}_z=dx\,dy\,\mathbf{k}.
\]

常用算符为

\[
\nabla u
=
\frac{\partial u}{\partial x}\mathbf{i}
+
\frac{\partial u}{\partial y}\mathbf{j}
+
\frac{\partial u}{\partial z}\mathbf{k},
\]

\[
\nabla\cdot\mathbf{A}
=
\frac{\partial A_x}{\partial x}
+
\frac{\partial A_y}{\partial y}
+
\frac{\partial A_z}{\partial z},
\]

\[
\nabla^2u
=
\frac{\partial^2u}{\partial x^2}
+
\frac{\partial^2u}{\partial y^2}
+
\frac{\partial^2u}{\partial z^2}.
\]

直角坐标适合平板、矩形边界、均匀场叠加等问题；若题目具有明显球对称或柱对称，通常不应优先使用直角坐标。

\subsection{\texorpdfstring{三、 圆柱坐标
\((\rho,\phi,z)\)}{三、 圆柱坐标 (\textbackslash rho,\textbackslash phi,z)}}\label{m1:ux4e09-ux5706ux67f1ux5750ux6807-rhophiz}

圆柱坐标与直角坐标的关系为

\[
x=\rho\cos\phi,\qquad y=\rho\sin\phi,\qquad z=z.
\]

\begin{center}
\includegraphics[width=\linewidth,height=.40\textheight,keepaspectratio]{\detokenize{assets/module-01/tikz-m1-coordinate-02-cylindrical-element.png}}
\par\smallskip{\small 图示：圆柱坐标微元中切向真实长度是 rho dphi}
\end{center}

尺度因子为

\[
h_{\rho}=1,\qquad h_{\phi}=\rho,\qquad h_z=1.
\]

微小位移为

\[
d\mathbf{l}
=
d\rho\,\mathbf{e}_{\rho}
+
\rho\,d\phi\,\mathbf{e}_{\phi}
+
dz\,\mathbf{e}_z.
\]

体积元为

\[
dV=\rho\,d\rho\,d\phi\,dz.
\]

三个常用有向面积元为

\[
d\mathbf{S}_{\rho}
=
\rho\,d\phi\,dz\,\mathbf{e}_{\rho},
\]

\[
d\mathbf{S}_{\phi}
=
d\rho\,dz\,\mathbf{e}_{\phi},
\]

\[
d\mathbf{S}_z
=
\rho\,d\rho\,d\phi\,\mathbf{e}_z.
\]

对标量场 \(u(\rho,\phi,z)\)，

\[
\nabla u
=
\frac{\partial u}{\partial \rho}\mathbf{e}_{\rho}
+
\frac{1}{\rho}\frac{\partial u}{\partial\phi}\mathbf{e}_{\phi}
+
\frac{\partial u}{\partial z}\mathbf{e}_z.
\]

对矢量场

\[
\mathbf{A}=A_{\rho}\mathbf{e}_{\rho}+A_{\phi}\mathbf{e}_{\phi}+A_z\mathbf{e}_z,
\]

散度为

\[
\nabla\cdot\mathbf{A}
=
\frac{1}{\rho}\frac{\partial}{\partial\rho}(\rho A_{\rho})
+
\frac{1}{\rho}\frac{\partial A_{\phi}}{\partial\phi}
+
\frac{\partial A_z}{\partial z}.
\]

旋度为

\[
\nabla\times\mathbf{A}
=
\left(
\frac{1}{\rho}\frac{\partial A_z}{\partial\phi}
-
\frac{\partial A_{\phi}}{\partial z}
\right)\mathbf{e}_{\rho}
+
\left(
\frac{\partial A_{\rho}}{\partial z}
-
\frac{\partial A_z}{\partial\rho}
\right)\mathbf{e}_{\phi}
+
\frac{1}{\rho}
\left[
\frac{\partial}{\partial\rho}(\rho A_{\phi})
-
\frac{\partial A_{\rho}}{\partial\phi}
\right]\mathbf{e}_z.
\]

标量拉普拉斯算符为

\[
\nabla^2u
=
\frac{1}{\rho}\frac{\partial}{\partial\rho}
\left(
\rho\frac{\partial u}{\partial\rho}
\right)
+
\frac{1}{\rho^2}\frac{\partial^2u}{\partial\phi^2}
+
\frac{\partial^2u}{\partial z^2}.
\]

若问题具有轴对称性，则 \(\partial/\partial\phi=0\)；若还具有沿 \(z\)
方向平移不变性，则
\(\partial/\partial z=0\)。这时许多二维柱对称静电题会大幅简化。

\subsection{\texorpdfstring{四、 球坐标
\((r,\theta,\phi)\)}{四、 球坐标 (r,\textbackslash theta,\textbackslash phi)}}\label{m1:ux56db-ux7403ux5750ux6807-rthetaphi}

球坐标与直角坐标的关系为

\[
x=r\sin\theta\cos\phi,\qquad
y=r\sin\theta\sin\phi,\qquad
z=r\cos\theta.
\]

\begin{center}
\includegraphics[width=\linewidth,height=.40\textheight,keepaspectratio]{\detokenize{assets/module-01/tikz-m1-coordinate-03-spherical-element.png}}
\par\smallskip{\small 图示：球坐标微元中两个角向长度分别是 r dtheta 与 r sin theta dphi}
\end{center}

尺度因子为

\[
h_r=1,\qquad h_{\theta}=r,\qquad h_{\phi}=r\sin\theta.
\]

微小位移为

\[
d\mathbf{l}
=
dr\,\mathbf{e}_r
+
r\,d\theta\,\mathbf{e}_{\theta}
+
r\sin\theta\,d\phi\,\mathbf{e}_{\phi}.
\]

体积元为

\[
dV=r^2\sin\theta\,dr\,d\theta\,d\phi.
\]

三个常用有向面积元为

\[
d\mathbf{S}_r
=
r^2\sin\theta\,d\theta\,d\phi\,\mathbf{e}_r,
\]

\[
d\mathbf{S}_{\theta}
=
r\sin\theta\,dr\,d\phi\,\mathbf{e}_{\theta},
\]

\[
d\mathbf{S}_{\phi}
=
r\,dr\,d\theta\,\mathbf{e}_{\phi}.
\]

对标量场 \(u(r,\theta,\phi)\)，

\[
\nabla u
=
\frac{\partial u}{\partial r}\mathbf{e}_r
+
\frac{1}{r}\frac{\partial u}{\partial\theta}\mathbf{e}_{\theta}
+
\frac{1}{r\sin\theta}\frac{\partial u}{\partial\phi}\mathbf{e}_{\phi}.
\]

对矢量场

\[
\mathbf{A}=A_r\mathbf{e}_r+A_{\theta}\mathbf{e}_{\theta}+A_{\phi}\mathbf{e}_{\phi},
\]

散度为

\[
\nabla\cdot\mathbf{A}
=
\frac{1}{r^2}\frac{\partial}{\partial r}(r^2A_r)
+
\frac{1}{r\sin\theta}\frac{\partial}{\partial\theta}(\sin\theta A_{\theta})
+
\frac{1}{r\sin\theta}\frac{\partial A_{\phi}}{\partial\phi}.
\]

旋度为

\[
\nabla\times\mathbf{A}
=
\frac{1}{r\sin\theta}
\left[
\frac{\partial}{\partial\theta}(\sin\theta A_{\phi})
-
\frac{\partial A_{\theta}}{\partial\phi}
\right]\mathbf{e}_r
\]

\[
\quad+
\frac{1}{r}
\left[
\frac{1}{\sin\theta}\frac{\partial A_r}{\partial\phi}
-
\frac{\partial}{\partial r}(rA_{\phi})
\right]\mathbf{e}_{\theta}
+
\frac{1}{r}
\left[
\frac{\partial}{\partial r}(rA_{\theta})
-
\frac{\partial A_r}{\partial\theta}
\right]\mathbf{e}_{\phi}.
\]

标量拉普拉斯算符为

\[
\nabla^2u
=
\frac{1}{r^2}\frac{\partial}{\partial r}
\left(
r^2\frac{\partial u}{\partial r}
\right)
+
\frac{1}{r^2\sin\theta}
\frac{\partial}{\partial\theta}
\left(
\sin\theta\frac{\partial u}{\partial\theta}
\right)
+
\frac{1}{r^2\sin^2\theta}
\frac{\partial^2u}{\partial\phi^2}.
\]

若问题具有球对称性，则 \(u=u(r)\)，拉普拉斯算符简化为

\[
\nabla^2u
=
\frac{1}{r^2}\frac{d}{dr}
\left(
r^2\frac{du}{dr}
\right).
\]

这正是点电荷电势、球壳电势、球形导体边值问题中最常用的形式。

\subsection{\texorpdfstring{五、 平面极坐标
\((\rho,\phi)\)}{五、 平面极坐标 (\textbackslash rho,\textbackslash phi)}}\label{m1:ux4e94-ux5e73ux9762ux6781ux5750ux6807-rhophi}

平面极坐标可以看成圆柱坐标去掉 \(z\) 方向。微小位移和面积元为

\[
d\mathbf{l}
=
d\rho\,\mathbf{e}_{\rho}
+
\rho\,d\phi\,\mathbf{e}_{\phi},
\qquad
dS=\rho\,d\rho\,d\phi.
\]

对标量场 \(u(\rho,\phi)\)，

\[
\nabla u
=
\frac{\partial u}{\partial\rho}\mathbf{e}_{\rho}
+
\frac{1}{\rho}\frac{\partial u}{\partial\phi}\mathbf{e}_{\phi}.
\]

对二维矢量场

\[
\mathbf{A}=A_{\rho}\mathbf{e}_{\rho}+A_{\phi}\mathbf{e}_{\phi},
\]

散度为

\[
\nabla\cdot\mathbf{A}
=
\frac{1}{\rho}\frac{\partial}{\partial\rho}(\rho A_{\rho})
+
\frac{1}{\rho}\frac{\partial A_{\phi}}{\partial\phi}.
\]

标量拉普拉斯算符为

\[
\nabla^2u
=
\frac{1}{\rho}\frac{\partial}{\partial\rho}
\left(
\rho\frac{\partial u}{\partial\rho}
\right)
+
\frac{1}{\rho^2}\frac{\partial^2u}{\partial\phi^2}.
\]

若场只与 \(\rho\) 有关，则

\[
\nabla^2u
=
\frac{1}{\rho}\frac{d}{d\rho}
\left(
\rho\frac{du}{d\rho}
\right).
\]

这正是无限长带电直线、同轴圆柱电容器、二维柱坐标分离变量法中的基础公式。

\subsection{六、坐标系的选择判据}\label{m1:ux516dux5750ux6807ux7cfbux7684ux9009ux62e9ux5224ux636e}

\begin{center}
\includegraphics[width=\linewidth,height=.40\textheight,keepaspectratio]{\detokenize{assets/module-01/figure-111-revised.pdf}}
\par\smallskip{\small 图示：选坐标系应先看源和边界，再看对称性与坐标面}
\end{center}

\begin{enumerate}
\def\labelenumi{\arabic{enumi}.}
\tightlist
\item
  \textbf{球对称}：点电荷、带电球壳、导体球、电偶极远场，优先用球坐标。
\item
  \textbf{柱对称}：无限长直线电荷、长直导线、同轴圆柱、螺线管近似，优先用圆柱坐标。
\item
  \textbf{平面对称或矩形边界}：平行板、电容器边缘忽略、矩形盒分离变量，优先用直角坐标。
\item
  \textbf{二维圆形边界}：圆柱导体、电像法中的圆柱问题、保角变换前后的圆边界，常用平面极坐标。
\item
  \textbf{先看边界，再选坐标}。边界面若恰好是某坐标系的坐标面，变量分离和积分计算通常会明显简化。
\end{enumerate}

本附录可随具体对称性反复查阅，以建立``几何对称性''与``算符形式''的对应关系。

\section{模块一小结}\label{m1:ux6a21ux5757ux4e00ux5c0fux7ed3}

{\def\LTcaptype{none} % do not increment counter
\begin{longtable}[]{@{}
  >{\raggedright\arraybackslash}p{(\linewidth - 4\tabcolsep) * \real{0.3333}}
  >{\raggedright\arraybackslash}p{(\linewidth - 4\tabcolsep) * \real{0.3333}}
  >{\raggedright\arraybackslash}p{(\linewidth - 4\tabcolsep) * \real{0.3333}}@{}}
\toprule\noalign{}
\begin{minipage}[b]{\linewidth}\raggedright
检验问题
\end{minipage} & \begin{minipage}[b]{\linewidth}\raggedright
结论或判断
\end{minipage} & \begin{minipage}[b]{\linewidth}\raggedright
回看位置
\end{minipage} \\
\midrule\noalign{}
\endhead
\bottomrule\noalign{}
\endlastfoot
能否先判断
\(\nabla f\)、\(\nabla\cdot\mathbf A\)、\(\nabla\times\mathbf A\)
的输出类型？ & 由输入对象和连接方式决定 & 专题二 \\
标量三重积为负是否表示体积为负？ & 负号表示取向，体积取绝对值 &
专题一 \\
梯度为何垂直等值面？ & \(df=\nabla f\cdot d\mathbf r\) & 专题三 \\
散度为零是否保证所有闭曲面通量为零？ & 仅在区域光滑且没有遗漏奇点时 &
专题四、七 \\
旋度方向如何与回路正向配对？ & 按右手定则同时选择 & 专题五、六 \\
积分式何时可以局域化？ & 场在所考察邻域内足够光滑 & 专题六、七 \\
\end{longtable}
}

\section{模块一综合题：奇异径向场和旋转场}\label{m1:ux6a21ux5757ux4e00ux7efcux5408ux9898ux5947ux5f02ux5f84ux5411ux573aux548cux65cbux8f6cux573a}

\par\addvspace{.4\baselineskip}
\noindent\begin{minipage}{\linewidth}

\subsection{第一题：含点奇点的径向场}\label{m1:ux7b2cux4e00ux9898ux542bux70b9ux5947ux70b9ux7684ux5f84ux5411ux573a}

\textbf{题号：M1-73}

\textbf{层级：} 复赛　　\textbf{方法：} 分布与奇性

取三维球坐标，\(\mathbf e_r\)
沿外径向，\(R>0\)；球面通量一律使用球体外法向。求全空间结果时须将原点可能存在的集中源与普通体源分开。

在 \(r>0\) 的区域给定标量场

\[
\varphi(r)=\frac{\alpha}{2}r^2+\frac{\beta}{r},
\qquad
\mathbf E=-\nabla\varphi,
\]

其中 \(\alpha\)、\(\beta\) 为常量。

\begin{enumerate}
\def\labelenumi{\arabic{enumi}.}
\tightlist
\item
  用球坐标求 \(\mathbf E\)、\(\nabla\cdot\mathbf E\) 和
  \(\nabla\times\mathbf E\)。
\item
  计算半径为 \(R\) 的球面通量，并写出 \(\nabla\cdot\mathbf E\)
  在全空间的分布形式。
\item
  若把 \(\mathbf E\)
  解释为静电场，求对应的电荷密度，并说明普通微分式与积分式为什么并不矛盾。
\end{enumerate}


\subsubsection{解答}\label{m1:ux89e3ux7b54}

\begin{center}
\includegraphics[width=\linewidth,height=.40\textheight,keepaspectratio]{\detokenize{assets/module-01/examples/M1-73.pdf}}
\par\smallskip{\small 图 M1-73：点奇点与均匀体源的叠加}
\end{center}
\end{minipage}\par

图
M1-73中把场拆成橙色的线性径向项与蓝色的平方反比项，两类源必须分别记账。前者在球内贡献
\(-4\pi\alpha R^3\)，后者在任意包围原点的球面贡献 \(4\pi\beta\)。先在
\(r>0\) 计算普通散度 \(-3\alpha\)，再用图示与半径无关的剩余通量补上
\(4\pi\beta\delta^{(3)}\)，二者之和才与总球面通量一致。

\[
\boxed{
\mathbf E
=\left(-\alpha r+\frac{\beta}{r^2}\right)\mathbf e_r
}.
\]

在 \(r>0\)，

\[
\nabla\cdot\mathbf E
=\frac1{r^2}\frac{d}{dr}\left(r^2E_r\right)
=\frac1{r^2}\frac{d}{dr}\left(-\alpha r^3+\beta\right)
=-3\alpha,
\]

而径向场由单值势函数产生，所以

\[
\nabla\times\mathbf E=\mathbf0,
\qquad r>0.
\]

球面上的法向就是 \(\mathbf e_r\)，故

\[
\Phi_E(R)
=4\pi R^2E_r(R)
=4\pi\left(-\alpha R^3+\beta\right).
\]

图 M1-73 中线性径向项的普通散度 \(-3\alpha\) 在球内的体积分只给出
\(-4\pi\alpha R^3\)，还缺少与半径无关的
\(4\pi\beta\)。缺少的通量集中在原点，因此全空间分布形式为

\[
\boxed{
\nabla\cdot\mathbf E
=-3\alpha+4\pi\beta\,\delta^{(3)}(\mathbf r)
}.
\]

由 \(\nabla\cdot\mathbf E=\rho/\varepsilon_0\)，对应源为

\[
\boxed{
\rho(\mathbf r)
=-3\varepsilon_0\alpha
+4\pi\varepsilon_0\beta\,\delta^{(3)}(\mathbf r)
}.
\]

第一项是均匀体电荷，第二项是位于原点、电量为 \(4\pi\varepsilon_0\beta\)
的点电荷。\(r>0\)
的普通微分式没有看见被挖去的奇点；把原点补回并采用分布散度后，高斯积分式与微分式完全一致。

\par\addvspace{.4\baselineskip}
\noindent\begin{minipage}{\linewidth}

\subsection{第二题：旋转场的环量与散度}\label{m1:ux7b2cux4e8cux9898ux65cbux8f6cux573aux7684ux73afux91cfux4e0eux6563ux5ea6}

\textbf{题号：M1-74}

\textbf{层级：} 复赛　　\textbf{方法：} 高斯定理、旋度与环量

采用右手单位正交基，\(R>0\)；圆盘法向与圆周方向按右手定则配对，所有导数都在给定光滑场的普通定义域内进行。

设 \(a\) 为常量，

\[
\mathbf A=\frac a2(-y\,\mathbf i+x\,\mathbf j).
\]

取 \(xy\) 平面内以原点为圆心、半径为 \(R\) 的圆盘 \(S\)，其边界 \(C\)
按从 \(+z\) 方向观察的逆时针方向取正向。

\begin{enumerate}
\def\labelenumi{\arabic{enumi}.}
\tightlist
\item
  求 \(\nabla\cdot\mathbf A\) 与 \(\nabla\times\mathbf A\)。
\item
  分别用直接积分和斯托克斯定理计算
  \(\oint_C\mathbf A\cdot d\mathbf l\)。
\item
  说明为什么``散度为零''不意味着环量为零。
\end{enumerate}


\end{minipage}\par
\subsubsection{解答}\label{m1:ux89e3ux7b54-1}


图 M1-74中的积分边界固定为逆时针，纸外是正法向；场箭头仅用 \(a>0\)
的情形示意。圆周上场与切向同向，系数为
\(aR/2\)，所以直接积分是它乘以周长。斯托克斯路线把法向旋度 \(a\)
乘以圆盘面积，得到同一个 \(a\pi R^2\)；即使
\(a<0\)，也只让场及结果反号，不改变已选定的边界正向。
\newpage
\par\noindent\begin{minipage}{\linewidth}
\begin{center}
\includegraphics[width=\linewidth,height=.40\textheight,keepaspectratio]{\detokenize{assets/module-01/examples/M1-74.pdf}}
\par\smallskip{\small 图 M1-74：旋转场的环量、散度与旋度}
\end{center}
\end{minipage}\par

\[
\nabla\cdot\mathbf A=0,
\qquad
\nabla\times\mathbf A=a\mathbf k.
\]

沿图 M1-74 的正向圆周，令 \(x=R\cos\phi\)、\(y=R\sin\phi\)，则

\[
d\mathbf l=(-R\sin\phi\,\mathbf i+R\cos\phi\,\mathbf j)\,d\phi,
\]

\[
\mathbf A=\frac{aR}{2}(-\sin\phi\,\mathbf i+\cos\phi\,\mathbf j),
\]

所以

\[
\oint_C\mathbf A\cdot d\mathbf l
=\int_0^{2\pi}\frac{aR^2}{2}\,d\phi
=a\pi R^2.
\]

另一方面，\(S\) 的正法向为 \(\mathbf k\)，斯托克斯定理给出

\[
\int_S(\nabla\times\mathbf A)\cdot d\mathbf S
=\int_Sa\,dS=a\pi R^2,
\]

与直接积分相同。散度控制闭曲面的净通量，旋度控制开曲面边界的环量密度；该场无源但有旋，说明两种局域信息彼此独立。
\clearpage
\endgroup
